% Generated mechanically from frozen input inputs/p02/ja/tex/homology.tex.
% Do not edit this generated file; edit the bound locale source instead.

% OK, start here.
%

\setcounter{chapter}{11}
\chapter{ホモロジー代数}



\phantomsection
\label{homology-section-phantom}


\section{序論}
\label{homology-section-introduction}

\noindent
本章ではホモロジー代数の基礎を説明する。この分野には明らかに
無尽蔵の題材があるため、他の部で必要になるたびに内容を補っていく。
参考文献として \cite{Maclane} がある。

\section{基礎的事項}
\label{homology-section-topology-basic}

\noindent
以下の概念は基礎的なものとみなし、定義も証明も与えない。
これは、それらがすべて必ずしも容易または周知であるという意味ではない。

\begin{enumerate}
\item 今のところ何もない。
\end{enumerate}

\section{前加法圏と加法圏}
\label{homology-section-additive-categories}

\noindent
まず前加法圏を定義する。

\begin{definition}
\label{homology-definition-preadditive}
圏 $\mathcal{A}$ が {\it 前加法圏} であるとは、各射集合
$\Mor_\mathcal{A}(x, y)$ にアーベル群の構造が与えられ、合成写像
$$
\Mor(x, y) \times \Mor(y, z)
\longrightarrow
\Mor(x, z)
$$
が双線形であることをいう。前加法圏の間の関手
$F : \mathcal{A} \to \mathcal{B}$ が {\it 加法的} であるとは、任意の
$x, y \in \Ob(\mathcal{A})$ に対して
$F : \Mor(x, y) \to \Mor(F(x), F(y))$ がアーベル群の準同型であることをいう。
\end{definition}

\noindent
特に、任意の $x, y$ に対して少なくとも一つの射 $x \to y$、すなわち
零射が存在する。

\begin{lemma}
\label{homology-lemma-preadditive-zero}
$\mathcal{A}$ を前加法圏とし、$x$ を $\mathcal{A}$ の対象とする。
次は同値である。
\begin{enumerate}
\item $x$ は始対象である。
\item $x$ は終対象である。
\item $\Mor_\mathcal{A}(x, x)$ において $\text{id}_x = 0$ である。
\end{enumerate}
さらに、このような対象 $0$ が存在するとき、射 $\alpha : y \to z$ が
$0$ を経由して分解するための必要十分条件は $\alpha = 0$ である。
\end{lemma}

\begin{proof}
まず、$x$ が (1) 始対象または (2) 終対象であると仮定する。いずれの場合も
$\Mor(x,x)$ は $\text{id}_x$ を含む自明なアーベル群となるので、
$\Mor(x, x)$ において $\text{id}_x = 0$ である。したがって (1) と (2) は
それぞれ (3) を含意する。

\medskip\noindent
次に $\Mor(x,x)$ において $\text{id}_x = 0$ と仮定する。
$y$ を $\mathcal{A}$ の任意の対象とし、$f \in \Mor(x ,y)$ とする。
合成写像を
$C : \Mor(x,x) \times \Mor(x,y) \to \Mor(x,y)$ と書く。このとき
$f = C(0, f)$ であり、$C$ は双線形なので $C(0, f) = 0$ である。
よって $f = 0$ であり、$x$ は $\mathcal{A}$ の始対象である。
$f \in \Mor(y, x)$ に対する同様の議論により、$x$ は終対象でもある。
したがって (3) は (1) と (2) の両方を含意する。
\end{proof}

\begin{definition}
\label{homology-definition-zero-object}
前加法圏 $\mathcal{A}$ において、補題 \ref{homology-lemma-preadditive-zero} のように
始対象かつ終対象である対象を {\it 零対象} といい、$0$ と書く。
\end{definition}

\begin{lemma}
\label{homology-lemma-preadditive-direct-sum}
$\mathcal{A}$ を前加法圏とし、$x, y \in \Ob(\mathcal{A})$ とする。
積 $x \times y$ が存在すれば余積 $x \amalg y$ も存在し、
余積 $x \amalg y$ が存在すれば積 $x \times y$ も存在する。
この場合さらに $x \amalg y \cong x \times y$ である。
\end{lemma}

\begin{proof}
$z = x \times y$ とし、その射影を $p : z \to x$ および
$q : z \to y$ とする。$(1, 0)$ に対応する射を $i : x \to z$、
$(0, 1)$ に対応する射を $j : y \to z$ と書く。このとき、対角方向の
合成が零である可換図式
$$
\xymatrix{
x \ar[rr]^1 \ar[rd]^i & & x \\
& z \ar[ru]^p \ar[rd]^q & \\
y \ar[rr]^1 \ar[ru]^j & & y
}
$$
を得る。$i \circ p + j \circ q : z \to z$ は、$p$ と合成すれば $p$、
$q$ と合成すれば $q$ になる射なので、恒等射である。
射 $a : x \to w$ と $b : y \to w$ が与えられたとする。このとき
$a \circ p + b \circ q : z \to w$ を作ることができる。これにより
$\Mor(z, w) = \Mor(x, w) \times \Mor(y, w)$ という全単射が得られ、
$z = x \amalg y$ であることが分かる。

\medskip\noindent
積の代わりに余積 $x \amalg y$ が与えられた場合に射 $p, q$ を構成することは
読者に委ねる。
\end{proof}

\begin{definition}
\label{homology-definition-direct-sum}
前加法圏 $\mathcal{A}$ の対象の組 $x, y$ に対し、$x$ と $y$ の
{\it 直和} $x \oplus y$ とは、上の補題 \ref{homology-lemma-preadditive-direct-sum}
における射 $i, j, p, q$ を備えた直積 $x \times y$ のことである。
\end{definition}

\begin{remark}
\label{homology-remark-direct-sum}
補題 \ref{homology-lemma-preadditive-direct-sum} の証明から、$p$ と $q$ が与えられると
射 $i$, $j$ は
$p \circ i = \text{id}_x$、$q \circ j = \text{id}_y$、
$p \circ j = 0$、$q \circ i = 0$ という条件によって一意に定まる。
さらに自動的に
$i \circ p + j \circ q = \text{id}_{x \oplus y}$ が成り立つ。
同様に、$i$, $j$ が与えられると射 $p$, $q$ は一意に定まる。
最後に、対象 $x, y, z$ と射
$i : x \to z$, $j : y \to z$, $p : z \to x$, $q : z \to y$ が
$p \circ i = \text{id}_x$、$q \circ j = \text{id}_y$、
$p \circ j = 0$、$q \circ i = 0$、かつ
$i \circ p + j \circ q = \text{id}_z$ を満たすなら、$z$ は $x$ と $y$ の
直和であり、その四つの構造射は $i, j, p, q$ に等しい。
\end{remark}

\begin{lemma}
\label{homology-lemma-additive-additive}
$\mathcal{A}$, $\mathcal{B}$ を前加法圏とし、
$F : \mathcal{A} \to \mathcal{B}$ を加法関手とする。このとき $F$ は
直和を直和に、零対象を零対象に写す。
\end{lemma}

\begin{proof}
$F$ は加法的であるとする。$x$ と $y$ の直和 $z$ は、
注意 \ref{homology-remark-direct-sum} によれば、射
$i : x \to z$, $j : y \to z$, $p : z \to x$, $q : z \to y$ と
関係式
$p \circ i = \text{id}_x$、$q \circ j = \text{id}_y$、
$p \circ j = 0$、$q \circ i = 0$、
$i \circ p + j \circ q = \text{id}_z$ によって特徴づけられる。
加法性により、$F(x), F(y), F(z)$ と射
$F(i), F(j), F(p), F(q)$ はまったく同じ関係式を満たすので、$F(z)$ は
$F(x)$ と $F(y)$ の直和である。したがって $F$ は直和を直和に写す。

\medskip\noindent
$F$ が零対象を零対象に写すことは、補題 \ref{homology-lemma-preadditive-zero} の
特徴づけ (3) から従う。
\end{proof}

\begin{definition}
\label{homology-definition-additive-category}
圏 $\mathcal{A}$ が {\it 加法圏} であるとは、それが前加法圏であり、
有限積が存在すること、言い換えれば零対象と直和をもつことをいう。
\end{definition}

\noindent
実際、空積は有限積であり、存在すれば終対象である。

\begin{definition}
\label{homology-definition-kernel}
$\mathcal{A}$ を前加法圏とし、$f : x \to y$ を射とする。
\begin{enumerate}
\item $f$ の {\it 核} とは、次を満たす射 $i : z \to x$ のことである：
(a) $f \circ i = 0$ であり、(b) $f \circ i' = 0$ を満たす任意の
$i' : z' \to x$ に対して、$i' = i \circ g$ を満たす一意な射
$g : z' \to z$ が存在する。
\item $f$ の核が存在するとき、それを $\Ker(f) \to x$ と書く。
\item $f$ の {\it 余核} とは、次を満たす射 $p : y \to z$ のことである：
(a) $p \circ f = 0$ であり、(b) $p' \circ f = 0$ を満たす任意の
$p' : y \to z'$ に対して、$p' = g \circ p$ を満たす一意な射
$g : z \to z'$ が存在する。
\item $f$ の余核が存在するとき、それを $y \to \Coker(f)$ と書く。
\item $f$ の核が存在するとき、射 $\Ker(f) \to x$ の余核を
$f$ の {\it 余像} という。
\item 核と余像が存在するとき、それを $x \to \Coim(f)$ と書く。
\item $f$ の余核が存在するとき、射 $y \to \Coker(f)$ の核を
$f$ の {\it 像} という。
\item $f$ の余核と像が存在するとき、それを $\Im(f) \to y$ と書く。
\end{enumerate}
\end{definition}

\noindent
上の定義では「核」および「余核」と呼び、それらが一意な同型を除いて
一意であることを暗黙に用いた。これは米田の補題
（Categories, Section \ref{categories-section-opposite}）から従う。
実際、$f : x \to y$ の核は、対象 $z$ を集合
$\Ker(\Mor_\mathcal{A}(z, x) \to \Mor_\mathcal{A}(z, y))$ に写す関手を
表現する。余核の場合は双対である。

\medskip\noindent
まず直和と核の関係を次のように述べる。

\begin{lemma}
\label{homology-lemma-additive-cat-biproduct-kernel}
$\mathcal{C}$ を前加法圏とする。
補題 \ref{homology-lemma-preadditive-direct-sum} の射 $i, j, p, q$ を備えた
$x \oplus y$ を $\mathcal{C}$ における直和とする。このとき
$i : x \to x \oplus y$ は $q : x \oplus y \rightarrow y$ の核である。
双対的に、$p$ は $j$ の余核である。
\end{lemma}

\begin{proof}
$q \circ f = 0$ を満たす射 $f : z' \to x \oplus y$ を取る。
$f = i \circ g$ を満たす一意な射 $g : z' \to x$ が存在することを
示せばよい。$i \circ p + j \circ q$ は $x \oplus y$ 上の恒等射なので
$$
f = (i \circ p + j \circ q) \circ f = i \circ p \circ f
$$
であり、したがって $g = p \circ f$ とすればよい。一意性は
$p \circ i$ が $x$ 上の恒等射であることから従う。第二の主張の証明は
双対である。
\end{proof}

\begin{lemma}
\label{homology-lemma-kernel-mono}
$\mathcal{C}$ を前加法圏とし、$f : x \to y$ を $\mathcal{C}$ の射とする。
\begin{enumerate}
\item $f$ の核が存在すれば、その核はモノ射である。
\item $f$ の余核が存在すれば、その余核はエピ射である。
\item $f$ の核と余像が存在すれば、その余像はエピ射である。
\item $f$ の余核と像が存在すれば、その像はモノ射である。
\end{enumerate}
\end{lemma}

\begin{proof}
(1) は核の定義で要求される一意性から直ちに従う。(2) の証明は双対である。
(3) は余像が余核であるため (2) から従う。同様に (4) は (1) から従う。
\end{proof}

\begin{lemma}
\label{homology-lemma-coim-im-map}
$f : x \to y$ を、核、余核、像、余像がすべて存在する前加法圏の射とする。
このとき $f$ は
$x \to \Coim(f) \to \Im(f) \to y$ と一意に分解できる。
\end{lemma}

\begin{proof}
$\Ker(f) \to x \to y$ が零なので、標準的な射 $\Coim(f) \to y$ がある。
合成 $\Coim(f) \to y \to \Coker(f)$ は零である。実際、これは零である射
$x \to y \to \Coker(f)$ を生じさせる一意な射である
（一意性は補題 \ref{homology-lemma-kernel-mono} (3) から従う）。
したがって $\Coim(f) \to y$ は $\Im(f) \to y$ を通じて一意に分解し、
求める射が得られる。
\end{proof}

\begin{example}
\label{homology-example-not-abelian}
$k$ を体とする。$k$ 上のフィルター付きベクトル空間の圏を考える
（定義 \ref{homology-definition-filtered} を参照）。$V = W = k$ とし、
$$
F^iV
=
\left\{
\begin{matrix}
V & \text{if} & i < 0 \\
0 & \text{if} & i \geq 0
\end{matrix}
\right.
\text{ and }
F^iW
=
\left\{
\begin{matrix}
W & \text{if} & i \leq 0 \\
0 & \text{if} & i > 0
\end{matrix}
\right.
$$
であるフィルター付きベクトル空間 $(V, F)$ と $(W, F)$ を考える。
台となるベクトル空間上で $\text{id}_k$ に対応する写像 $f : V \to W$ は、
核と余核が自明であるにもかかわらず同型ではない。また
$\Coim(f) = V$ および $\Im(f) = W$ である。したがって、$k$ 上の
フィルター付きベクトル空間の圏はアーベル圏ではない。
\end{example}

\noindent
この節の最後に、冪等自己準同型の像の分裂に関する補題をいくつか述べる。
自己準同型 $f$ が冪等であるとは $f \circ f = f$ であることをいう。

\begin{remark}
\label{homology-remark-idempotent-symmetry}
$\mathcal{A}$ を前加法圏、$x$ を $\mathcal{A}$ の対象、
$f : x \to x$ を冪等射とする。$g = \text{id} _ x - f$ と置くと、
$f \circ g = 0$、$g \circ f = 0$、$g \circ g = g$ である。
\end{remark}

\begin{lemma}
\label{homology-lemma-idempotent-kernel-cokernel}
$\mathcal{A}$ を前加法圏、$x$ を $\mathcal{A}$ の対象、
$f : x \to x$ を冪等射とする。$g = \text{id} _ x - f$ と書き、
注意 \ref{homology-remark-idempotent-symmetry} により
$f \circ g = g \circ f = 0$ であることを思い出す。
\begin{enumerate}
\item $f$ が核 $i : \Ker (f) \to x$ をもち、$g = i \circ p$ を満たす
一意な射を $p$ とすると、$p \circ i = \text{id} _ {\Ker (f)}$ であり、
$p$ は $f$ の余核である。
\item $f$ が余核 $p : x \to \Coker (f)$ をもち、$g = i \circ p$ を満たす
一意な射を $i$ とすると、$p \circ i = \text{id} _ {\Coker (f)}$ であり、
$i$ は $f$ の核である。
\item $g$ が核 $j : \Ker (g) \to x$ をもち、$f = j \circ q$ を満たす
一意な射を $q$ とすると、$q \circ j = \text{id} _ {\Ker (g)}$ であり、
$q$ は $g$ の余核である。
\item $g$ が余核 $q : x \to \Coker (g)$ をもち、$f = j \circ q$ を満たす
一意な射を $j$ とすると、$q \circ j = \text{id} _ {\Coker (g)}$ であり、
$j$ は $g$ の核である。
\end{enumerate}
\end{lemma}

\begin{proof}
(1) について、まず $i = (f + g) \circ i = g \circ i$ を計算する。
$p$ の構成により、これは
$i \circ p \circ i = g \circ i = i \circ \text{id} _ {\Ker (f)}$ を
含意する。補題 \ref{homology-lemma-kernel-mono} により
$p \circ i = \text{id} _ {\Ker (f)}$ を得る。

\medskip\noindent
残るのは、$p : x \to \Ker (f)$ が $f$ の余核の普遍性を満たすこと、
すなわち $a \circ f = 0$ を満たす射 $a : x \to y$ が与えられたとき、
$a = b \circ p$ を満たす一意な射 $b : \Ker (f) \to y$ が存在することを
示すことである。存在は、まず $a = a \circ (f + g) = a \circ g$ を計算し、
$p$ の構成から $a = a \circ g = a \circ i \circ p$ と結論することで従う
（$b = a \circ i$ と取ればよい）。一意性は
$b = b \circ p \circ i = a \circ i$ から従う
（したがって $a$ が $b$ を決定する）。

\medskip\noindent
双対性、対称性（すなわち $g$ 自身も冪等であること；注意
\ref{homology-remark-idempotent-symmetry} を参照）、および双対性と対称性を
それぞれ用いれば、(2)、(3)、(4) は (1) と同様に従う。
\end{proof}

\begin{lemma}
\label{homology-lemma-idempotent-splitting}
$\mathcal{A}$ を前加法圏、$x$ を $\mathcal{A}$ の対象、
$f : x \to x$ を冪等射とし、$g = \text{id} _ x - f$ と書く。
$f$ と $g$ がそれぞれ核または余核の少なくとも一方をもち、さらに
$i , j , p , q$ が補題 \ref{homology-lemma-idempotent-kernel-cokernel} の記述どおりに
構成されるとする。このとき四つの（余）核はすべて存在し、
\begin{align*}
x
& \simeq \Ker (f) \oplus \Ker (g) \\
& \simeq \Coker (f) \oplus \Ker (g) \\
& \simeq \Ker (f) \oplus \Coker (g) \\
& \simeq \Coker (f) \oplus \Coker (g)
\end{align*}
である。各直積の構造は $i , j , p , q$ によって与えられる。
\end{lemma}

\begin{proof}
補題 \ref{homology-lemma-idempotent-kernel-cokernel} により
$p \circ i = \text{id} _ {\Ker (f)} = \text{id} _ {\Coker (f)}$ かつ
$q \circ j = \text{id} _ {\Ker (g)} = \text{id} _ {\Coker (g)}$ である。
構成から $i \circ p + j \circ q = f + g = \text{id} _ x$ である。
注意 \ref{homology-remark-direct-sum} により主張が従う。
\end{proof}

\begin{lemma}
\label{homology-lemma-split-morphism-kernel-cokernel}
$\mathcal{A}$ を前加法圏、$x$ と $y$ を $\mathcal{A}$ の対象とし、
$q \circ j = \text{id} _ y$ を満たす射 $j : y \to x$ と $q : x \to y$ を
取る。このとき
\begin{enumerate}
\item $\text{id} _ x - j \circ q$ は核 $j$ をもつ。
\item $\text{id} _ x - j \circ q$ は余核 $q$ をもつ。
\end{enumerate}
\end{lemma}

\begin{proof}
(1) について、まず $(\text{id} _ x - j \circ q) \circ j = 0$ である。
$j$ が $\text{id} _ {x} - j \circ q$ の核の普遍性を満たすことを示す。
すなわち、$(\text{id} _ x - j \circ q) \circ a = 0$ を満たす射
$a : z \to x$ が与えられたとき、$a = j \circ b$ を満たす一意な射
$b : z \to y$ が存在することを示せばよい。存在は
$a = j \circ q \circ a$ から従う（$b = q \circ a$ と取ればよい）。
一意性は $b = q \circ j \circ b = q \circ a$ から従う
（したがって $a$ が $b$ を決定する）。

\medskip\noindent
双対性により (2) は (1) と同様に従う。
\end{proof}

\begin{lemma}
\label{homology-lemma-split-morphism-splitting}
$\mathcal{A}$ を前加法圏、$x$ と $y$ を $\mathcal{A}$ の対象とし、
$q \circ j = \text{id} _ y$ を満たす射 $j : y \to x$ と $q : x \to y$ を
取る。$j \circ q$ が核または余核の少なくとも一方をもつなら、
両方の（余）核が存在し、
\begin{align*}
x
& \simeq \Ker (j \circ q) \oplus y \\
& \simeq \Coker (j \circ q) \oplus y
\end{align*}
である。いずれの直積の構造にも $j , q$ と標準的な（余）核射が含まれ、
残りの射は一意に定まる。
\end{lemma}

\begin{proof}
$f = j \circ q$ と置き、$f \circ f = f$ を計算して、
補題 \ref{homology-lemma-idempotent-splitting} と
補題 \ref{homology-lemma-split-morphism-kernel-cokernel} を適用する。
\end{proof}




\section{カルービ圏}
\label{homology-section-karoubian}

\noindent
初読ではこの節を飛ばしてよい。

\begin{definition}
\label{homology-definition-karoubian}
$\mathcal{C}$ を前加法圏とする。$\mathcal{C}$ の各対象のすべての
冪等自己準同型が核をもつとき、$\mathcal{C}$ は {\it カルービ圏} であるという。
\end{definition}

\noindent
双対的な概念は、対象のすべての冪等自己準同型が余核をもつことになる。
しかし、次の補題の双対を考えれば、これは同値な概念である。

\begin{lemma}
\label{homology-lemma-karoubian}
$\mathcal{C}$ を前加法圏とする。次は同値である。
\begin{enumerate}
\item $\mathcal{C}$ はカルービ圏である。
\item $\mathcal{C}$ の対象のすべての冪等自己準同型は余核をもつ。
\item $\mathcal{C}$ の冪等自己準同型 $p : z \to z$ が与えられると、
$p$ が $y$ への射影に対応するような直和分解 $z = x \oplus y$ が存在する。
\end{enumerate}
\end{lemma}

\begin{proof}
補題 \ref{homology-lemma-idempotent-kernel-cokernel}、
\ref{homology-lemma-idempotent-splitting}、および
\ref{homology-lemma-additive-cat-biproduct-kernel} から直ちに従う。
\end{proof}

\begin{lemma}
\label{homology-lemma-projectors-have-images}
$\mathcal{D}$ を前加法圏とする。
\begin{enumerate}
\item $\mathcal{D}$ が可算積と、右逆射をもつ射の核をもつなら、
$\mathcal{D}$ はカルービ圏である。
\item $\mathcal{D}$ が可算余積と、左逆射をもつ射の余核をもつなら、
$\mathcal{D}$ はカルービ圏である。
\end{enumerate}
\end{lemma}

\begin{proof}
$X$ を $\mathcal{D}$ の対象とし、$e : X \to X$ を冪等射とする。関手
$$
W \longmapsto \Ker(
\Mor_\mathcal{D}(W, X)
\xrightarrow{e}
\Mor_\mathcal{D}(W, X)
)
$$
が表現可能であることと、$e$ が核をもつことは同値である。
任意のアーベル群 $A$ と冪等自己準同型 $e : A \to A$ に対して
$$
\Ker(e : A \to A)
= \Ker(\Phi :
\prod\nolimits_{n \in \mathbf{N}} A
\to
\prod\nolimits_{n \in \mathbf{N}} A
)
$$
であることに注意する。ここで
$$
\Phi(a_1, a_2, a_3, \ldots) = (ea_1 + (1 - e)a_2, ea_2 + (1 - e)a_3, \ldots)
$$
である。さらに、$\Phi$ は右逆射
$$
\Psi(a_1, a_2, a_3, \ldots) =
(a_1, (1 - e)a_1 + ea_2, (1 - e)a_2 + ea_3, \ldots).
$$
をもつ。したがって (1) が成り立つ。(2) の証明は双対である
（カルービ圏の双対的な定義、すなわち補題 \ref{homology-lemma-karoubian} の条件 (2) を
用いる）。
\end{proof}







\section{アーベル圏}
\label{homology-section-abelian-categories}

\noindent
アーベル圏とは、蛇の補題が成り立つためにちょうど十分な公理を満たす圏である。
（ときに忘れられる）公理の一つは、補題 \ref{homology-lemma-coim-im-map} の標準射
$\Coim(f) \to \Im(f)$ が常に同型であるというものである。
例 \ref{homology-example-not-abelian} は、この公理が必要であることを示している。

\begin{definition}
\label{homology-definition-abelian-category}
圏 $\mathcal{A}$ が {\it アーベル圏} であるとは、それが加法圏であり、
すべての核と余核が存在し、さらに $\mathcal{A}$ の任意の射 $f$ に対する
自然な写像 $\Coim(f) \to \Im(f)$ が同型であることをいう。
\end{definition}

\begin{lemma}
\label{homology-lemma-abelian-opposite}
$\mathcal{A}$ を前加法圏とする。射集合上の加法により
$\mathcal{A}^{opp}$ は前加法圏となる。さらに、$\mathcal{A}$ が加法圏である
ための必要十分条件は $\mathcal{A}^{opp}$ が加法圏であることであり、
$\mathcal{A}$ がアーベル圏であるための必要十分条件は
$\mathcal{A}^{opp}$ がアーベル圏であることである。
\end{lemma}

\begin{proof}
最初の主張は直ちに分かる。$\mathcal{A}$ が加法圏であることと
$\mathcal{A}^{opp}$ が加法圏であることが同値であることを見るには、加法性が
零対象と直和の存在によって特徴づけられ、どちらも反対圏に移っても保たれる
ことを思い出せばよい。最後に、$\mathcal{A}$ がアーベル圏であることと
$\mathcal{A}^{opp}$ がアーベル圏であることが同値であることは、
$\mathcal{A}^{opp}$ における核、余核、像、余像が、それぞれ
$\mathcal{A}$ における余核、核、余像、像に対応することから従う。
\end{proof}

\begin{definition}
\label{homology-definition-injective-surjective}
$f : x \to y$ をアーベル圏の射とする。
\begin{enumerate}
\item $\Ker(f) = 0$ のとき $f$ は {\it 単射} であるという。
\item $\Coker(f) = 0$ のとき $f$ は {\it 全射} であるという。
\end{enumerate}
$x \to y$ が単射なら、$x$ は $y$ の {\it 部分対象} であるといい、
$x \subset y$ と書く。$x \to y$ が全射なら、$y$ は $x$ の
{\it 商} であるという。
\end{definition}

\begin{lemma}
\label{homology-lemma-characterize-injective}
$f : x \to y$ をアーベル圏 $\mathcal{A}$ の射とする。このとき
\begin{enumerate}
\item $f$ が単射であることと $f$ がモノ射であることは同値である。
\item $f$ が全射であることと $f$ がエピ射であることは同値である。
\item $f$ が同型であることと $f$ が単射かつ全射であることは同値である。
\end{enumerate}
\end{lemma}

\begin{proof}
(1) を証明する。定義 \ref{homology-definition-kernel} を参照すると、$\Ker(f)$ は
$z$ を
$\Ker(\Mor_\mathcal{A}(z, x) \to \Mor_\mathcal{A}(z, y))$ に写す関手を
表現する対象である。したがって、$\Ker(f)$ が $0$ であること、任意の $z$ に
対して $\Mor_\mathcal{A}(z, x) \to \Mor_\mathcal{A}(z, y)$ が単射であること、
および $f$ がモノ射であることは同値である。(2) の証明も同様である。
(3) について、同型はモノ射かつエピ射なので、(1), (2) により $f$ は
単射かつ全射である。逆に、$f$ が単射かつ全射なら
$x = \Coim(f)$ および $y = \Im(f)$ であるから、$f$ は同型である。
\end{proof}

\noindent
アーベル圏において $x \subset y$ が部分対象であるとき、
$$
y/x = \Coker(x \to y).
$$
と書く。

\begin{lemma}
\label{homology-lemma-colimit-abelian-category}
$\mathcal{A}$ をアーベル圏とする。$\mathcal{A}$ ではすべての有限極限と
有限余極限が存在する。
\end{lemma}

\begin{proof}
有限極限が存在することを示すには、有限積と等化子が存在することを示せば
十分である。Categories, Lemma \ref{categories-lemma-finite-limits-exist}
を参照されたい。有限積は定義により存在し、$a, b : x \to y$ の等化子は
$a - b$ の核である。有限余極限についても双対的な同様の議論が成り立つ。
\end{proof}

\begin{example}
\label{homology-example-fibre-product-pushouts}
$\mathcal{A}$ をアーベル圏とする。$\mathcal{A}$ における押し出しと
ファイバー積は次のように簡単に記述できる。
\begin{enumerate}
\item $a : x \to y$, $b : z \to y$ が $\mathcal{A}$ の射なら、
ファイバー積は
$x \times_y z = \Ker((a, -b) : x \oplus z \to y)$ である。
\item $a : y \to x$, $b : y \to z$ が $\mathcal{A}$ の射なら、
押し出しは
$x \amalg_y z = \Coker((a, -b) : y \to x \oplus z)$ である。
\end{enumerate}
\end{example}

\begin{definition}
\label{homology-definition-exact}
$\mathcal{A}$ を加法圏とし、$\mathcal{A}$ における射の列
$$
\ldots \to x \to y \to z \to \ldots
\quad\text{or}\quad
x_1 \to x_2 \to \ldots \to x_n
$$
を考える。図示された任意の連続する二つの矢印の合成が零であるとき、この列を
{\it 複体} という。$\mathcal{A}$ がアーベル圏であるとする。上の第一の型の
複体が $y$ において {\it 完全} であるとは
$\Im(x \to y) = \Ker(y \to z)$ であることをいう。第二の型の複体が、
$1 < i < n$ なる $x_i$ において {\it 完全} であるとは
$\Im(x_{i - 1} \to x_i) = \Ker(x_i \to x_{i + 1})$ であることをいう。
上のような列が複体であり、第一の場合には各対象で、第二の場合にはすべての
$1 < i < n$ に対する $x_i$ で完全であるとき、その列は {\it 完全} である、
{\it 完全列} である、または {\it 完全複体} であるという。
これらの概念には、次の形の列に対する変種もある。
$$
\ldots \to x_{-3} \to x_{-2} \to x_{-1}
\quad\text{and}\quad
x_1 \to x_2 \to x_3 \to \ldots
$$
{\it 短完全列} とは、次の形の完全複体である。
$$
0 \to A  \to B \to C \to 0.
$$
\end{definition}

\noindent
次の補題では、アーベル群の列が完全であることの意味を読者が知っているものと
仮定する。

\begin{lemma}
\label{homology-lemma-check-exactness}
$\mathcal{A}$ をアーベル圏とし、
$0 \to M_1 \to M_2 \to M_3 \to 0$ を $\mathcal{A}$ の複体とする。
\begin{enumerate}
\item $M_1 \to M_2 \to M_3 \to 0$ が完全であるための必要十分条件は、
$\mathcal{A}$ の任意の対象 $N$ に対して
$$
0 \to \Hom_\mathcal{A}(M_3, N) \to
\Hom_\mathcal{A}(M_2, N) \to \Hom_\mathcal{A}(M_1, N)
$$
がアーベル群の完全列であることである。
\item $0 \to M_1 \to M_2 \to M_3$ が完全であるための必要十分条件は、
$\mathcal{A}$ の任意の対象 $N$ に対して
$$
0 \to \Hom_\mathcal{A}(N, M_1) \to \Hom_\mathcal{A}(N, M_2) \to
\Hom_\mathcal{A}(N, M_3)
$$
がアーベル群の完全列であることである。
\end{enumerate}
\end{lemma}

\begin{proof}
省略する。ヒント：Algebra, Lemma \ref{algebra-lemma-hom-exact} を参照。
\end{proof}

\begin{definition}
\label{homology-definition-ses-split}
$\mathcal{A}$ をアーベル圏とする。$i : A \to B$ と $q : B \to C$ を
$\mathcal{A}$ の射とし、$0 \to A \to B \to C \to 0$ が短完全列であると
する。$(B, i, j, p, q)$ が $A$ と $C$ の直和となるような射
$j : C \to B$ と $p : B \to A$ が存在するとき、この短完全列は
{\it 分裂する} という。
\end{definition}

\begin{lemma}
\label{homology-lemma-ses-split}
$\mathcal{A}$ をアーベル圏とし、
$0 \to A \to B \to C \to 0$ を短完全列とする。
\begin{enumerate}
\item $B \to C$ の右逆射 $s : C \to B$ が与えられると、
定義 \ref{homology-definition-ses-split} の意味で $(s, \pi)$ が短完全列を分裂させる
ような一意な $\pi : B \to A$ が存在する。
\item $A \to B$ の左逆射 $\pi : B \to A$ が与えられると、
定義 \ref{homology-definition-ses-split} の意味で $(s, \pi)$ が短完全列を分裂させる
ような一意な $s : C \to B$ が存在する。
\end{enumerate}
\end{lemma}

\begin{proof}
補題 \ref{homology-lemma-split-morphism-splitting} の直ちに従う系である。
\end{proof}

\begin{lemma}
\label{homology-lemma-characterize-cartesian}
$\mathcal{A}$ をアーベル圏とする。可換図式
$$
\xymatrix{
w\ar[r]^f\ar[d]_g
& y\ar[d]^h\\
x\ar[r]^k
& z
}
$$
を考える。
\begin{enumerate}
\item この図式がデカルト図式であるための必要十分条件は
$$
0 \to w \xrightarrow{(g, f)} x \oplus y \xrightarrow{(k, -h)} z
$$
が完全であることである。
\item この図式が余デカルト図式であるための必要十分条件は
$$
w \xrightarrow{(g, -f)} x \oplus y \xrightarrow{(k, h)} z \to 0
$$
が完全であることである。
\end{enumerate}
\end{lemma}

\begin{proof}
$u = (g, f) : w \to x \oplus y$ および
$v = (k, -h) : x \oplus y \to z$ と置く。標準射影を
$p : x \oplus y \to x$ と $q : x \oplus y \to y$ とし、標準単射を
$i : \Ker(v) \to x \oplus y$ とする。
例 \ref{homology-example-fibre-product-pushouts} により、この図式がデカルト図式で
あるための必要十分条件は、$f \circ r = q \circ i$ および
$g \circ r = p \circ i$ を満たす同型 $r : \Ker(v) \to w$ が存在することで
ある。列 $0 \to w \overset{u} \to x \oplus y \overset{v} \to z$ が完全で
あるための必要十分条件は、$u \circ r = i$ を満たす同型
$r : \Ker(v) \to w$ が存在することである。一方、
$r : \Ker(v) \to w$ が与えられたとき、
$f \circ r = q \circ i$ および $g \circ r = p \circ i$ であることは、
$q \circ u \circ r= f \circ r = q \circ i$ および
$p \circ u \circ r = g \circ r = p \circ i$ であることと同値であり、
したがって $u \circ r = i$ であることと同値である。これで (1) が示され、
(2) は双対性から従う。
\end{proof}

\begin{lemma}
\label{homology-lemma-cartesian-kernel}
$\mathcal{A}$ をアーベル圏とし、可換図式
$$
\xymatrix{
w\ar[r]^f\ar[d]_g
& y\ar[d]^h\\
x\ar[r]^k
& z
}
$$
を考える。
\begin{enumerate}
\item 図式がデカルト図式なら、$g$ によって誘導される射
$\Ker(f)\to\Ker(k)$ は同型である。
\item 図式が余デカルト図式なら、$h$ によって誘導される射
$\Coker(f)\to\Coker(k)$ は同型である。
\end{enumerate}
\end{lemma}

\begin{proof}
図式がデカルト図式であるとする。$g$ によって誘導される射を
$e:\Ker(f)\to\Ker(k)$ とし、標準単射を
$i:\Ker(f)\to w$ および $j:\Ker(k)\to x$ とする。
$f\circ t=0$ と $g\circ t=j$ を満たす $t:\Ker(k)\to w$ が存在する。
したがって、$i\circ u=t$ を満たす $u:\Ker(k)\to\Ker(f)$ が存在する。
このとき
$g\circ i\circ u\circ e=g\circ t\circ e=j\circ e=g\circ i$ かつ
$f\circ i\circ u\circ e=0=f\circ i$ なので $i\circ u\circ e=i$ である。
$i$ はモノ射だから $u\circ e=\text{id}_{\Ker(f)}$ を得る。
さらに $j\circ e\circ u=g\circ i\circ u=g\circ t=j$ である。
$j$ はモノ射だから $e\circ u=\text{id}_{\Ker(k)}$ を得る。
これで (1) が示された。(2) は双対性から従う。
\end{proof}

\begin{lemma}
\label{homology-lemma-cartesian-cocartesian}
$\mathcal{A}$ をアーベル圏とし、可換図式
$$
\xymatrix{
w\ar[r]^f\ar[d]_g
& y\ar[d]^h\\
x\ar[r]^k
& z
}
$$
を考える。
\begin{enumerate}
\item 図式がデカルト図式で $k$ がエピ射なら、図式は余デカルト図式であり、
$f$ はエピ射である。
\item 図式が余デカルト図式で $g$ がモノ射なら、図式はデカルト図式であり、
$h$ はモノ射である。
\end{enumerate}
\end{lemma}

\begin{proof}
図式がデカルト図式であり、$k$ がエピ射であるとする。
$u = (g, f) : w \to x \oplus y$ と置き、
$v = (k, -h) : x \oplus y \to z$ と置く。$k$ がエピ射なので $v$ もエピ射で
ある。したがって補題 \ref{homology-lemma-characterize-cartesian} により、列
$0\to w\overset{u}\to x\oplus y\overset{v}\to z\to 0$ は完全である。
よって補題 \ref{homology-lemma-characterize-cartesian} により図式は余デカルト図式で
ある。最後に、補題 \ref{homology-lemma-cartesian-kernel} と
補題 \ref{homology-lemma-characterize-injective} により $f$ はエピ射である。
これで (1) が示され、(2) は双対性から従う。
\end{proof}






\begin{lemma}
\label{homology-lemma-epimorphism-universal-abelian-category}
$\mathcal{A}$ をアーベル圏とする。
\begin{enumerate}
\item $x \to y$ が全射なら、任意の $z \to y$ に対して射影
$x \times_y z \to z$ は全射である。
\item $x \to y$ が単射なら、任意の $x \to z$ に対して射
$z \to z \amalg_x y$ は単射である。
\end{enumerate}
\end{lemma}

\begin{proof}
補題 \ref{homology-lemma-characterize-injective} と
補題 \ref{homology-lemma-cartesian-cocartesian} から直ちに従う。
\end{proof}

\begin{lemma}
\label{homology-lemma-check-exactness-fibre-product}
$\mathcal{A}$ をアーベル圏とする。$f:x\to y$ および $g:y\to z$ を
$g\circ f=0$ を満たす射とする。このとき次は同値である。
\begin{enumerate}
\item 列 $x\overset{f}\to y\overset{g}\to z$ は完全である。
\item $g\circ h=0$ を満たす任意の $h:w\to y$ に対して、対象 $v$、
エピ射 $k:v\to w$、および $h\circ k=f\circ l$ を満たす射
$l:v\to x$ が存在する。
\end{enumerate}
\end{lemma}

\begin{proof}
標準単射を $i:\Ker(g)\to y$、標準射影を $p:x\to\Coim(f)$、
標準単射を $j:\Im(f)\to\Ker(g)$ とする。

\medskip\noindent
(1) が成り立つとする。$g\circ h=0$ を満たす $h:w\to y$ を取る。
$i\circ c=h$ を満たす $c:w\to\Ker(g)$ が存在する。
$v=x\times_{\Ker(g)}w$ とし、その標準射影を $k:v\to w$ および
$l:v\to x$ とすると、$c\circ k=j\circ p\circ l$ である。したがって
$h\circ k=i\circ c\circ k=i\circ j\circ p\circ l=f\circ l$ である。
仮定により $j\circ p$ はエピ射なので、補題
\ref{homology-lemma-cartesian-cocartesian} により $k$ はエピ射である。
よって (2) が従う。

\medskip\noindent
(2) が成り立つとする。このとき $g\circ i=0$ である。したがって対象 $w$、
エピ射 $k:w\to\Ker(g)$、および $f\circ l=i\circ k$ を満たす射
$l:w\to x$ が存在する。すると
$i\circ j\circ p\circ l=f\circ l=i\circ k$ である。$i$ はモノ射なので
$j\circ p\circ l=k$ はエピ射である。よって $j$ はエピ射であり、したがって
同型である。これから (1) が従う。
\end{proof}

\begin{lemma}
\label{homology-lemma-exact-kernel-sequence}
$\mathcal{A}$ をアーベル圏とし、可換図式
$$
\xymatrix{
x \ar[r]^f \ar[d]^\alpha &
y \ar[r]^g \ar[d]^\beta &
z \ar[d]^\gamma\\
u \ar[r]^k & v \ar[r]^l & w
}
$$
を考える。
\begin{enumerate}
\item 第一行が完全で $k$ がモノ射なら、誘導される列
$\Ker(\alpha) \to \Ker(\beta) \to \Ker(\gamma)$ は完全である。
\item 第二行が完全で $g$ がエピ射なら、誘導される列
$\Coker(\alpha) \to \Coker(\beta) \to \Coker(\gamma)$ は完全である。
\end{enumerate}
\end{lemma}

\begin{proof}
第一行が完全で $k$ がモノ射であるとする。誘導される射を
$a:\Ker(\alpha)\to\Ker(\beta)$ および
$b:\Ker(\beta)\to\Ker(\gamma)$ とする。標準単射を
$h:\Ker(\alpha)\to x$、$i:\Ker(\beta)\to y$、
$j:\Ker(\gamma)\to z$ とする。$j$ はモノ射なので $b\circ a=0$ である。
$b\circ c=0$ を満たす $c:s\to\Ker(\beta)$ を取る。このとき
$g\circ i\circ c=j\circ b\circ c=0$ である。
補題 \ref{homology-lemma-check-exactness-fibre-product} により、対象 $t$、
エピ射 $d:t\to s$、および $i\circ c\circ d=f\circ e$ を満たす射
$e:t\to x$ が存在する。すると
$k\circ \alpha\circ e=\beta\circ f\circ e=\beta\circ i\circ c\circ d=0$
である。$k$ はモノ射なので $\alpha\circ e=0$ を得る。
したがって $h\circ m=e$ を満たす $m:t\to\Ker(\alpha)$ が存在する。
すると $i\circ a\circ m=f\circ h\circ m=f\circ e=i\circ c\circ d$ である。
$i$ はモノ射なので $a\circ m=c\circ d$ を得る。ゆえに補題
\ref{homology-lemma-check-exactness-fibre-product} から (1) が従い、(2) は双対性から
従う。
\end{proof}

\begin{lemma}
\label{homology-lemma-snake}
$\mathcal{A}$ をアーベル圏とし、各行が完全である可換図式
$$
\xymatrix{
& x \ar[r]^f \ar[d]^\alpha &
y \ar[r]^g \ar[d]^\beta &
z \ar[r] \ar[d]^\gamma &
0 \\
0 \ar[r] & u \ar[r]^k & v \ar[r]^l & w
}
$$
を考える。
\begin{enumerate}
\item 図式
$$
\xymatrix{
y \ar[d]_\beta &
y \times_z \Ker(\gamma) \ar[l]_{\pi'} \ar[r]^{\pi} &
\Ker(\gamma) \ar[d]^\delta \\
v \ar[r]^{\iota'} & \Coker(\alpha) \amalg_u v &
\Coker(\alpha) \ar[l]_\iota
}
$$
が可換となる一意な射 $\delta : \Ker(\gamma) \to \Coker(\alpha)$ が存在する。
ここで $\pi$ と $\pi'$ は標準射影、$\iota$ と $\iota'$ は標準余射入である。
\item 誘導される列
$$
\Ker(\alpha) \xrightarrow{f'} \Ker(\beta) \xrightarrow{g'}
\Ker(\gamma) \xrightarrow{\delta} \Coker(\alpha) \xrightarrow{k'}
\Coker(\beta) \xrightarrow{l'} \Coker(\gamma)
$$
は完全である。$f$ が単射なら $f'$ も単射であり、$l$ が全射なら
$l'$ も全射である。
\end{enumerate}
\end{lemma}

\begin{proof}
補題 \ref{homology-lemma-cartesian-cocartesian} により $\pi$ はエピ射、$\iota$ は
モノ射なので、$\delta$ の一意性は明らかである。
$p = y \times_z \Ker(\gamma)$ および
$q = \Coker(\alpha) \amalg_u v$ と置く。標準単射を
$h : \Ker(\beta) \to y$、$i : \Ker(\gamma) \to z$、
$j : \Ker(\pi) \to p$ とする。また、標準射影を
$\pi'' : u \to \Coker(\alpha)$ とする。補題
\ref{homology-lemma-cartesian-cocartesian} を念頭に置くと、完全な行をもつ可換図式
$$
\xymatrix{
0 \ar[r] &
\Ker(\pi) \ar[r]^j &
p \ar[r]^{\pi} \ar[d]_{\pi'} &
\Ker(\gamma) \ar[d]_i \ar[r] & 0 \\
& x \ar[r]^f \ar[d]_\alpha & y \ar[r]^g \ar[d]_\beta &
z \ar[d]_\gamma \ar[r] & 0 \\
0 \ar[r] & u \ar[r]^k \ar[d]_{\pi''} &
v \ar[r]^l \ar[d]_{\iota'} & w & \\
0 \ar[r] & \Coker(\alpha) \ar[r]^\iota & q & &
}
$$
を得る。$l \circ \beta \circ \pi' = \gamma \circ i \circ \pi = 0$ であり、
上の図式の第三行は完全なので、$k \circ a = \beta \circ \pi'$ を満たす
$a : p \to u$ が存在する。上右の四角形はデカルト図式なので、補題
\ref{homology-lemma-cartesian-kernel} により、$\pi' \circ j \circ b = f$ を満たす
エピ射 $b : x \to \Ker(\pi)$ が存在する。したがって
$k \circ a \circ j \circ b = \beta \circ \pi' \circ j \circ b =
\beta \circ f = k \circ \alpha$ である。$k$ はモノ射なので
$a \circ j \circ b = \alpha$ を得る。よって
$\pi'' \circ a \circ j \circ b = \pi'' \circ \alpha = 0$ である。
$b$ はエピ射なので、これは $\pi'' \circ a \circ j = 0$ を含意する。
したがって、上の図式の第一行が完全であることから
$\delta \circ \pi = \pi'' \circ a$ を満たす
$\delta : \Ker(\gamma) \to \Coker(\alpha)$ が存在する。すると
$\iota \circ \delta \circ \pi = \iota \circ \pi'' \circ a =
\iota' \circ k \circ a = \iota' \circ \beta \circ \pi'$ となり、
求める可換性が得られる。

\medskip\noindent
上の図式の上右の四角形はデカルト図式なので、
$\pi' \circ c = h$ および $\pi \circ c = g'$ を満たす
$c : \Ker(\beta) \to p$ が存在する。したがって
$\iota \circ \delta \circ g' = \iota \circ \delta \circ \pi \circ c =
\iota' \circ \beta \circ \pi' \circ c = \iota' \circ \beta \circ h = 0$
である。$\iota$ はモノ射なので $\delta \circ g' = 0$ を得る。

\medskip\noindent
次に $\delta \circ d = 0$ を満たす $d : r \to \Ker(\gamma)$ を取る。
完全列 $p \xrightarrow{\pi} \Ker(\gamma) \to 0$ と $d$ に補題
\ref{homology-lemma-check-exactness-fibre-product} を適用すると、対象 $s$、
エピ射 $m : s \to r$、および $\pi \circ n = d \circ m$ を満たす射
$n : s \to p$ を得る。$\pi'' \circ a \circ n = \delta \circ d \circ m = 0$
なので、完全列 $x \xrightarrow{\alpha} u \xrightarrow{p} \Coker(\alpha)$ と
$a \circ n$ に補題 \ref{homology-lemma-check-exactness-fibre-product} を適用すると、
対象 $t$、エピ射 $\varepsilon : t \to s$、および
$a \circ n \circ \varepsilon = \alpha \circ \zeta$ を満たす射
$\zeta : t \to x$ を得る。このとき
$\beta \circ \pi' \circ n \circ \varepsilon =
k \circ \alpha \circ \zeta = \beta \circ f \circ \zeta$ である。
$\eta = \pi' \circ n \circ \varepsilon - f \circ \zeta : t \to y$ と置くと、
$\beta \circ \eta = 0$ である。したがって
$\eta = h \circ \vartheta$ を満たす $\vartheta : t \to \Ker(\beta)$ が存在する。
また
$i \circ g' \circ \vartheta = g \circ h \circ \vartheta =
g \circ \pi' \circ n \circ \varepsilon - g \circ f \circ \zeta =
i \circ \pi \circ n \circ \varepsilon = i \circ d \circ m \circ \varepsilon$
である。$i$ はモノ射なので
$g' \circ \vartheta = d \circ m \circ \varepsilon$ を得る。
$m \circ \varepsilon$ はエピ射であるから、補題
\ref{homology-lemma-check-exactness-fibre-product} により
$\Ker(\beta) \xrightarrow{g'} \Ker(\gamma) \xrightarrow{\delta} \Coker(\alpha)$
は完全である。あとは補題 \ref{homology-lemma-exact-kernel-sequence} と双対性から主張が
従う。
\end{proof}

\begin{lemma}
\label{homology-lemma-snake-natural}
$\mathcal{A}$ をアーベル圏とし、完全な行をもつ可換図式
$$
\xymatrix{
& & & x\ar[ld]\ar[rr]\ar[dd]^(.4)\alpha
& & y\ar[ld]\ar[rr]\ar[dd]^(.4)\beta
& & z\ar[ld]\ar[rr]\ar[dd]^(.4)\gamma
& & 0\\
& & x'\ar[rr]\ar[dd]^(.4){\alpha'}
& & y'\ar[rr]\ar[dd]^(.4){\beta'}
& & z'\ar[rr]\ar[dd]^(.4){\gamma'}
& & 0
& \\
& 0\ar[rr]
& & u\ar[ld]\ar[rr]
& & v\ar[ld]\ar[rr]
& & w\ar[ld]
& & \\
0\ar[rr]
& & u'\ar[rr]
& & v'\ar[rr]
& & w'
& & &
}
$$
を考える。このとき誘導される図式
$$
\xymatrix@C=15pt{
\Ker(\alpha) \ar[r] \ar[d] &
\Ker(\beta) \ar[r] \ar[d] &
\Ker(\gamma) \ar[r]^(.45){\delta} \ar[d] &
\Coker(\alpha) \ar[r] \ar[d] &
\Coker(\beta) \ar[r] \ar[d] &
\Coker(\gamma) \ar[d] \\
\Ker(\alpha') \ar[r] &
\Ker(\beta') \ar[r] &
\Ker(\gamma') \ar[r]^(.45){\delta'} &
\Coker(\alpha') \ar[r] &
\Coker(\beta') \ar[r] &
\Coker(\gamma')
}
$$
は可換である。
\end{lemma}

\begin{proof}
省略する。
\end{proof}

\begin{lemma}
\label{homology-lemma-four-lemma}
$\mathcal{A}$ をアーベル圏とし、完全な行をもつ可換図式
$$
\xymatrix{
w \ar[r] \ar[d]^\alpha & x \ar[r] \ar[d]^\beta & y \ar[r] \ar[d]^\gamma &
z \ar[d]^\delta \\
w' \ar[r] & x' \ar[r] & y' \ar[r] & z'
}
$$
を考える。
\begin{enumerate}
\item $\alpha, \gamma$ が全射で $\delta$ が単射なら、$\beta$ は全射である。
\item $\beta, \delta$ が単射で $\alpha$ が全射なら、$\gamma$ は単射である。
\end{enumerate}
\end{lemma}

\begin{proof}
$\alpha, \gamma$ が全射で $\delta$ が単射であるとする。
$w'$ を $\Im(w' \to x')$ で置き換えることにより、$w' \to x'$ が単射であると
仮定してよい。$z$ を $\Im(y \to z)$ で置き換えることにより、
$y \to z$ が全射であると仮定してよい。そこで補題 \ref{homology-lemma-snake} を
$$
\xymatrix{
& \Ker(y \to z) \ar[r] \ar[d] & y \ar[r] \ar[d] & z \ar[r] \ar[d] & 0 \\
0 \ar[r] & \Ker(y' \to z') \ar[r] & y' \ar[r] & z'
}
$$
に適用すると、$\Ker(y \to z) \to \Ker(y' \to z')$ が全射であると分かる。
最後に補題 \ref{homology-lemma-snake} を
$$
\xymatrix{
& w \ar[r] \ar[d] & x \ar[r] \ar[d] & \Ker(y \to z) \ar[r] \ar[d] & 0 \\
0 \ar[r] & w' \ar[r] & x' \ar[r] & \Ker(y' \to z')
}
$$
に適用すると、$x \to x'$ が全射であると分かる。これで (1) が示された。
(2) の証明は双対である。
\end{proof}

\begin{lemma}
\label{homology-lemma-five-lemma}
\begin{reference}
\cite[Lemma 4.5 page 16]{Eilenberg-Steenrod}
\end{reference}
$\mathcal{A}$ をアーベル圏とし、完全な行をもつ可換図式
$$
\xymatrix{
v \ar[r] \ar[d]^\alpha &
w \ar[r] \ar[d]^\beta &
x \ar[r] \ar[d]^\gamma &
y \ar[r] \ar[d]^\delta &
z \ar[d]^\epsilon \\
v' \ar[r] & w' \ar[r] & x' \ar[r] & y' \ar[r] & z'
}
$$
を考える。$\beta, \delta$ が同型、$\epsilon$ が単射、$\alpha$ が全射なら、
$\gamma$ は同型である。
\end{lemma}

\begin{proof}
補題 \ref{homology-lemma-four-lemma} から直ちに従う。
\end{proof}







\section{拡大}
\label{homology-section-extensions}

\begin{definition}
\label{homology-definition-extension}
$\mathcal{A}$ をアーベル圏とし、$A, B \in \Ob(\mathcal{A})$ とする。
$A$ による $B$ の {\it 拡大 $E$} とは、短完全列
$$
0 \to A \to E \to B \to 0.
$$
のことである。二つの拡大 $0 \to A \to E \to B \to 0$ と
$0 \to A \to F \to B \to 0$ の間の {\it 拡大の射} とは、図式
$$
\xymatrix{
0 \ar[r] &
A \ar[r] \ar[d]^{\text{id}} &
E \ar[r] \ar[d]^f &
B \ar[r] \ar[d]^{\text{id}} &
0 \\
0 \ar[r] &
A \ar[r] &
F \ar[r] &
B \ar[r] &
0
}
$$
を可換にする $\mathcal{A}$ の射 $f : E \to F$ のことである。
したがって、$A$ による $B$ の拡大は圏をなす。
\end{definition}

\noindent
言葉の濫用により、射 $A \to E$ と $E \to B$ への言及をしばしば省くが、
これらは明らかに拡大の構造の一部である。

\begin{definition}
\label{homology-definition-ext-group}
$\mathcal{A}$ をアーベル圏とし、$A, B \in \Ob(\mathcal{A})$ とする。
$A$ による $B$ の拡大の同型類の集合を
$$
\Ext_\mathcal{A}(B, A).
$$
と書き、{\it $\Ext$ 群} という。
\end{definition}

\noindent
この定義が意味をもつのは、我々の規約では $\Ob(\mathcal{A})$ が集合であり、
したがって $\Ext_\mathcal{A}(B, A)$ も集合だからである。
Categories, Remark \ref{categories-remark-big-categories} に挙げられた
「大きな」アーベル圏の各場合にも、$\Ext_\mathcal{A}(B, A)$ が集合であることを
直接確認できる。また、$\mathcal{A}$ が十分な射影的対象または十分な
入射的対象をもつ場合には、常にそうであることを後で見る。
将来の参照をここに挿入する。

\medskip\noindent
実際、$\Ext_\mathcal{A}(-, -)$ を次の関手にすることができる。
$$
\mathcal{A} \times \mathcal{A}^{opp} \longrightarrow \textit{Sets}, \quad
(A, B) \longmapsto \Ext_\mathcal{A}(B, A)
$$
その構成は次のとおりである。
\begin{enumerate}
\item 射 $B' \to B$ と、$A$ による $B$ の拡大 $E$ が与えられたとき、
$E' = E \times_B B'$ と定める。補題 \ref{homology-lemma-cartesian-kernel} により、
短完全列からなる次の可換図式が得られる。
$$
\xymatrix{
0 \ar[r] & A \ar[r] \ar[d] & E' \ar[r] \ar[d] & B' \ar[r] \ar[d] & 0 \\
0 \ar[r] & A \ar[r] & E \ar[r] & B \ar[r] & 0
}
$$
拡大 $E'$ を、$B' \to B$ による $E$ の {\it 引き戻し} という。
\item 射 $A \to A'$ と、$A$ による $B$ の拡大 $E$ が与えられたとき、
$E' = A' \amalg_A E$ と定める。補題 \ref{homology-lemma-cartesian-cocartesian} により、
短完全列からなる次の可換図式が得られる。
$$
\xymatrix{
0 \ar[r] & A \ar[r] \ar[d] & E \ar[r] \ar[d] & B \ar[r] \ar[d] & 0 \\
0 \ar[r] & A' \ar[r] & E' \ar[r] & B \ar[r] & 0
}
$$
拡大 $E'$ を、$A \to A'$ による $E$ の {\it 押し出し} という。
\end{enumerate}
これが上記の関手を定めることを見るには、いくつか確認すべきことがある。
まず、変数 $B$ に関する関手性には
$(E \times_B B') \times_{B'} B'' = E \times_B B''$ が必要であり、これは
ファイバー積の一般的な性質である。変数 $A$ に関する関手性は双対的に扱う。
最後に、$A \to A'$ と $B' \to B$ が与えられたとき、$A'$ による $B'$ の
拡大として
$$
A' \amalg_A (E \times_B B')
\cong
(A' \amalg_A E)\times_B B'
$$
であることを示す必要がある。$A' \amalg_A E$ は $A' \oplus E$ の商である
ことを思い出そう。したがって右辺は $A' \oplus E \times_B B'$ の商であり、
その核が左辺を得るためにまさに必要なものであることは直ちに分かる。

\medskip\noindent
$E_1$ と $E_2$ が $A$ による $B$ の拡大なら、$E_1\oplus E_2$ は
$A\oplus A$ による $B \oplus B$ の拡大である。和写像
$A \oplus A \to A$ によって押し出し、対角写像 $B \to B \oplus B$ によって
引き戻すと、$A$ による $B$ の拡大 $E_1 + E_2$ を得る。
$$
\xymatrix{
0 \ar[r] &
A \oplus A \ar[r] \ar[d]_\Sigma &
E_1 \oplus E_2 \ar[r] \ar[d] &
B \oplus B \ar[r] \ar[d] &
0 \\
0 \ar[r] &
A \ar[r] &
E' \ar[r] &
B \oplus B \ar[r] &
0\\
0 \ar[r] &
A \ar[r] \ar[u] &
E_1 + E_2 \ar[r] \ar[u] &
B \ar[r] \ar[u]^\Delta &
0
}
$$
拡大 $E_1 + E_2$ を、与えられた拡大の {\it ベール和} という。

\begin{lemma}
\label{homology-lemma-baer-sum}
上の構成 $(E_1, E_2) \mapsto E_1 + E_2$ は
$\Ext_\mathcal{A}(B, A)$ 上に可換群構造を定め、両変数に関して関手的である。
\end{lemma}

\begin{proof}
省略する。
\end{proof}

\begin{lemma}
\label{homology-lemma-six-term-sequence-ext}
$\mathcal{A}$ をアーベル圏とし、
$0 \to M_1 \to M_2 \to M_3 \to 0$ を $\mathcal{A}$ の短完全列とする。
\begin{enumerate}
\item $\mathcal{A}$ の任意の対象 $N$ に対して、アーベル群の標準的な
六項完全列
$$
\xymatrix{
0 \ar[r] &
\Hom_\mathcal{A}(M_3, N) \ar[r] &
\Hom_\mathcal{A}(M_2, N) \ar[r] &
\Hom_\mathcal{A}(M_1, N) \ar[lld] \\
& \Ext_\mathcal{A}(M_3, N) \ar[r] &
\Ext_\mathcal{A}(M_2, N) \ar[r] &
\Ext_\mathcal{A}(M_1, N)
}
$$
が存在する。
\item $\mathcal{A}$ の任意の対象 $N$ に対して、アーベル群の標準的な
六項完全列
$$
\xymatrix{
0 \ar[r] &
\Hom_\mathcal{A}(N, M_1) \ar[r] &
\Hom_\mathcal{A}(N, M_2) \ar[r] &
\Hom_\mathcal{A}(N, M_3) \ar[lld] \\
& \Ext_\mathcal{A}(N, M_1) \ar[r] &
\Ext_\mathcal{A}(N, M_2) \ar[r] &
\Ext_\mathcal{A}(N, M_3)
}
$$
が存在する。
\end{enumerate}
\end{lemma}

\begin{proof}
省略する。ヒント：境界写像は、与えられた短完全列の押し出しまたは引き戻しを
用いて定義される。
\end{proof}





\section{加法関手}
\label{homology-section-functors}

\noindent
最初に、加法圏の間の加法関手を特徴づける、ごく単純な補題を述べる。

\begin{lemma}
\label{homology-lemma-additive-functor}
$\mathcal{A}$ と $\mathcal{B}$ を加法圏とし、
$F : \mathcal{A} \to \mathcal{B}$ を関手とする。次の条件は同値である。
\begin{enumerate}
\item $F$ は加法的である。
\item 任意の $A, B \in \mathcal{A}$ に対して
$F(A) \oplus F(B) \to F(A \oplus B)$ は同型である。
\item 任意の $A, B \in \mathcal{A}$ に対して
$F(A \oplus B) \to F(A) \oplus F(B)$ は同型である。
\end{enumerate}
\end{lemma}

\begin{proof}
加法関手は補題 \ref{homology-lemma-additive-additive} により直和と可換するので、
(1) は (2) と (3) を含意する。(2) が成り立ち、$0$ が $\mathcal{A}$ の
零対象を表すとする。写像
$F(0) \oplus F(0) \to F(0 \oplus 0) \cong F(0)$ が同型であることから、
$\mathcal{B}$ の任意の $C$ に対して対角写像
$\Hom(F(0), C) \to \Hom(F(0), C) \oplus \Hom(F(0), C)$ は全単射である。
したがって $F(0)$ は $\mathcal{B}$ の零対象であり、$F$ は任意の二対象間の
零射を零射に写す。(3) が成り立つ場合も同様である。さらに合成
$F(A) \oplus F(B) \to F(A \oplus B) \to F(A) \oplus F(B)$ は行列
$$
\left(
\begin{matrix}
F(1) & F(0) \\
F(0) & F(1)
\end{matrix}
\right)
$$
で与えられ、したがって恒等写像であるから、(2) と (3) は同値である。
(2) と (3) が成り立つと仮定する。射 $f, g : A \to B$ を取る。
このとき $f + g$ は合成
$$
A \to A \oplus A \xrightarrow{\text{diag}(f, g)} B \oplus B \to B
$$
に等しい。関手 $F$ を適用して次の図式を考える。
$$
\xymatrix{
F(A) \ar[r] \ar[rd] &
F(A \oplus A) \ar[rr]_{F(\text{diag}(f, g))} & &
F(B \oplus B) \ar[r] \ar[d] &
F(B) \\
&
F(A) \oplus F(A) \ar[u] \ar[rr]^{\text{diag}(F(f), F(g))} & &
F(B) \oplus F(B) \ar[ru]
}
$$
この図式は可換である。中央の正方形については、誘導される四つの写像
$F(A) \to F(B)$ のそれぞれについて確認でき、これは $F$ が関手であり零射を
零射に写すことから従う。左の三角形については、
$F(A \oplus A) \to F(A) \oplus F(A)$ が同型であることを用いれば、この写像と
合成した後で確認すれば十分であり、その確認は直ちにできる。もう一方の三角形も
双対的である。したがって下側を回る合成は $F(f + g)$ に等しく、結論を得る。
\end{proof}

\noindent
Categories, Definition \ref{categories-definition-exact} では、有限（余）極限をもつ
圏の間の関手について「右完全」「左完全」および「完全」という概念を定義した。
したがって、これらは特にアーベル圏の間の関手に適用できる。

\begin{lemma}
\label{homology-lemma-exact-functor}
$\mathcal{A}$ と $\mathcal{B}$ をアーベル圏とし、
$F : \mathcal{A} \to \mathcal{B}$ を関手とする。
\begin{enumerate}
\item $F$ が左完全または右完全なら、$F$ は加法的である。
\item $F$ が左完全であるための必要十分条件は、任意の短完全列
$0 \to A \to B \to C \to 0$ に対して列
$0 \to F(A) \to F(B) \to F(C)$ が完全であることである。
\item $F$ が右完全であるための必要十分条件は、任意の短完全列
$0 \to A \to B \to C \to 0$ に対して列
$F(A) \to F(B) \to F(C) \to 0$ が完全であることである。
\item $F$ が完全であるための必要十分条件は、任意の短完全列
$0 \to A \to B \to C \to 0$ に対して列
$0 \to F(A) \to F(B) \to F(C) \to 0$ が完全であることである。
\end{enumerate}
\end{lemma}

\begin{proof}
$F$ が左完全、すなわち有限極限と可換するなら、$F$ は積を積に写し、したがって
直和を保つので、補題 \ref{homology-lemma-additive-functor} により $F$ は加法的である。
一方、任意の短完全列 $0 \to A \to B \to C \to 0$ に対して
$0 \to F(A) \to F(B) \to F(C)$ が完全であると仮定する。
$A, B$ を二つの対象とする。このとき、例えば
補題 \ref{homology-lemma-additive-cat-biproduct-kernel} により短完全列
$$
0 \to A \to A \oplus B \to B \to 0
$$
がある。仮定により、可換図式
$$
\xymatrix{
0 \ar[r] &
F(A) \ar[d] \ar[r] &
F(A) \oplus F(B) \ar[r] \ar[d] &
F(B) \ar[d] \ar[r] &
0 \\
0 \ar[r] &
F(A) \ar[r] &
F(A \oplus B) \ar[r] &
F(B)
}
$$
の下段は完全である。よって蛇の補題（補題 \ref{homology-lemma-snake}）により
$F(A) \oplus F(B) \to F(A \oplus B)$ は同型である。したがって、この場合も
$F$ は加法的である。以下の証明では $F$ が加法的であると仮定してよい。

\medskip\noindent
$f : B \to C$ を $B$ から $C$ への射とする。
$0 \to A \to B \to C$ が完全であることは $A = \Ker(f)$ にほかならない。
$f$ の核は明らかに $B$ から $C$ への二つの写像 $f$ と $0$ の等化子である。
したがって $F$ が極限と可換すれば
$F(\Ker(f)) = \Ker(F(f))$ であり、これはちょうど
$0 \to F(A) \to F(B) \to F(C)$ が完全であることを意味する。

\medskip\noindent
逆に、$F$ は加法的であり、任意の短完全列
$0 \to A \to B \to C \to 0$ を完全列
$0 \to F(A) \to F(B) \to F(C)$ に写すと仮定する。加法的であるため、$F$ は
直和、したがって $\mathcal{A}$ における有限積と可換する。ゆえに、有限極限と
可換することを示すには、等化子と可換することを示せば十分である。
アーベル圏の等化子は差の写像の核に等しいので、$F$ が核を取る操作と可換する
ことを示せばよい。$f : A \to B$ を射とする。$f$ を
$A \to I \to B$ と分解し、$f' : A \to I$ は全射、$i : I \to B$ は単射とする。
（これはアーベル圏の定義により可能である。）このとき明らかに
$\Ker(f) = \Ker(f')$ である。また
$0 \to \Ker(f') \to A \to I \to 0$ と
$0 \to I \to B \to B/I \to 0$ は短完全列である。$F$ に課した条件から
$0 \to F(\Ker(f')) \to F(A) \to F(I)$ と
$0 \to F(I) \to F(B) \to F(B/I)$ は完全である。したがって
$F(\Ker(f'))$ は写像 $F(A) \to F(B)$ の核でもあり、結論を得る。

\medskip\noindent
(3) の証明は (2) の証明と同様である。(4) は (2) と (3) を組み合わせればよい。
\end{proof}

\begin{lemma}
\label{homology-lemma-exact-functor-ext}
$\mathcal{A}$ と $\mathcal{B}$ をアーベル圏とし、
$F : \mathcal{A} \to \mathcal{B}$ を完全関手とする。
$\mathcal{A}$ の任意の対象の組 $A, B$ に対して、関手 $F$ はアーベル群準同型
$$
\Ext_\mathcal{A}(B, A)
\longrightarrow
\Ext_\mathcal{B}(F(B), F(A))
$$
を誘導し、これは拡大 $E$ を $F(E)$ に写す。
\end{lemma}

\begin{proof}
省略する。
\end{proof}

\noindent
次の補題は、サイト上のアーベル群の層の圏がアーベル圏であることの証明で用いる。
そこでは関手 $b$ は層化関手である。

\begin{lemma}
\label{homology-lemma-adjoint-get-abelian}
$a : \mathcal{A} \to \mathcal{B}$ および $b : \mathcal{B} \to \mathcal{A}$ を
関手とする。次を仮定する。
\begin{enumerate}
\item $\mathcal{A}$ と $\mathcal{B}$ は加法圏、$a$ と $b$ は加法関手であり、
$a$ は $b$ の右随伴である。
\item $\mathcal{B}$ はアーベル圏であり、$b$ は左完全である。
\item $ba \cong \text{id}_\mathcal{A}$ である。
\end{enumerate}
このとき $\mathcal{A}$ はアーベル圏である。
\end{lemma}

\begin{proof}
$\mathcal{B}$ はアーベル圏なので、補題 \ref{homology-lemma-colimit-abelian-category} により
$\mathcal{B}$ ではすべての有限極限と有限余極限が存在する。$b$ は左随伴なので
$b$ は右完全でもあり、したがって完全である。Categories, Lemma
\ref{categories-lemma-exact-adjoint} を参照されたい。
$\varphi : B_1 \to B_2$ を $\mathcal{B}$ の射とする。特に
$K = \Ker(B_1 \to B_2)$ とすれば、$K$ は $0$ と $\varphi$ の等化子であり、
したがって $bK$ は $0$ と $b\varphi$ の等化子である。それゆえ
$bK$ は $b\varphi$ の核である。
同様に $Q = \Coker(B_1 \to B_2)$ とすれば、$Q$ は $0$ と $\varphi$ の余等化子で
あり、したがって $bQ$ は $0$ と $b\varphi$ の余等化子である。それゆえ
$bQ$ は $b\varphi$ の余核である。よって $\mathcal{A}$ において $b\varphi$ の形をした
任意の射は核と余核をもつ。一方、$ba \cong \text{id}$ なので
$\mathcal{A}$ の任意の射はこの形をしており、$\mathcal{A}$ では核と余核が存在する。
実際、この議論から、射 $\psi : A_1 \to A_2$ に対して
$$
\Ker(\psi) = b\Ker(a\psi),
\quad\text{and}\quad
\Coker(\psi) = b\Coker(a\psi).
$$
が従う。あとは $\Coim(\psi)= \Im(\psi)$ を示さなければならない。
次のように示す。まず、$\mathcal{A}$ は核と余核をもつので、すべての有限極限と
有限余極限をもつことに注意する（補題 \ref{homology-lemma-colimit-abelian-category} の証明を
参照）。したがって Categories, Lemma \ref{categories-lemma-exact-adjoint} により
$a$ は左完全であり、それゆえ核（すなわち等化子）を核に写す。
\begin{align*}
\Coim(\psi)
& =
\Coker(\Ker(\psi) \to A_1)
& \text{by definition} \\
& =
b\Coker(a(\Ker(\psi) \to A_1))
& \text{by formula above} \\
& =
b\Coker(\Ker(a\psi) \to aA_1))
& a\text{ preserves kernels} \\
& =
b\Coim(a\psi)
& \text{by definition} \\
& =
b\Im(a\psi)
& \mathcal{B}\text{ is abelian} \\
& =
b\Ker(aA_2 \to \Coker(a\psi))
& \text{by definition} \\
& =
\Ker(baA_2 \to b\Coker(a\psi))
& b\text{ preserves kernels} \\
& =
\Ker(A_2 \to b\Coker(a\psi))
& ba = \text{id}_\mathcal{A} \\
& =
\Ker(A_2 \to \Coker(\psi))
& \text{by formula above} \\
& =
\Im(\psi)
& \text{by definition}
\end{align*}
したがって補題が成り立つ。
\end{proof}





\section{局所化}
\label{homology-section-localization}

\noindent
この節では、ガブリエル--ジスマン局所化と圏の加法構造との関係を述べる。

\begin{lemma}
\label{homology-lemma-localization-preadditive}
$\mathcal{C}$ を前加法圏とし、$S$ を左または右乗法系とする。
局所化関手 $Q : \mathcal{C} \to S^{-1}\mathcal{C}$ が加法的となるような
$S^{-1}\mathcal{C}$ 上の前加法構造が一意に存在する。
\end{lemma}

\begin{proof}
$S$ が左乗法系である場合を証明する。$S$ が右乗法系である場合は双対的である。
$X, Y$ を $\mathcal{C}$ の対象とし、
$\alpha, \beta : X \to Y$ を $S^{-1}\mathcal{C}$ の射とする。
Categories, Lemma \ref{categories-lemma-morphisms-left-localization} により、
これらを共通分母 $s$ をもつ組 $s^{-1}f, s^{-1}g$ で表すことができる。
$S^{-1}\mathcal{C}$ に前加法構造があり、$Q$ が加法的であると仮定する。このとき
\begin{align*}
s^{-1}f+s^{-1}g
&=Q(s)^{-1}\circ Q(f)+Q(s)^{-1}\circ Q(g)\\
&=Q(s)^{-1}\circ(Q(f)+Q(g))\\
&=Q(s)^{-1}\circ Q(f+g)\\
&=s^{-1}(f+g).
\end{align*}
したがって、$S^{-1}\mathcal{C}$ 上の前加法構造は、存在すれば一意である。
そこで $\alpha + \beta$ を $s^{-1}(f + g)$ の同値類として定め、以下ではこれが
代表元によらず、合成が双線形であることを示す。これが示されれば $Q$ が
加法関手であることは明らかである。

\medskip\noindent
まず上の構成が良定義であることを示す。抽象的には、アーベル群のフィルター余極限が
集合のフィルター余極限と一致することを述べ、Categories, Equation
(\ref{categories-equation-left-localization-morphisms-colimit}) を用いればよい。
次のようにもう少し詳しく確かめることもできる。
$s : Y \to Y_1$ および $f, g : X \to Y_1$ とする。
$\alpha, \beta$ の第二の表示が $(s')^{-1}f', (s')^{-1}g'$ であり、
$s' : Y \to Y_2$ および $f', g' : X \to Y_2$ であるとする。
Categories, Remark \ref{categories-remark-left-localization-morphisms-colimit} により、
射 $s_3 : Y \to Y_3$ と射
$a_1 : Y_1 \to Y_3$, $a_2 : Y_2 \to Y_3$ であって
$a_1 \circ s = s_3 = a_2 \circ s'$、かつ
$a_1 \circ f = a_2 \circ f'$ および $a_1 \circ g = a_2 \circ g'$ を満たすものが
存在する。したがって $s^{-1}(f + g)$ は
\begin{align*}
s_3^{-1}(a_1 \circ (f + g)) & =
s_3^{-1}(a_1 \circ f + a_1 \circ g) \\
& = s_3^{-1}(a_2 \circ f' + a_2 \circ g') \\
& = s_3^{-1}(a_2 \circ (f' + g'))
\end{align*}
と同値であり、これは $(s')^{-1}(f' + g')$ と同値である。

\medskip\noindent
$s : Y \to Y'$ および $f, g : X \to Y'$ を固定し、
$S^{-1}\mathcal{C}$ における射 $X \to Y$ として
$\alpha = s^{-1}f$ および $\beta = s^{-1}g$ であるとする。
合成が双線形であることを示すため、まず $S^{-1}\mathcal{C}$ における射
$\gamma : Y \to Z$ の場合を考える。ある $h : Y \to Z'$ と $S$ に属する
$t : Z \to Z'$ によって $\gamma = t^{-1}h$ と書けるとする。LMS2 を用いて、
$a \circ s = t' \circ h$ を満たす射 $a : Y' \to Z''$ と $S$ に属する
$t' : Z' \to Z''$ を選ぶ。図式で表すと
$$
\xymatrix{
& & Z \ar[d]^t \\
& Y \ar[r]^h \ar[d]^s & Z' \ar[d]^{t'} \\
X \ar[r]^{f, g} & Y' \ar[r]^a & Z''
}
$$
となる。このとき
$\gamma \circ \alpha = (t' \circ t)^{-1}(a \circ f)$ および
$\gamma \circ \beta = (t' \circ t)^{-1}(a \circ g)$ である。
したがって $\gamma \circ (\alpha + \beta)$ は
$(t' \circ t)^{-1}(a \circ (f + g)) =
(t' \circ t)^{-1}(a \circ f + a \circ g)$ によって表され、これは
$\gamma \circ \alpha + \gamma \circ \beta$ を表す。

\medskip\noindent
最後に、$\delta : W \to X$ を $S^{-1}\mathcal{C}$ の別の射とする。
ある $i : W \to X'$ と $S$ に属する $r : X \to X'$ によって
$\delta = r^{-1}i$ と書けるとする。次の図式が可換となるような
$S$ に属する射 $s' : Y' \to Y''$ と射 $a'', b'' : X' \to Y''$ が存在すると
主張する。
$$
\xymatrix{
& & & Y \ar[d]^s \\
& X \ar[rr]^{f, g, f + g} \ar[d]^r & & Y' \ar[d]^{s'} \\
W \ar[r]^i & X' \ar[rr]^{a'', b'', a'' + b''} & & Y''
}
$$
実際、まず LMS2 を用いて $S$ に属する
$s_1 : Y' \to Y_1$, $s_2 : Y' \to Y_2$ と射
$a : X' \to Y_1$, $b : X' \to Y_2$ であって
% P02-E0007: frozen-source candidate retained; see control/ERRATA_CANDIDATES.jsonl.
$a \circ r = s_1 \circ f$ および $b \circ r = s_2 \circ f$ を満たすものを選ぶ。
次に、圏 $Y'/S$ がフィルター圏であること
（Categories, Remark \ref{categories-remark-left-localization-morphisms-colimit}）を
用いて、$s' = a' \circ s_1$ および $s' = b' \circ s_2$ を満たす
$s' : Y' \to Y''$ と射 $a' : Y_1 \to Y''$, $b' : Y_2 \to Y''$ を取る。
$a'' = a' \circ a$ および $b'' = b' \circ b$ と置けばよい。このとき合成
$\alpha \circ \delta$ および $\beta \circ \delta$ はそれぞれ
$(s' \circ s)^{-1}(a'' \circ i)$ および $(s' \circ s)^{-1}(b'' \circ i)$ によって
表される。したがって $\alpha \circ \delta + \beta \circ \delta$ は
$(s' \circ s)^{-1}(a'' \circ i + b'' \circ i) =
(s' \circ s)^{-1}((a'' + b'') \circ i)$ によって表され、これは再び図式により
$(\alpha + \beta) \circ \delta$ の代表元である。
\end{proof}

\begin{lemma}
\label{homology-lemma-localization-additive}
$\mathcal{C}$ を加法圏とし、$S$ を左または右乗法系とする。このとき
$S^{-1}\mathcal{C}$ は加法圏であり、局所化関手
$Q : \mathcal{C} \to S^{-1}\mathcal{C}$ は加法的である。
\end{lemma}

\begin{proof}
補題 \ref{homology-lemma-localization-preadditive} により、$S^{-1}\mathcal{C}$ は前加法圏で
$Q$ は加法的である。補題 \ref{homology-lemma-additive-additive} により
$S^{-1}\mathcal{C}$ は零対象と直和をもつ。
\end{proof}

\noindent
次の補題は、乗法系を逆にした場合の局所化関手の「核」を記述する。

\begin{lemma}
\label{homology-lemma-kernel-localization}
$\mathcal{C}$ を加法圏、$S$ を乗法系とし、$X$ を $\mathcal{C}$ の対象とする。
次の条件は同値である。
\begin{enumerate}
\item $S^{-1}\mathcal{C}$ において $Q(X) = 0$ である。
\item $0 : X \to Y$ が $S$ の元となるような
$Y \in \Ob(\mathcal{C})$ が存在する。
\item $0 : Z \to X$ が $S$ の元となるような
$Z \in \Ob(\mathcal{C})$ が存在する。
\end{enumerate}
\end{lemma}

\begin{proof}
(2) が成り立つなら $0 = Q(0) : Q(X) \to Q(Y)$ は同型である。
加法圏 $S^{-1}\mathcal{C}$ では、これは $Q(X) = 0$ を含意する。
したがって (2) $\Rightarrow$ (1) である。同様に (3) $\Rightarrow$ (1) である。
$Q(X) = 0$ と仮定する。これは射 $f : 0 \to X$ が
$S^{-1}\mathcal{C}$ において同型に写されることを意味する。したがって
Categories, Lemma \ref{categories-lemma-what-gets-inverted} により、
$fg \in S$ を満たす射 $g : Z \to 0$ が存在する。これで
(1) $\Rightarrow$ (3) が示された。同様に (1) $\Rightarrow$ (2) である。
\end{proof}

\begin{lemma}
\label{homology-lemma-localization-abelian}
$\mathcal{A}$ をアーベル圏とする。
\begin{enumerate}
\item $S$ が左乗法系なら、圏 $S^{-1}\mathcal{A}$ は余核をもち、関手
$Q : \mathcal{A} \to S^{-1}\mathcal{A}$ は余核と可換する。
\item $S$ が右乗法系なら、圏 $S^{-1}\mathcal{A}$ は核をもち、関手
$Q : \mathcal{A} \to S^{-1}\mathcal{A}$ は核と可換する。
\item $S$ が乗法系なら、圏 $S^{-1}\mathcal{A}$ はアーベル圏であり、関手
$Q : \mathcal{A} \to S^{-1}\mathcal{A}$ は完全である。
\end{enumerate}
\end{lemma}

\begin{proof}
$S$ を左乗法系とする。$a : X \to Y$ を $S^{-1}\mathcal{A}$ の射とする。
このとき、$S$ に属するある $s : Y \to Y'$ と $f : X \to Y'$ によって
$a = s^{-1}f$ と書ける。$Q(s)$ は同型なので、
$\Coker(a : X \to Y)$ の存在は
$\Coker(Q(f) : X \to Y')$ の存在と同値である。$\Coker(Q(f))$ は $0$ と
$Q(f)$ の余等化子であるから、Categories, Lemma
\ref{categories-lemma-left-localization-limits} により
$\Coker(Q(f))$ は $Q(\Coker(f))$ によって表される。これで (1) が示された。

\medskip\noindent
(2) は (1) の双対である。

\medskip\noindent
$S$ が乗法系なら、$S$ は左乗法系であると同時に右乗法系でもある。したがって
$S^{-1}\mathcal{A}$ は核と余核をもち、$Q$ は核および余核と可換する。(3) の証明を
完了するには、$S^{-1}\mathcal{A}$ において $\Coim = \Im$ であることを示す
必要がある。$S^{-1}\mathcal{A}$ の任意の矢印がある矢印 $Q(f)$ と同型であることを
再び用いれば、この結果は $\mathcal{A}$ における結果から従う。
\end{proof}





\section{ジョルダン--ヘルダー}
\label{homology-section-jordan-holder}

\noindent
ジョルダン--ヘルダーの補題は補題 \ref{homology-lemma-jordan-holder} である。
まずいくつかの定義を述べる。

\begin{definition}
\label{homology-definition-simple}
$\mathcal{A}$ をアーベル圏とする。$\mathcal{A}$ の対象 $A$ が零でなく、
$A$ の部分対象が $0$ と $A$ だけであるとき、その対象は {\it 単純} であるという。
\end{definition}

\begin{definition}
\label{homology-definition-Artinian}
$\mathcal{A}$ をアーベル圏とする。
\begin{enumerate}
\item $\mathcal{A}$ の対象 $A$ が部分対象について降鎖条件を満たすとき、またその
ときに限り、その対象は {\it アルティン的} であるという。
\item $\mathcal{A}$ のすべての対象がアルティン的であるとき、$\mathcal{A}$ は
{\it アルティン的} であるという。
\end{enumerate}
\end{definition}

\begin{definition}
\label{homology-definition-Noetherian}
$\mathcal{A}$ をアーベル圏とする。
\begin{enumerate}
\item $\mathcal{A}$ の対象 $A$ が部分対象について昇鎖条件を満たすとき、またその
ときに限り、その対象は {\it ネーター的} であるという。
\item $\mathcal{A}$ のすべての対象がネーター的であるとき、$\mathcal{A}$ は
{\it ネーター的} であるという。
\end{enumerate}
\end{definition}

\begin{lemma}
\label{homology-lemma-ses-artinian}
$\mathcal{A}$ をアーベル圏とし、
$0 \to A_1 \to A_2 \to A_3 \to 0$ を $\mathcal{A}$ の短完全列とする。
$A_2$ がアルティン的であるための必要十分条件は、$A_1$ と $A_3$ が
アルティン的であることである。
\end{lemma}

\begin{proof}
省略する。
\end{proof}

\begin{lemma}
\label{homology-lemma-ses-noetherian}
$\mathcal{A}$ をアーベル圏とし、
$0 \to A_1 \to A_2 \to A_3 \to 0$ を $\mathcal{A}$ の短完全列とする。
$A_2$ がネーター的であるための必要十分条件は、$A_1$ と $A_3$ が
ネーター的であることである。
\end{lemma}

\begin{proof}
省略する。
\end{proof}

\begin{lemma}
\label{homology-lemma-finite-length}
$\mathcal{A}$ をアーベル圏とし、$A$ を $\mathcal{A}$ の対象とする。
次の条件は同値である。
\begin{enumerate}
\item $A$ はアルティン的かつネーター的である。
\item 部分対象によるフィルトレーション
$0 \subset A_1 \subset A_2 \subset \ldots \subset A_n = A$ であって、
$i = 1, \ldots, n$ に対して $A_i/A_{i - 1}$ が単純であるものが存在する。
\end{enumerate}
\end{lemma}

\begin{proof}
(1) を仮定する。$A$ が零なら (2) が成り立つ。$A$ が零でなければ、
アルティン性により最小の非零対象 $A_1 \subset A$ が存在する。もちろん $A_1$ は
単純である。$A_1 = A$ なら証明は終わる。そうでなければ、
$A_2 \not = A_1$ を満たすもののうち極小の
$A_1 \subset A_2 \subset A$ を取ることができる。このとき $A_2/A_1$ は
単純である。この手続きを続けると、$A_i/A_{i - 1}$ が単純となる $A$ の部分対象の列
$0 \subset A_1 \subset A_2 \subset \ldots $ が得られる。$A$ はネーター的なので
この過程は停止する。したがって (2) が従う。

\medskip\noindent
(2) を仮定する。$n$ に関する帰納法で (1) を示す。$n = 1$ なら $A$ は単純で、
明らかにネーター的かつアルティン的である。$n - 1$ に対して結果が成り立つなら、
短完全列 $0 \to A_{n - 1} \to A_n \to A_n/A_{n - 1} \to 0$ と
補題 \ref{homology-lemma-ses-artinian} および補題 \ref{homology-lemma-ses-noetherian} を用いて
$n$ に対する結論を得る。
\end{proof}

\begin{lemma}[Jordan-H\"older]
\label{homology-lemma-jordan-holder}
$\mathcal{A}$ をアーベル圏とする。$A$ を、補題 \ref{homology-lemma-finite-length} の
同値な条件を満たす $\mathcal{A}$ の対象とする。二つのフィルトレーション
$$
0 \subset A_1 \subset A_2 \subset \ldots \subset A_n = A
\quad\text{and}\quad
0 \subset B_1 \subset B_2 \subset \ldots \subset B_m = A
$$
が与えられ、$S_i = A_i/A_{i - 1}$ と $T_j = B_j/B_{j - 1}$ が単純対象であると
する。このとき $n = m$ であり、任意の $i \in \{1, \ldots, n\}$ に対して
$S_i \cong T_{\sigma(i)}$ を満たす $\{1, \ldots, n\}$ の置換 $\sigma$ が存在する。
\end{lemma}

\begin{proof}
$A_1 \subset B_j$ を満たす最小の添字を $j$ とする。このとき写像
$S_1 = A_1 \to B_j/B_{j - 1} = T_j$ は同型である。さらに対象
$A/A_1 = A_n/A_1 = B_m/A_1$ は、二つのフィルトレーション
$$
0 \subset A_2/A_1 \subset A_3/A_1 \subset \ldots \subset A_n/A_1
$$
および
$$
0 \subset (B_1 + A_1)/A_1 \subset \ldots \subset
(B_{j - 1} + A_1)/A_1 = B_j/A_1 \subset B_{j + 1}/A_1
\subset \ldots \subset B_m/A_1
$$
をもつ。帰納法により結論を得る。
\end{proof}






\section{セール部分圏}
\label{homology-section-serre-subcategories}

\noindent
\cite[Chapter I, Section 1]{Serre_homotopie_classes} では、アーベル群の
「類」という概念が定義されている。この概念は多くの著者によって
（少しずつ異なる仕方で）アーベル圏へ拡張された。ここではセールの元の定義と
実質的に同一な次の変種を用いる。

\begin{definition}
\label{homology-definition-serre-subcategory}
\begin{reference}
\cite[Condition (I) on page 259]{Serre_homotopie_classes}
\end{reference}
$\mathcal{A}$ をアーベル圏とする。
\begin{enumerate}
\item $\mathcal{A}$ の {\it セール部分圏} とは、$\mathcal{A}$ の空でない
充満部分圏 $\mathcal{C}$ であって、完全列\footnote{定義
\ref{homology-definition-exact} により、これは $\Im(A \to B) = \Ker(B \to C)$ を意味する。}
$$
A \to B \to C
$$
において $A, C \in \Ob(\mathcal{C})$ なら
$B \in \Ob(\mathcal{C})$ でもあるものをいう。
\item $\mathcal{A}$ の {\it 弱セール部分圏} とは、$\mathcal{A}$ の空でない
充満部分圏 $\mathcal{C}$ であって、完全列
$$
A_0 \to A_1 \to A_2 \to A_3 \to A_4
$$
において $A_0, A_1, A_3, A_4$ が $\mathcal{C}$ に属すれば $A_2$ も
$\mathcal{C}$ に属するものをいう。
\end{enumerate}
\end{definition}

\noindent
文献によっては第二の概念を「厚い」部分圏と呼び、別の文献では第一の概念を
「厚い」部分圏と呼ぶ。しかし、セール部分圏という概念については、上の定義が
普遍的に受け入れられているようである。いずれの場合も圏 $\mathcal{C}$ は
アーベル圏であり、包含関手 $\mathcal{C} \to \mathcal{A}$ は充満忠実な完全関手で
あることに注意する。これらの型の部分圏をさらに詳しく特徴づけよう。

\begin{lemma}
\label{homology-lemma-characterize-serre-subcategory}
$\mathcal{A}$ をアーベル圏とし、$\mathcal{C}$ を $\mathcal{A}$ の部分圏とする。
$\mathcal{C}$ がセール部分圏であるための必要十分条件は、次の条件が成り立つことで
ある。
\begin{enumerate}
\item $0 \in \Ob(\mathcal{C})$ である。
\item $\mathcal{C}$ は $\mathcal{A}$ の厳密充満部分圏である。
\item $\mathcal{C}$ の対象の任意の部分対象および商は $\mathcal{C}$ の対象である。
\item $A \in \Ob(\mathcal{A})$ が $\mathcal{C}$ の対象の拡大なら、
$A \in \Ob(\mathcal{C})$ でもある。
\end{enumerate}
さらに、セール部分圏はアーベル圏であり、包含関手は完全である。
\end{lemma}

\begin{proof}
省略する。
\end{proof}

\begin{lemma}
\label{homology-lemma-characterize-weak-serre-subcategory}
$\mathcal{A}$ をアーベル圏とし、$\mathcal{C}$ を $\mathcal{A}$ の部分圏とする。
$\mathcal{C}$ が弱セール部分圏であるための必要十分条件は、次の条件が成り立つこと
である。
\begin{enumerate}
\item $0 \in \Ob(\mathcal{C})$ である。
\item $\mathcal{C}$ は $\mathcal{A}$ の厳密充満部分圏である。
\item $\mathcal{C}$ の対象の間の射について $\mathcal{A}$ で取った核と余核は
$\mathcal{C}$ に属する。
\item $A \in \Ob(\mathcal{A})$ が $\mathcal{C}$ の対象の拡大なら、
$A \in \Ob(\mathcal{C})$ でもある。
\end{enumerate}
さらに、弱セール部分圏はアーベル圏であり、包含関手は完全である。
\end{lemma}

\begin{proof}
省略する。
\end{proof}

\begin{lemma}
\label{homology-lemma-kernel-exact-functor}
$\mathcal{A}$, $\mathcal{B}$ をアーベル圏とし、
$F : \mathcal{A} \to \mathcal{B}$ を完全関手とする。このとき
$F(C) = 0$ を満たす $\mathcal{A}$ の対象 $C$ からなる充満部分圏は
$\mathcal{A}$ のセール部分圏をなす。
\end{lemma}

\begin{proof}
省略する。
\end{proof}

\begin{definition}
\label{homology-definition-kernel-category}
$\mathcal{A}$, $\mathcal{B}$ をアーベル圏とし、
$F : \mathcal{A} \to \mathcal{B}$ を完全関手とする。
$F(C) = 0$ を満たす $\mathcal{A}$ の対象 $C$ からなる充満部分圏を
{\it 関手 $F$ の核} といい、$\Ker(F)$ と書くことがある。
\end{definition}

\noindent
アーベル圏の任意のセール部分圏は、ある完全関手の核である。
Examples, Section \ref{examples-section-serre-quotient-modulo-torsion-modules} では、
セールの元来の例であるねじれ群についてこれを論じる。

\begin{lemma}
\label{homology-lemma-serre-subcategory-is-kernel}
$\mathcal{A}$ をアーベル圏とし、
$\mathcal{C} \subset \mathcal{A}$ をセール部分圏とする。アーベル圏
$\mathcal{A}/\mathcal{C}$ と完全関手
$$
F : \mathcal{A} \longrightarrow \mathcal{A}/\mathcal{C}
$$
であって、本質的全射であり、その核が $\mathcal{C}$ であるものが存在する。
圏 $\mathcal{A}/\mathcal{C}$ と関手 $F$ は次の普遍性によって特徴づけられる。
$\mathcal{C} \subset \Ker(G)$ を満たす任意の完全関手
$G : \mathcal{A} \to \mathcal{B}$ に対して、一意な完全関手
$H : \mathcal{A}/\mathcal{C} \to \mathcal{B}$ による分解
$G = H \circ F$ が存在する。
\end{lemma}

\begin{proof}
$\mathcal{A}$ の矢印の集合を次の式で定める。
$$
S = \{f \in \text{Arrows}(\mathcal{A}) \mid
\Ker(f), \Coker(f) \in \Ob(\mathcal{C}) \}.
$$
$S$ は乗法系であると主張する。このためには MS1, MS2, MS3 を確認すればよい。
Categories, Definition \ref{categories-definition-multiplicative-system} を参照されたい。

\medskip\noindent
恒等射が $S$ の元であることは明らかである。
$f : A \to B$ と $g : B \to C$ を $S$ の元とする。完全列
$$
\begin{matrix}
0 \to \Ker(f) \to \Ker(gf) \to \Ker(g) \\
\Coker(f) \to \Coker(gf) \to \Coker(g) \to 0
\end{matrix}
$$
がある。したがって $gf \in S$ である。これで MS1 が示された。
（実際、同様の議論により $S$ が飽和乗法系であることも示せる。Categories,
Definition \ref{categories-definition-saturated-multiplicative-system} を参照。）

\medskip\noindent
実線部分が与えられた図式
$$
\xymatrix{
A \ar[d]_t \ar[r]_g & B \ar@{..>}[d]^s \\
C \ar@{..>}[r]^f & C \amalg_A B
}
$$
を考え、$t \in S$ とする。
$W = C \amalg_A B =  \Coker((t, -g) : A \to C \oplus B)$ と置く。
このとき $\Ker(t) \to \Ker(s)$ は全射であり、
$\Coker(t) \to \Coker(s)$ は同型である。したがって $s$ は $S$ の元である。
これで LMS2 が示され、RMS2 の証明は双対的である。

\medskip\noindent
最後に射 $f, g : B \to C$ と、$S$ に属し
$f \circ s = g \circ s$ を満たす射 $s : A \to B$ を考える。これは
$(f - g) \circ s = 0$ を意味する。さらに、これは
$I = \Im(f - g) \subset C$ が $\Coker(s)$ の商であり、したがって
$\mathcal{C}$ の対象であることを意味する。よって
$t : C \to C' = C/I$ は $S$ の元であり、$t \circ (f - g) = 0$、すなわち
$t \circ f = t \circ g$ を満たす。これで LMS3 が示され、RMS3 の証明は
双対的である。

\medskip\noindent
$S$ が乗法系であることを示したので、
$\mathcal{A}/\mathcal{C} = S^{-1}\mathcal{A}$ と置き、$F$ を局所化関手 $Q$ とする。
補題 \ref{homology-lemma-localization-abelian} により、圏 $\mathcal{A}/\mathcal{C}$ は
アーベル圏で $F$ は完全である。$X$ が $F = Q$ の核に属するなら、
補題 \ref{homology-lemma-kernel-localization} により $0 : X \to Z$ は $S$ の元であり、
したがって $X$ は $\mathcal{C}$ の対象である。すなわち $F$ の核は
$\mathcal{C}$ である。最後に $G$ が補題の主張にある関手なら、$G$ は $S$ の
すべての元を同型に写す。したがって局所化の普遍性、すなわち Categories,
Lemma \ref{categories-lemma-properties-left-localization} により関手
$H : \mathcal{A}/\mathcal{C} \to \mathcal{B}$ を得る。なお $H$ が完全であることを
示さなければならない。そのためには、$H$ が核と余核を取る操作と可換することを
示せば十分である。補題 \ref{homology-lemma-exact-functor} を参照されたい。
$A \to B$ を $\mathcal{A}/\mathcal{C}$ の射とする。この $A \to B$ は、$S$ に属する
$s : A' \to A$ と $\mathcal{A}$ の任意の射 $f : A' \to B$ によって
$fs^{-1}$ と表すことができる。$F = Q$ は $s$ を商圏
$\mathcal{A}/\mathcal{C}$ の同型に写すので、$\mathcal{A}$ の射
$f : A \to B$ の核と余核を取る操作と $H$ が可換することを示せば十分である。
しかしここでは $H(f) = G(f)$ であり、$G$ が完全であることから結果が従う。
\end{proof}

\begin{lemma}
\label{homology-lemma-quotient-by-kernel-exact-functor}
$\mathcal{A}$, $\mathcal{B}$ をアーベル圏とし、
$F : \mathcal{A} \to \mathcal{B}$ を完全関手とする。
$\mathcal{C} \subset \mathcal{A}$ を、$F$ の核に含まれるセール部分圏とする。
$\mathcal{C} = \Ker(F)$ であるための必要十分条件は、誘導される関手
$\overline{F} : \mathcal{A}/\mathcal{C} \to \mathcal{B}$
（補題 \ref{homology-lemma-serre-subcategory-is-kernel}）が忠実であることである。
\end{lemma}

\begin{proof}
補題 \ref{homology-lemma-serre-subcategory-is-kernel} の結果を以下では断りなく用いる。
「ならば」の向きは、構成により $\overline{F}$ の核が零であることから従う。
実際、$f : X \to Y$ が $\mathcal{A}/\mathcal{C}$ の射で
$\overline{F}(f) = 0$ を満たすなら、
$\overline{F}(\Im(f)) = \Im(\overline{F}(f)) = 0$ であるから、$F$ の核に関する
仮定により $\Im(f) = 0$ である。したがって $f = 0$ である。

\medskip\noindent
「逆ならば」の向きについて、$X$ を $F(X)
= 0$ を満たす $\mathcal{A}$ の対象とする。このとき
$\overline{F}(\text{id}_X) = \text{id}_{\overline{F}(X)} = 0$ なので、
$\overline{F}$ の忠実性により $\mathcal{A}/\mathcal{C}$ において
$\text{id}_X = 0$ である。したがって $\mathcal{A}/\mathcal{C}$ において
$X = 0$、すなわち $X \in
\Ob(\mathcal{C})$ である。
\end{proof}





\section{K 群}
\label{homology-section-K-groups}

\noindent
アーベル圏の $K_0$ について少しだけ述べる。

\begin{definition}
\label{homology-definition-K-zero}
$\mathcal{A}$ をアーベル圏とする。$K_0(\mathcal{A})$ を
{\it $\mathcal{A}$ の第零 $K$ 群} という。これは次のように構成される
アーベル群である。$\mathcal{A}$ の対象上の自由アーベル群を取り、各短完全列
$0 \to A \to B \to C \to 0$ に対して関係式
$[B] - [A] - [C] = 0$ を課す。
\end{definition}

\noindent
別の言い方をすれば、$K_0(\mathcal{A})$ は表示
$$
\bigoplus_{A \to B \to C\text{ ses}}
\mathbf{Z}[A \to B \to C]
\longrightarrow
\bigoplus_{A \in \Ob(\mathcal{A})}
\mathbf{Z}[A]
\longrightarrow
K_0(\mathcal{A})
\longrightarrow
0
$$
をもち、$[A \to B \to C] \mapsto [B] - [A] - [C]$ である。
短完全列 $0 \to 0 \to 0 \to 0 \to 0$ から
$K_0(\mathcal{A})$ における関係式 $[0] = 0$ が得られる。特に断らない限り
すべての圏を「小さい」としているので、集合論上の問題はない。
環上の加群の圏についての $K$ 群の例は
Algebra, Section \ref{algebra-section-K-groups} で計算した。

\begin{lemma}
\label{homology-lemma-exact-functor-K-groups}
$F : \mathcal{A} \to \mathcal{B}$ をアーベル圏の間の完全関手とする。
このとき $F$ は、$K_0(F)([A]) = [F(A)]$ と置くことによって $K$ 群の準同型
$K_0(F) : K_0(\mathcal{A}) \to K_0(\mathcal{B})$ を誘導する。
\end{lemma}

\begin{proof}
自明である。
\end{proof}

\noindent
アーベル圏 $\mathcal{A}$ の対象 $M$ と、次の形の複体が与えられたとする。
\begin{equation}
\label{homology-equation-cyclic-complex}
\xymatrix{
\ldots \ar[r] &
M \ar[r]^\varphi &
M \ar[r]^\psi &
M \ar[r]^\varphi &
M \ar[r] & \ldots
}
\end{equation}
この状況で
$$
H^0(M, \varphi, \psi) = \Ker(\psi)/\Im(\varphi)
, \quad\text{and}\quad
H^1(M, \varphi, \psi) = \Ker(\varphi)/\Im(\psi).
$$
と定める。

\begin{lemma}
\label{homology-lemma-serre-subcategory-K-groups}
$\mathcal{A}$ をアーベル圏とし、$\mathcal{C} \subset \mathcal{A}$ を
セール部分圏として $\mathcal{B} = \mathcal{A}/\mathcal{C}$ と置く。
\begin{enumerate}
\item 完全関手 $\mathcal{C} \to \mathcal{A}$ および
$\mathcal{A} \to \mathcal{B}$ は $K$ 群の完全列
$$
K_0(\mathcal{C}) \to
K_0(\mathcal{A}) \to
K_0(\mathcal{B}) \to
0
$$
を誘導する。
\item $K_0(\mathcal{C}) \to K_0(\mathcal{A})$ の核は、
$$
[H^0(M, \varphi, \psi)] - [H^1(M, \varphi, \psi)]
$$
の形の元全体に等しい。ここで $(M, \varphi, \psi)$ は
(\ref{homology-equation-cyclic-complex}) のような複体で、$\mathcal{B}$ では完全になるという
性質、言い換えれば $H^0(M, \varphi, \psi)$ と $H^1(M, \varphi, \psi)$ が
$\mathcal{C}$ の対象であるという性質をもつ。
\end{enumerate}
\end{lemma}

\begin{proof}
(1) を証明する。$K_0(\mathcal{A}) \to K_0(\mathcal{B})$ が全射であり、合成
$K_0(\mathcal{C}) \to K_0(\mathcal{A}) \to K_0(\mathcal{B})$ が零であることは
明らかである。$x \in K_0(\mathcal{A})$ を $K_0(\mathcal{B})$ で零に写る元とする。
$\mathcal{A}$ の $A, A'$ によって $x = [A] - [A']$ と書ける（楽しい演習）。
対応する $\mathcal{B}$ の対象を $B, B'$ と書く。$x$ が
$K_0(\mathcal{B})$ で零に写ることは、有限集合 $I = I^+ \amalg I^{-}$ と、
各 $i \in I$ に対する $\mathcal{B}$ の短完全列
$$
0 \to B_i \to B'_i \to B''_i \to 0
$$
であって、$\mathcal{B}$ の対象の同型類上の自由アーベル群において
$$
[B] - [B'] = \sum\nolimits_{i \in I^{+}} ([B'_i] - [B_i] - [B''_i])
-
\sum\nolimits_{i \in I^{-}} ([B'_i] - [B_i] - [B''_i])
$$
を満たすものが存在することを意味する。これは
$$
[B]
+ \sum\nolimits_{i \in I^{+}} ([B_i] + [B''_i])
+ \sum\nolimits_{i \in I^{-}} [B'_i]
=
[B']
+ \sum\nolimits_{i \in I^{-}} ([B_i] + [B''_i])
+ \sum\nolimits_{i \in I^{+}} [B'_i].
$$
と書き直せる。左辺と右辺は $\mathcal{B}$ の同じ対象の同型類を重複度込みで
含まなければならないので、全単射
$$
\tau :
\{B\} \amalg \{B_i, B''_i; i \in I^+\} \amalg \{B'_i; i \in I^-\}
\longrightarrow
\{B'\} \amalg \{B_i, B''_i; i \in I^-\} \amalg \{B'_i; i \in I^+\}
$$
であって、$N$ と $\tau(N)$ が $\mathcal{B}$ で同型となるものが存在する。
補題 \ref{homology-lemma-serre-subcategory-is-kernel} と
補題 \ref{homology-lemma-localization-abelian} の証明により、各 $i \in I$ に対して
$\mathcal{A}$ の短完全列
$$
0 \to A_i \to A'_i \to A''_i \to 0
$$
を、$B_i, B'_i, B''_i$ が $\mathcal{B}$ における
$A_i, A'_i, A''_i$ の像と同型となるように選ぶことができる。このことから、対応する
全単射
$$
\tau :
\{A\} \amalg \{A_i, A''_i; i \in I^+\} \amalg \{A'_i; i \in I^-\}
\longrightarrow
\{A'\} \amalg \{A_i, A''_i; i \in I^-\} \amalg \{A'_i; i \in I^+\}
$$
は次の性質を満たす。$M$ と $\tau(M)$ は $\mathcal{A}$ の対象で、
$\mathcal{B}$ において同型となる。これは
$[M] - [\tau(M)]$ が $K_0(\mathcal{C}) \to K_0(\mathcal{A})$ の像に属することを
意味する。実際、$\mathcal{B}$ における同型は、$\mathcal{A}$ における図式
$M \leftarrow M' \rightarrow \tau(M)$ で与えられ、$M' \to M$ と
$M' \to \tau(M)$ はともに核と余核が $\mathcal{C}$ に属する。
逆にたどると $x = [A] - [A']$ は
$K_0(\mathcal{C}) \to K_0(\mathcal{A})$ の像に属し、(1) の証明が完了する。

\medskip\noindent
(2) を証明する。証明は (1) と同様だが、帳簿付けが少し増える。まず、
$[H^0(M, \varphi, \psi)] - [H^1(M, \varphi, \psi)]$ の形の任意の類が、計算
\begin{align*}
0 & = [M] - [M] \\
& =  [\Ker(\varphi)] + [\Im(\varphi)]
- [\Ker(\psi)] - [\Im(\psi)] \\
& =
[\Ker(\varphi)/\Im(\psi)] -
[\Ker(\psi)/\Im(\varphi)] \\
& = [H^1(M, \varphi, \psi)] - [H^0(M, \varphi, \psi)]
\end{align*}
により $K_0(\mathcal{A})$ で零になることに注意する。したがって
$K_0(\mathcal{C}) \to K_0(\mathcal{A})$ の核の任意の元がこの形であることを
示せば十分である。

\medskip\noindent
$K_0(\mathcal{C})$ の任意の元 $x$ は、$\mathcal{C}$ の二対象の差
$x = [P] - [Q]$ として表せる（楽しい演習）。この元が
$K_0(\mathcal{A})$ で零に写ると仮定する。これは次のデータが存在することを
意味する。
\begin{enumerate}
\item 有限集合 $I = I^{+} \amalg I^{-}$。
\item 各 $i \in I$ に対する $\mathcal{A}$ の短完全列
$0 \to A_i \to B_i \to C_i \to 0$。
\end{enumerate}
これらは、$\mathcal{A}$ の対象上の自由アーベル群において
$$
[P] - [Q] =
\sum\nolimits_{i \in I^{+}} ([B_i] - [A_i] - [C_i])
-
\sum\nolimits_{i \in I^{-}} ([B_i] - [A_i] - [C_i])
$$
を満たす。これは
$$
[P]
+ \sum\nolimits_{i \in I^{+}} ([A_i] + [C_i])
+ \sum\nolimits_{i \in I^{-}} [B_i]
=
[Q]
+ \sum\nolimits_{i \in I^{-}} ([A_i] + [C_i])
+ \sum\nolimits_{i \in I^{+}} [B_i].
$$
と書き直せる。左辺と右辺は $\mathcal{A}$ の同じ対象を重複度込みで含まなければ
ならないので、上に現れる項の間に全単射 $\tau$ が存在する。次のように置く。
$$
T^{+} =
\{p\}\ \amalg\ \{a, c\} \times I^{+}\ \amalg\ \{b\} \times I^{-}
$$
および
$$
T^{-} =
\{q\}\ \amalg\ \{a, c\} \times I^{-}\ \amalg\ \{b\} \times I^{+}.
$$
$T = T^{+} \amalg T^{-} = \{p, q\} \amalg \{a, b, c\} \times I$ と置く。
$t \in T$ に対して
$$
O(t)
=
\left\{
\begin{matrix}
P & \text{if} & t = p \\
Q & \text{if} & t = q \\
A_i & \text{if} & t = (a, i) \\
B_i & \text{if} & t = (b, i) \\
C_i & \text{if} & t = (c, i)
\end{matrix}
\right.
$$
と定める。すると $\tau : T^{+} \to T^{-}$ を、すべての $t \in T^{+}$ に対して
$O(t) = O(\tau(t))$ を満たす全単射とみなせる。
$t^{-}_0 = \tau(p)$ とし、$\tau(t^{+}_0) = q$ を満たす一意な元
$t^{+}_0 \in T^{+}$ を取る。対象
$$
M^{+} = \bigoplus\nolimits_{t \in T^{+}} O(t)
$$
を考える。$\tau$ を用いると、これは対象
$$
M^{-} = \bigoplus\nolimits_{t \in T^{-}} O(t)
$$
に等しい。写像
$$
\varphi : M^{+} \longrightarrow M^{-}
$$
を次のように定める。$t = (a, i)$, $i \in I^{+}$ に対応する直和因子
$O(t) = A_i$ 上では、写像 $A_i \to B_i$ を $M^{-}$ の直和因子
$O((b, i)) = B_i$ への写像として用いる。$(b, i)$, $i \in I^{-}$ に対応する
直和因子 $O(t) = B_i$ 上では、写像 $B_i \to C_i$ を $M^{-}$ の直和因子
$O((c, i)) = C_i$ への写像として用いる。$p$ および
$(c, i)$, $i \in I^{+}$ に対応する直和因子上では零写像とする。
同様に写像
$$
\psi : M^{-} \longrightarrow M^{+}
$$
を次のように定める。$t = (a, i)$, $i \in I^{-}$ に対応する直和因子
$O(t) = A_i$ 上では、写像 $A_i \to B_i$ を $M^{+}$ の直和因子
$O((b, i)) = B_i$ への写像として用いる。$(b, i)$, $i \in I^{+}$ に対応する
直和因子 $O(t) = B_i$ 上では、写像 $B_i \to C_i$ を $M^{+}$ の直和因子
$O((c, i)) = C_i$ への写像として用いる。$q$ および
$(c, i)$, $i \in I^{-}$ に対応する直和因子上では零写像とする。

\medskip\noindent
$\varphi$ の核は、直和因子 $P$、直和因子
$O((c, i)) = C_i$, $i \in I^{+}$、および直和因子
$O((b, i)) = B_i$, $i \in I^{-}$ の内部の部分対象 $A_i$ の直和に等しい。
$\psi$ の像は、直和因子 $O((c, i)) = C_i$, $i \in I^{+}$、および直和因子
$O((b, i)) = B_i$, $i \in I^{-}$ の内部の部分対象 $A_i$ の直和に等しい。
言い換えれば
$$
P \cong \Ker(\varphi)/\Im(\psi).
$$
である。まったく同様に
$$
Q \cong \Ker(\psi)/\Im(\varphi).
$$
である。上で述べたとおり、全単射 $\tau$ の存在から
$M^{+} = M^{-}$ が従うので、補題が従う。
\end{proof}





\section{コホモロジー的デルタ関手}
\label{homology-section-cohomological-delta-functor}

\begin{definition}
\label{homology-definition-cohomological-delta-functor}
$\mathcal{A}, \mathcal{B}$ をアーベル圏とする。$\mathcal{A}$ から
$\mathcal{B}$ への {\it コホモロジー的 $\delta$ 関手}、または単に
{\it $\delta$ 関手} とは、次のデータのことである。
\begin{enumerate}
\item 加法関手の族 $F^n : \mathcal{A} \to \mathcal{B}$, $n \geq 0$。
\item $\mathcal{A}$ の各短完全列 $0 \to A \to B \to C \to 0$ に対する
$\mathcal{B}$ の射の族
$\delta_{A \to B \to C} : F^n(C) \to F^{n + 1}(A)$, $n \geq 0$。
\end{enumerate}
これらのデータは次の公理を満たすものとする。
\begin{enumerate}
\item 上の各短完全列に対して、列
$$
\xymatrix{
0 \ar[r] &
F^0(A) \ar[r] &
F^0(B) \ar[r] &
F^0(C) \ar[lld]^{\delta_{A \to B \to C}} \\
 &
F^1(A) \ar[r] &
F^1(B) \ar[r] &
F^1(C) \ar[lld]^{\delta_{A \to B \to C}} \\
 &
F^2(A) \ar[r] &
F^2(B) \ar[r] &
\ldots
}
$$
は完全である。
\item $\mathcal{A}$ の短完全列の各射
$(A \to B \to C) \to (A' \to B' \to C')$ に対して、図式
$$
\xymatrix{
F^n(C) \ar[d] \ar[rr]_{\delta_{A \to B \to C}} & & F^{n + 1}(A) \ar[d] \\
F^n(C') \ar[rr]^{\delta_{A' \to B' \to C'}} & & F^{n + 1}(A')
}
$$
は可換である。
\end{enumerate}
\end{definition}

\noindent
特に、このことから $F^0$ は左完全である。

\begin{definition}
\label{homology-definition-morphism-delta-functors}
$\mathcal{A}, \mathcal{B}$ をアーベル圏とする。
$(F^n, \delta_F)$ と $(G^n, \delta_G)$ を $\mathcal{A}$ から
$\mathcal{B}$ への $\delta$ 関手とする。{\it $F$ から $G$ への $\delta$ 関手の射}
とは、関手の変換の族 $t^n : F^n \to G^n$, $n \geq 0$ であって、
$\mathcal{A}$ の各短完全列 $0 \to A \to B \to C \to 0$ に対して図式
$$
\xymatrix{
F^n(C) \ar[d]_{t^n} \ar[rr]_{\delta_{F, A \to B \to C}} &
& F^{n + 1}(A) \ar[d]^{t^{n + 1}} \\
G^n(C) \ar[rr]^{\delta_{G, A \to B \to C}} & & G^{n + 1}(A)
}
$$
が可換となるものをいう。
\end{definition}

\begin{definition}
\label{homology-definition-universal-delta-functor}
$\mathcal{A}, \mathcal{B}$ をアーベル圏とし、
$F = (F^n, \delta_F)$ を $\mathcal{A}$ から $\mathcal{B}$ への
$\delta$ 関手とする。任意の $\delta$ 関手 $G = (G^n, \delta_G)$ と任意の関手の射
$t : F^0 \to G^0$ に対して、$t = t^0$ を満たす一意な $\delta$ 関手の射
$\{t^n\}_{n \geq 0} : F \to G$ が存在するとき、またそのときに限り、$F$ を
{\it 普遍 $\delta$ 関手} という。
\end{definition}

\begin{lemma}
\label{homology-lemma-efface-implies-universal}
$\mathcal{A}, \mathcal{B}$ をアーベル圏とし、
$F = (F^n, \delta_F)$ を $\mathcal{A}$ から $\mathcal{B}$ への
$\delta$ 関手とする。任意の $n > 0$ と任意の $A \in \Ob(\mathcal{A})$ に対して、
$F^n(u) : F^n(A) \to F^n(B)$ が零となる単射 $u : A \to B$
（$A$ と $n$ に依存してよい）が存在すると仮定する。このとき $F$ は普遍
$\delta$ 関手である。
\end{lemma}

\begin{proof}
$G = (G^n, \delta_G)$ を $\mathcal{A}$ から $\mathcal{B}$ への $\delta$ 関手とし、
$t : F^0 \to G^0$ を関手の射とする。$t = t^0$ を満たす一意な
$\delta$ 関手の射 $\{t^n\}_{n \geq 0} : F \to G$ が存在することを示さなければ
ならない。$n$ に関する帰納法で $t^n$ を構成する。$n = 0$ では $t^0 = t$ と置く。
$i \leq n$ に対する関手の変換 $t^i$ の一意な列が既に構成され、次数
$\leq n$ の写像 $\delta$ と両立していると仮定する。

\medskip\noindent
$A \in \Ob(\mathcal{A})$ とする。仮定により $F^{n + 1}(u) = 0$ を満たす
埋め込み $u : A \to B$ を選べる。$C = B/u(A)$ と置く。
短完全列 $0 \to A \to B \to C \to 0$ と $\delta$ 関手 $F$ に対する
長完全コホモロジー列から、$u$ の選び方により
$F^{n + 1}(A) = \Coker(F^n(B) \to F^n(C))$ を得る。$t^n$ は既に定義されているので、
$$
t^{n + 1}_A : F^{n + 1}(A) \to G^{n + 1}(A)
$$
を、図式
$$
\xymatrix{
\Coker(F^n(B) \to F^n(C)) \ar[r]_{t^n}
\ar[d]_{\delta_{F, A \to B \to C}} &
\Coker(G^n(B) \to G^n(C))
\ar[d]^{\delta_{G, A \to B \to C}} \\
F^{n + 1}(A) \ar[r]^{t^{n + 1}_A} &
G^{n + 1}(A)
}
$$
を可換にする一意な写像として定める。これは課した要件によって明らかに一意に
定まる。これが関手の変換を定めることの確認は省略する。
\end{proof}

\begin{lemma}
\label{homology-lemma-uniqueness-universal-delta-functor}
$\mathcal{A}, \mathcal{B}$ をアーベル圏とし、
$F : \mathcal{A} \to \mathcal{B}$ を関手とする。
$F^0 = F$ を満たす $\mathcal{A}$ から $\mathcal{B}$ への普遍 $\delta$ 関手
$(F^n, \delta_F)$ が存在するなら、それは $\delta$ 関手の一意な同型を除いて
一意に定まる。
\end{lemma}

\begin{proof}
定義から直ちに従う。
\end{proof}






\section{複体}
\label{homology-section-complexes}

\noindent
もちろん鎖複体と余鎖複体の概念は双対であり、この節の二つの部分のうち一方だけを
読めばよい。好みの方を選ばれたい。（実際には、番号付けの規約が鎖複体と余鎖複体の
間を簡単に移行できるようにはなっていないため、これは完全にはうまくいかない。）

\medskip\noindent
加法圏 $\mathcal{A}$ における {\it 鎖複体 $A_\bullet$} とは、
$\mathcal{A}$ の複体
$$
\ldots \to
A_{n + 1} \xrightarrow{d_{n + 1}}
A_n \xrightarrow{d_n}
A_{n - 1} \to
\ldots
$$
のことである。言い換えれば、各 $i \in \mathbf{Z}$ に対して
$\mathcal{A}$ の対象 $A_i$ が与えられ、また各 $i \in \mathbf{Z}$ に対して
射 $d_i : A_i \to A_{i - 1}$ が与えられ、
すべての $i$ に対して $d_{i - 1} \circ d_i = 0$ が成り立つ。
{\it 鎖複体の射 $f : A_\bullet \to B_\bullet$} とは、すべての図式
$$
\xymatrix{
A_i \ar[r]_{d_i} \ar[d]_{f_i} & A_{i - 1} \ar[d]^{f_{i - 1}} \\
B_i \ar[r]^{d_i} & B_{i - 1}
}
$$
が可換となる射の族 $f_i : A_i \to B_i$ のことである。
{\it $\mathcal{A}$ の鎖複体の圏} を $\text{Ch}(\mathcal{A})$ と書く。
次の形の対象からなる充満部分圏
$$
\ldots \to A_2 \to A_1 \to A_0 \to 0 \to 0 \to \ldots
$$
を $\text{Ch}_{\geq 0}(\mathcal{A})$ と書く。言い換えれば、鎖複体
$A_\bullet$ が $\text{Ch}_{\geq 0}(\mathcal{A})$ に属するための必要十分条件は、
すべての $i < 0$ に対して $A_i = 0$ であることである。

\medskip\noindent
加法圏 $\mathcal{A}$ が与えられたとき、関手
$$
\mathcal{A} \longrightarrow \text{Ch}(\mathcal{A}),\quad
A \longmapsto (\ldots \to 0 \to A \to 0 \to \ldots)
$$
によって $\mathcal{A}$ を、次数 $0$ 以外では零である鎖複体からなる
$\text{Ch}(\mathcal{A})$ の充満部分圏と同一視する。記法を濫用して右辺の対象を
単に $A$ と書くことが多い。$A$ を鎖複体として見ていることを強調したければ、
$A[0]$ という記法を用いることもある。節 \ref{homology-section-homotopy-shift} を参照。

\medskip\noindent
鎖複体の射の組 $f, g : A_\bullet \to B_\bullet$ の間の
{\it ホモトピー $h$} とは、すべての $i$ に対して
$$
f_i - g_i = d_{i + 1} \circ h_i + h_{i - 1} \circ d_i
$$
を満たす射の族 $h_i : A_i \to B_{i + 1}$ のことである。
$f$ と $g$ の間のホモトピーが存在するとき、二つの射
$f, g : A_\bullet \to B_\bullet$ は {\it ホモトピック} であるという。
鎖複体、鎖複体の射、および鎖複体の射の間のホモトピーという概念は、前加法圏でも
明らかに意味をもつ。

\begin{lemma}
\label{homology-lemma-compose-homotopy}
\begin{slogan}
$\text{Ch}(\mathcal{A})$ の Hom 関手はホモトピー関係を保つ。
\end{slogan}
$\mathcal{A}$ を加法圏とする。$f, g : B_\bullet \to C_\bullet$ を鎖複体の射とし、
鎖複体の射 $a : A_\bullet \to B_\bullet$ および
$c : C_\bullet \to D_\bullet$ が与えられているとする。
$\{h_i : B_i \to C_{i + 1}\}$ が $f$ と $g$ の間のホモトピーを定めるなら、
$\{c_{i + 1} \circ h_i \circ a_i\}$ は $c \circ f \circ a$ と
$c \circ g \circ a$ の間のホモトピーを定める。
\end{lemma}

\begin{proof}
省略する。
\end{proof}

\noindent
特に、これは射をホモトピーまで同一視した鎖複体の圏を定義できることを意味する。
これについては後で再び扱う。

\begin{definition}
\label{homology-definition-homotopy-equivalent}
$\mathcal{A}$ を加法圏とする。射 $a : A_\bullet \to B_\bullet$ に対して、
$b \circ a$ と $\text{id}_A$ の間、および $a \circ b$ と $\text{id}_B$ の間に
それぞれホモトピーが存在するような射 $b : B_\bullet \to A_\bullet$ が存在するとき、
その射を {\it ホモトピー同値} という。$A_\bullet$ と $B_\bullet$ の間にそのような
射が存在するとき、$A_\bullet$ と $B_\bullet$ は {\it ホモトピー同値} であるという。
\end{definition}

\noindent
言い換えれば、二つの複体がホモトピー同値であるとは、射をホモトピーまで同一視した
複体の圏においてそれらが同型になることである。

\begin{lemma}
\label{homology-lemma-cat-chain-abelian}
$\mathcal{A}$ をアーベル圏とする。
\begin{enumerate}
\item $\mathcal{A}$ における鎖複体の圏はアーベル圏である。
\item 複体の射 $f : A_\bullet \to B_\bullet$ が単射であるための必要十分条件は、
各 $f_n : A_n \to B_n$ が単射であることである。
\item 複体の射 $f : A_\bullet \to B_\bullet$ が全射であるための必要十分条件は、
各 $f_n : A_n \to B_n$ が全射であることである。
\item 鎖複体の列
$$
A_\bullet \xrightarrow{f} B_\bullet \xrightarrow{g} C_\bullet
$$
が $B_\bullet$ において完全であるための必要十分条件は、各列
$$
A_i \xrightarrow{f_i} B_i \xrightarrow{g_i} C_i
$$
が $B_i$ において完全であることである。
\end{enumerate}
\end{lemma}

\begin{proof}
省略する。
\end{proof}

\noindent
アーベル圏の鎖複体 $A_\bullet$ の、任意の $i \in \mathbf{Z}$ に対する第 $i$ 次
{\it ホモロジー群} を式
$$
H_i(A_\bullet) = \Ker(d_i)/\Im(d_{i + 1}).
$$
で定める。$f : A_\bullet \to B_\bullet$ が $\mathcal{A}$ の鎖複体の射なら、明らかに
$f_i(\Ker(d_i : A_i \to A_{i - 1})) \subset
\Ker(d_i : B_i \to B_{i - 1})$ であり、$\Im(d_{i + 1})$ についても同様なので、
誘導される射 $H_i(f) : H_i(A_\bullet) \to H_i(B_\bullet)$ を得る。
したがって関手
$$
H_i : \text{Ch}(\mathcal{A}) \longrightarrow \mathcal{A}.
$$
を得る。

\begin{definition}
\label{homology-definition-quasi-isomorphism}
$\mathcal{A}$ をアーベル圏とする。
\begin{enumerate}
\item 鎖複体の射 $f : A_\bullet \to B_\bullet$ によって誘導される写像
$H_i(f) : H_i(A_\bullet) \to H_i(B_\bullet)$ がすべての
$i \in \mathbf{Z}$ に対して同型であるとき、その射を {\it 擬同型} という。
\item 鎖複体 $A_\bullet$ のすべてのホモロジー対象 $H_i(A_\bullet)$ が零であるとき、
その鎖複体を {\it 非輪状} という。
\end{enumerate}
\end{definition}


\begin{lemma}
\label{homology-lemma-map-homology-homotopy}
$\mathcal{A}$ をアーベル圏とする。
\begin{enumerate}
\item 写像 $f, g : A_\bullet \to B_\bullet$ がホモトピックなら、誘導される写像
$H_i(f)$ と $H_i(g)$ は等しい。
\item 写像 $f : A_\bullet \to B_\bullet$ がホモトピー同値なら、$f$ は擬同型である。
\end{enumerate}
\end{lemma}

\begin{proof}
省略する。
\end{proof}

\begin{lemma}
\label{homology-lemma-long-exact-sequence-chain}
$\mathcal{A}$ をアーベル圏とする。
$$
0 \to
A_\bullet \to
B_\bullet \to
C_\bullet \to
0
$$
を $\mathcal{A}$ の鎖複体の短完全列とする。このとき標準的な長完全
ホモロジー列
$$
\xymatrix{
\ldots & \ldots & \ldots \ar[lld] \\
H_i(A_\bullet) \ar[r] & H_i(B_\bullet) \ar[r] & H_i(C_\bullet) \ar[lld] \\
H_{i - 1}(A_\bullet) \ar[r] &
H_{i - 1}(B_\bullet) \ar[r] &
H_{i - 1}(C_\bullet) \ar[lld] \\
\ldots & \ldots & \ldots \\
}
$$
が存在する。
\end{lemma}

\begin{proof}
省略する。写像は、図式
$$
\xymatrix{
&
A_i/\Im(d_{A, i + 1}) \ar[r] \ar[d]^{d_{A, i}} &
B_i/\Im(d_{B, i + 1}) \ar[r] \ar[d]^{d_{B, i}} &
C_i/\Im(d_{C, i + 1}) \ar[r] \ar[d]^{d_{C, i}} &
0 \\
0 \ar[r] &
\Ker(d_{A, i - 1}) \ar[r] &
\Ker(d_{B, i - 1}) \ar[r] &
\Ker(d_{C, i - 1}) &
}
$$
に適用した蛇の補題 \ref{homology-lemma-snake} から得られる。
\end{proof}

\noindent
加法圏 $\mathcal{A}$ における {\it 余鎖複体 $A^\bullet$} とは、
$\mathcal{A}$ の複体
$$
\ldots \to
A^{n - 1} \xrightarrow{d^{n - 1}}
A^n \xrightarrow{d^n}
A^{n + 1} \to
\ldots
$$
のことである。言い換えれば、各 $i \in \mathbf{Z}$ に対して
$\mathcal{A}$ の対象 $A^i$ が与えられ、各 $i \in \mathbf{Z}$ に対して
射 $d^i : A^i \to A^{i + 1}$ が与えられ、$d^{i + 1} \circ d^i = 0$ が
すべての $i$ に対して成り立つ。
{\it 余鎖複体の射 $f : A^\bullet \to B^\bullet$} とは、すべての図式
$$
\xymatrix{
A^i \ar[r]_{d^i} \ar[d]_{f^i} & A^{i + 1} \ar[d]^{f^{i + 1}} \\
B^i \ar[r]^{d^i} & B^{i + 1}
}
$$
が可換となる射の族 $f^i : A^i \to B^i$ のことである。
{\it $\mathcal{A}$ の余鎖複体の圏} を $\text{CoCh}(\mathcal{A})$ と書く。
次の形の対象からなる充満部分圏
$$
\ldots \to 0 \to 0 \to A^0 \to A^1 \to A^2 \to \ldots
$$
を $\text{CoCh}_{\geq 0}(\mathcal{A})$ と書く。言い換えれば、余鎖複体
$A^\bullet$ が部分圏 $\text{CoCh}_{\geq 0}(\mathcal{A})$ に属するための必要十分
条件は、すべての $i < 0$ に対して $A^i = 0$ であることである。

\medskip\noindent
加法圏 $\mathcal{A}$ が与えられたとき、関手
$$
\mathcal{A} \longrightarrow \text{CoCh}(\mathcal{A}),\quad
A \longmapsto (\ldots \to 0 \to A \to 0 \to \ldots)
$$
によって $\mathcal{A}$ を、次数 $0$ 以外では零である余鎖複体からなる
$\text{CoCh}(\mathcal{A})$ の充満部分圏と同一視する。記法を濫用して右辺の対象を
単に $A$ と書くことが多い。$A$ を余鎖複体として見ていることを強調したければ、
$A[0]$ という記法を用いることもある。節 \ref{homology-section-homotopy-shift} を参照。

\medskip\noindent
余鎖複体の射の組 $f, g : A^\bullet \to B^\bullet$ の間の
{\it ホモトピー $h$} とは、すべての $i$ に対して
$$
f^i - g^i = d^{i - 1} \circ h^i + h^{i + 1} \circ d^i
$$
を満たす射の族 $h^i : A^i \to B^{i - 1}$ のことである。
$f$ と $g$ の間のホモトピーが存在するとき、二つの射
$f, g : A^\bullet \to B^\bullet$ は {\it ホモトピック} であるという。
余鎖複体、余鎖複体の射、および余鎖複体の射の間のホモトピーという概念は、
前加法圏でも明らかに意味をもつ。

\begin{lemma}
\label{homology-lemma-compose-homotopy-cochain}
$\mathcal{A}$ を加法圏とする。$f, g : B^\bullet \to C^\bullet$ を余鎖複体の射とし、
余鎖複体の射 $a : A^\bullet \to B^\bullet$ および
$c : C^\bullet \to D^\bullet$ が与えられているとする。
$\{h^i : B^i \to C^{i - 1}\}$ が $f$ と $g$ の間のホモトピーを定めるなら、
$\{c^{i - 1} \circ h^i \circ a^i\}$ は $c \circ f \circ a$ と
$c \circ g \circ a$ の間のホモトピーを定める。
\end{lemma}

\begin{proof}
省略する。
\end{proof}

\noindent
特に、これは射をホモトピーまで同一視した余鎖複体の圏を定義できることを意味する。
これについては後で再び扱う。

\begin{definition}
\label{homology-definition-homotopy-equivalent-cochain}
$\mathcal{A}$ を加法圏とする。射 $a : A^\bullet \to B^\bullet$ に対して、
$b \circ a$ と $\text{id}_A$ の間、および $a \circ b$ と $\text{id}_B$ の間に
それぞれホモトピーが存在するような射 $b : B^\bullet \to A^\bullet$ が存在するとき、
その射を {\it ホモトピー同値} という。$A^\bullet$ と $B^\bullet$ の間にそのような
射が存在するとき、$A^\bullet$ と $B^\bullet$ は {\it ホモトピー同値} であるという。
\end{definition}

\noindent
言い換えれば、二つの複体がホモトピー同値であるとは、射をホモトピーまで同一視した
複体の圏においてそれらが同型になることである。

\begin{lemma}
\label{homology-lemma-cat-cochain-abelian}
$\mathcal{A}$ をアーベル圏とする。
\begin{enumerate}
\item $\mathcal{A}$ における余鎖複体の圏はアーベル圏である。
\item 余鎖複体の射 $f : A^\bullet \to B^\bullet$ が単射であるための必要十分条件は、
各 $f^n : A^n \to B^n$ が単射であることである。
\item 余鎖複体の射 $f : A^\bullet \to B^\bullet$ が全射であるための必要十分条件は、
各 $f^n : A^n \to B^n$ が全射であることである。
\item 余鎖複体の列
$$
A^\bullet \xrightarrow{f} B^\bullet \xrightarrow{g} C^\bullet
$$
が $B^\bullet$ において完全であるための必要十分条件は、各列
$$
A^i \xrightarrow{f^i} B^i \xrightarrow{g^i} C^i
$$
が $B^i$ において完全であることである。
\end{enumerate}
\end{lemma}

\begin{proof}
省略する。
\end{proof}

\noindent
余鎖複体 $A^\bullet$ の、任意の $i \in \mathbf{Z}$ に対する第 $i$ 次
{\it コホモロジー群} を式
$$
H^i(A^\bullet) = \Ker(d^i)/\Im(d^{i - 1}).
$$
で定める。$f : A^\bullet \to B^\bullet$ が $\mathcal{A}$ の余鎖複体の射なら、
明らかに
$f^i(\Ker(d^i : A^i \to A^{i + 1})) \subset
\Ker(d^i : B^i \to B^{i + 1})$ であり、$\Im(d^{i - 1})$ についても同様なので、
誘導される射 $H^i(f) : H^i(A^\bullet) \to H^i(B^\bullet)$ を得る。
したがって関手
$$
H^i : \text{CoCh}(\mathcal{A}) \longrightarrow \mathcal{A}.
$$
を得る。

\begin{definition}
\label{homology-definition-quasi-isomorphism-cochain}
$\mathcal{A}$ をアーベル圏とする。
\begin{enumerate}
\item $\mathcal{A}$ の余鎖複体の射 $f : A^\bullet \to B^\bullet$ によって
誘導される写像 $H^i(f) : H^i(A^\bullet) \to H^i(B^\bullet)$ が、すべての
$i \in \mathbf{Z}$ に対して同型であるとき、その射を {\it 擬同型} という。
\item 余鎖複体 $A^\bullet$ のすべてのコホモロジー対象 $H^i(A^\bullet)$ が零である
とき、その余鎖複体を {\it 非輪状} という。
\end{enumerate}
\end{definition}

\begin{lemma}
\label{homology-lemma-map-cohomology-homotopy-cochain}
$\mathcal{A}$ をアーベル圏とする。
\begin{enumerate}
\item 写像 $f, g : A^\bullet \to B^\bullet$ がホモトピックなら、誘導される写像
$H^i(f)$ と $H^i(g)$ は等しい。
\item $f : A^\bullet \to B^\bullet$ がホモトピー同値なら、$f$ は擬同型である。
\end{enumerate}
\end{lemma}

\begin{proof}
省略する。
\end{proof}

\begin{lemma}
\label{homology-lemma-long-exact-sequence-cochain}
\begin{slogan}
複体の短完全列から（コ）ホモロジーの長完全列が得られる。
\end{slogan}
$\mathcal{A}$ をアーベル圏とする。
$$
0 \to
A^\bullet \to
B^\bullet \to
C^\bullet \to
0
$$
を $\mathcal{A}$ の余鎖複体の短完全列とする。このとき長完全コホモロジー列
$$
\xymatrix{
\ldots & \ldots & \ldots \ar[lld] \\
H^i(A^\bullet) \ar[r] &
H^i(B^\bullet) \ar[r] &
H^i(C^\bullet) \ar[lld] \\
H^{i + 1}(A^\bullet) \ar[r] &
H^{i + 1}(B^\bullet) \ar[r] &
H^{i + 1}(C^\bullet) \ar[lld] \\
\ldots & \ldots & \ldots \\
}
$$
が存在する。この構成は短完全列に関して関手的であり、定義
\ref{homology-definition-cohomology-shift} の意味でシフトと両立する長完全コホモロジー列を
与える。
\end{lemma}

\begin{proof}
横向きの写像 $H^i(A^\bullet) \to H^i(B^\bullet)$ と
$H^i(B^\bullet) \to H^i(C^\bullet)$ については、上で見たように $H^i$ が関手で
あることを用いる。「境界写像」$H^i(C^\bullet) \to H^{i + 1}(A^\bullet)$ には、
図式
$$
\xymatrix{
&
A^i/\Im(d_A^{i - 1}) \ar[r] \ar[d]^{d_A^i} &
B^i/\Im(d_B^{i - 1}) \ar[r] \ar[d]^{d_B^i} &
C^i/\Im(d_C^{i - 1}) \ar[r] \ar[d]^{d_C^i} &
0 \\
0 \ar[r] &
\Ker(d_A^{i + 1}) \ar[r] &
\Ker(d_B^{i + 1}) \ar[r] &
\Ker(d_C^{i + 1}) &
}
$$
に適用した蛇の補題 \ref{homology-lemma-snake} の写像 $\delta$ を用いる。右の縦写像の核は
$H^i(C^\bullet)$ に等しく、左の縦写像の余核は
$H^{i + 1}(A^\bullet)$ に等しいので、これは正しく定義される。長い列の完全性は
補題 \ref{homology-lemma-snake} の (2) の完全性である。関手性は補題
\ref{homology-lemma-snake-natural} であり、シフトとの両立性は定義から直ちに従う。
\end{proof}





\section{ホモトピーとシフト関手}
\label{homology-section-homotopy-shift}

\noindent
ホモロジー代数を論じるときには常に符号と添字を扱わなければならず、これは厄介な
特徴である\footnote{符号の誤りに気づいた場合、または規約を改善できる場合には、
ぜひ知らせてほしい。}。

\begin{definition}
\label{homology-definition-shift}
$\mathcal{A}$ を加法圏とし、$A_\bullet$ を境界写像
$d_{A, n} : A_n \to A_{n - 1}$ をもつ鎖複体とする。任意の
$k \in \mathbf{Z}$ に対して {\it $k$ シフト鎖複体 $A[k]_\bullet$} を次のように定める。
\begin{enumerate}
\item $A[k]_n = A_{n + k}$ と置く。
\item $d_{A[k], n} : A[k]_n \to A[k]_{n - 1}$ を
$d_{A[k], n} = (-1)^k d_{A, n + k}$ と置く。
\end{enumerate}
$f : A_\bullet \to B_\bullet$ が鎖複体の射なら、
$f[k] : A[k]_\bullet \to B[k]_\bullet$ を
$f[k]_n = f_{k + n}$ で与えられる鎖複体の射とする。
\end{definition}

\noindent
これは、互いに（関手の同型を介さず文字どおり）可換する関手
$[k] : \text{Ch}(\mathcal{A}) \to \text{Ch}(\mathcal{A})$ をもち、
$A[k][l]_\bullet = A[k + l]_\bullet$ および
$[0] = \text{id}_{\text{Ch}(\mathcal{A})}$ が成り立つことを意味する。

\medskip\noindent
節 \ref{homology-section-complexes} で見たように、$\mathcal{A}$ を
$\text{Ch}(\mathcal{A})$ の充満部分圏とみなすことを思い出す。したがって
$\mathcal{A}$ の任意の対象 $A$ について、記法 $A[k]$ は次数 $-k$ に $A$ をもち、
それ以外のすべての次数で零となる一意な鎖複体を表す。

\begin{definition}
\label{homology-definition-homology-shift}
$\mathcal{A}$ をアーベル圏とし、$A_\bullet$ を境界写像
$d_{A, n} : A_n \to A_{n - 1}$ をもつ鎖複体とする。任意の
$k \in \mathbf{Z}$ に対して、同一視 $A_{i + k} = A[k]_i$ により
{\it $H_{i + k}(A_\bullet) \rightarrow H_i(A[k]_\bullet)$} と同一視する。
\end{definition}

\noindent
この同一視は $A_\bullet$ に関して関手的である。この定義には符号が関与しないので、
実際、すべてのホモロジー対象 $H_{i - k}(A[k]_\bullet)$ の同一視からなる両立系を
得る。この系はさらに、同一視 $A[k][l]_\bullet = A[k + l]_\bullet$ および
$[0] = \text{id}_{\text{Ch}(\mathcal{A})}$ と両立する。

\medskip\noindent
$\mathcal{A}$ を加法圏とする。$A_\bullet$ と $B_\bullet$ を鎖複体、
$a, b : A_\bullet \to B_\bullet$ を鎖複体の射とし、
$\{h_i : A_i \to B_{i + 1}\}$ を $a$ と $b$ の間のホモトピーとする。これは
$a_i - b_i = d_{i + 1} \circ h_i + h_{i - 1} \circ d_i$ を意味することを思い出す。
$a = b$ なら、式 $0 = d_{i + 1} \circ h_i + h_{i - 1} \circ d_i$、
言い換えれば $ - d_{i + 1} \circ h_i = h_{i - 1} \circ d_i$ を得る。
上の定義により、これは族 $\{h_i\}$ が鎖複体の射
$$
A_\bullet \longrightarrow B[1]_\bullet.
$$
を定めることを意味する。上の注意により、これは射
$A[-1]_\bullet \to B_\bullet$ と同じものである。以上から次の補題を得る。

\begin{lemma}
\label{homology-lemma-homotopy-shift}
$\mathcal{A}$ を加法圏とし、$A_\bullet$ と $B_\bullet$ を鎖複体とする。
任意の鎖複体の射 $a : A_\bullet \to B_\bullet$ に対して、$a$ から $a$ への
ホモトピーの集合と
$\Mor_{\text{Ch}(\mathcal{A})}(A_\bullet, B[1]_\bullet)$ の間に全単射がある。
より一般に、$a$ と $b$ の間のホモトピーの集合は空であるか、または群
$\Mor_{\text{Ch}(\mathcal{A})}(A_\bullet, B[1]_\bullet)$ の下の主等質空間である。
\end{lemma}

\begin{proof}
上を参照。
\end{proof}

\begin{lemma}
\label{homology-lemma-ses-termwise-split}
$\mathcal{A}$ をアーベル圏とし、
$$
0 \to A_\bullet \to B_\bullet \to C_\bullet \to 0
$$
を複体の短完全列とする。短完全列
$0 \to A_n \to B_n \to C_n \to 0$ を分裂させる射の族
$\{s_n : C_n \to B_n\}$ が与えられているとする。対応する射影を
$\pi_n : B_n \to A_n$ とする。補題 \ref{homology-lemma-ses-split} を参照。このとき射の族
$$
\pi_{n - 1} \circ d_{B, n} \circ s_n
:
C_n \to A_{n - 1}
$$
は複体の射 $\delta(s) : C_\bullet \to A[-1]_\bullet$ を定める。
\end{lemma}

\begin{proof}
短完全列の複体の射を $i : A_\bullet \to B_\bullet$ および
$q : B_\bullet \to C_\bullet$ と書く。このとき
$i_{n - 1} \circ \pi_{n - 1} \circ d_{B, n} \circ s_n =
d_{B, n} \circ s_n - s_{n - 1} \circ d_{C, n}$ である。したがって所望のとおり
$i_{n - 2} \circ d_{A, n - 1} \circ \pi_{n - 1} \circ d_{B, n} \circ s_n =
d_{B, n - 1} \circ (d_{B, n} \circ s_n - s_{n - 1} \circ d_{C, n}) =
- d_{B, n - 1} \circ s_{n - 1} \circ d_{C, n}$ である。
\end{proof}

\begin{lemma}
\label{homology-lemma-ses-termwise-split-long}
記法と仮定は上の補題 \ref{homology-lemma-ses-termwise-split} と同じとする。
複体の射 $\delta(s) : C_\bullet \to A[-1]_\bullet$ は写像
$$
H_i(\delta(s)) :
H_i(C_\bullet) \longrightarrow H_i(A[-1]_\bullet) = H_{i - 1}(A_\bullet)
$$
を誘導し、これは補題 \ref{homology-lemma-long-exact-sequence-chain} によって鎖複体の
短完全列に付随する長完全ホモロジー列に現れる写像である。
\end{lemma}

\begin{proof}
省略する。
\end{proof}

\begin{lemma}
\label{homology-lemma-ses-termwise-split-homotopy}
記法と仮定は上の補題 \ref{homology-lemma-ses-termwise-split} と同じとする。
$\{s'_n : C_n \to B_n\}$ を第二の分裂の選択とする。一意な射
$h_n : C_n \to A_n$ により $s'_n = s_n + i_n \circ h_n$ と書く。
写像の族 $\{h_n : C_n \to A[-1]_{n + 1}\}$ は、対応する複体の射
$\delta(s), \delta(s') : C_\bullet \to A[-1]_\bullet$ の間のホモトピーである。
\end{lemma}

\begin{proof}
省略する。
\end{proof}

\begin{definition}
\label{homology-definition-shift-cochain}
$\mathcal{A}$ を加法圏とし、$A^\bullet$ を境界写像
$d_A^n : A^n \to A^{n + 1}$ をもつ余鎖複体とする。任意の
$k \in \mathbf{Z}$ に対して {\it $k$ シフト余鎖複体 $A[k]^\bullet$} を次のように
定める。
\begin{enumerate}
\item $A[k]^n = A^{n + k}$ と置く。
\item $d_{A[k]}^n : A[k]^n \to A[k]^{n + 1}$ を
$d_{A[k]}^n = (-1)^k d_A^{n + k}$ と置く。
\end{enumerate}
$f : A^\bullet \to B^\bullet$ が余鎖複体の射なら、
$f[k] : A[k]^\bullet \to B[k]^\bullet$ を
$f[k]^n = f^{k + n}$ で与えられる余鎖複体の射とする。
\end{definition}

\noindent
これは、互いに（関手の同型を介さず文字どおり）可換する関手
$[k] : \text{CoCh}(\mathcal{A}) \to \text{CoCh}(\mathcal{A})$ をもち、
$A[k][l]^\bullet = A[k + l]^\bullet$ および
$[0] = \text{id}_{\text{CoCh}(\mathcal{A})}$ が成り立つことを意味する。

\medskip\noindent
節 \ref{homology-section-complexes} で見たように、$\mathcal{A}$ を
$\text{CoCh}(\mathcal{A})$ の充満部分圏とみなすことを思い出す。したがって
$\mathcal{A}$ の任意の対象 $A$ について、記法 $A[k]$ は次数 $-k$ に $A$ をもち、
それ以外のすべての次数で零となる一意な余鎖複体を表す。

\begin{definition}
\label{homology-definition-cohomology-shift}
$\mathcal{A}$ をアーベル圏とし、$A^\bullet$ を境界写像
$d_A^n : A^n \to A^{n + 1}$ をもつ余鎖複体とする。任意の
$k \in \mathbf{Z}$ に対して、同一視 $A^{i + k} = A[k]^i$ により
{\it $H^{i + k}(A^\bullet) \longrightarrow H^i(A[k]^\bullet)$} と同一視する。
\end{definition}

\noindent
この同一視は $A^\bullet$ に関して関手的である。この定義には符号が関与しないので、
実際、すべてのホモロジー対象 $H^{i - k}(A[k]^\bullet)$ の同一視からなる両立系を
得る。この系はさらに、同一視 $A[k][l]^\bullet = A[k + l]^\bullet$ および
$[0] = \text{id}_{\text{CoCh}(\mathcal{A})}$ と両立する。

\medskip\noindent
$\mathcal{A}$ を加法圏とする。$A^\bullet$ と $B^\bullet$ を余鎖複体、
$a, b : A^\bullet \to B^\bullet$ を余鎖複体の射とし、
$\{h^i : A^i \to B^{i - 1}\}$ を $a$ と $b$ の間のホモトピーとする。これは
$a^i - b^i = d^{i - 1} \circ h^i + h^{i + 1} \circ d^i$ を意味することを思い出す。
$a = b$ なら、式 $0 = d^{i - 1} \circ h^i + h^{i + 1} \circ d^i$、
言い換えれば $ - d^{i - 1} \circ h^i = h^{i + 1} \circ d^i$ を得る。
上の定義により、これは族 $\{h^i\}$ が余鎖複体の射
$$
A^\bullet \longrightarrow B[-1]^\bullet.
$$
を定めることを意味する。上の注意により、これは射
$A[1]^\bullet \to B^\bullet$ と同じものである。以上から次の補題を得る。

\begin{lemma}
\label{homology-lemma-homotopy-shift-cochain}
$\mathcal{A}$ を加法圏とし、$A^\bullet$ と $B^\bullet$ を余鎖複体とする。
任意の余鎖複体の射 $a : A^\bullet \to B^\bullet$ に対して、$a$ から $a$ への
ホモトピーの集合と
$\Mor_{\text{CoCh}(\mathcal{A})}(A^\bullet, B[-1]^\bullet)$ の間に全単射がある。
より一般に、$a$ と $b$ の間のホモトピーの集合は空であるか、または群
$\Mor_{\text{CoCh}(\mathcal{A})}(A^\bullet, B[-1]^\bullet)$ の下の主等質空間である。
\end{lemma}

\begin{proof}
上を参照。
\end{proof}

\begin{lemma}
\label{homology-lemma-ses-termwise-split-cochain}
$\mathcal{A}$ を加法圏とし、
$$
0 \to A^\bullet \to B^\bullet \to C^\bullet \to 0
$$
を複体の複体（！）とする。表示された列の写像と両立する分裂
$B^n = A^n \oplus C^n$ が与えられているとする。対応する写像を
$s^n : C^n \to B^n$ および $\pi^n : B^n \to A^n$ とする。このとき射の族
$$
\pi^{n + 1} \circ d_B^n \circ s^n
:
C^n \to A^{n + 1}
$$
は複体の射 $\delta : C^\bullet \to A[1]^\bullet$ を定める。
\end{lemma}

\begin{proof}
短完全列の複体の射を $i : A^\bullet \to B^\bullet$ および
$q : B^\bullet \to C^\bullet$ と書く。このとき
$i^{n + 1} \circ \pi^{n + 1} \circ d_B^n \circ s^n =
d_B^n \circ s^n - s^{n + 1} \circ d_C^n$ である。したがって所望のとおり
$i^{n + 2} \circ d_A^{n + 1} \circ \pi^{n + 1} \circ d_B^n \circ s^n =
d_B^{n + 1} \circ (d_B^n \circ s^n - s^{n + 1} \circ d_C^n) =
- d_B^{n + 1} \circ s^{n + 1} \circ d_C^n$ である。
\end{proof}

\begin{lemma}
\label{homology-lemma-ses-termwise-split-long-cochain}
記法と仮定は上の補題 \ref{homology-lemma-ses-termwise-split-cochain} と同じとし、さらに
$\mathcal{A}$ はアーベル圏であると仮定する。複体の射
$\delta : C^\bullet \to A[1]^\bullet$ は写像
$$
H^i(\delta) :
H^i(C^\bullet) \longrightarrow H^i(A[1]^\bullet) = H^{i + 1}(A^\bullet)
$$
を誘導し、これは補題 \ref{homology-lemma-long-exact-sequence-cochain} によって余鎖複体の
短完全列に付随する長完全ホモロジー列に現れる写像である。
\end{lemma}

\begin{proof}
省略する。
\end{proof}

\begin{lemma}
\label{homology-lemma-ses-termwise-split-homotopy-cochain}
記法と仮定は補題 \ref{homology-lemma-ses-termwise-split-cochain} と同じとする。
$\alpha : A^\bullet \to B^\bullet$ と
$\beta : B^\bullet \to C^\bullet$ を与えられた複体の射とする。
$(s')^n : C^n \to B^n$ と $(\pi')^n : B^n \to A^n$ を第二の分裂の選択とする。
一意な射 $h^n : C^n \to A^n$ および $g^n : C^n \to A^n$ によって
$(s')^n = s^n + \alpha^n \circ h^n$ および
$(\pi')^n = \pi^n + g^n \circ \beta^n$ と書く。このとき
\begin{enumerate}
\item $g^n = - h^n$ である。
\item 写像の族 $\{g^n : C^n \to A[1]^{n - 1}\}$ は
$\delta, \delta' : C^\bullet \to A[1]^\bullet$ の間のホモトピーである。より正確には
$(\delta')^n = \delta^n + g^{n + 1} \circ d_C^n + d_{A[1]}^{n - 1} \circ g^n$ である。
\end{enumerate}
\end{lemma}

\begin{proof}
$(s')^n$ と $(\pi')^n$ は分裂なので $(\pi')^n \circ (s')^n = 0$ である。よって
$$
0 = ( \pi^n + g^n \circ \beta^n ) \circ ( s^n + \alpha^n \circ h^n ) =
g^n \circ \beta^n \circ s^n + \pi^n \circ \alpha^n \circ h^n =
g^n + h^n
$$
であり、(1) が示された。$(\delta')^n$ を次のように計算する。
$$
( \pi^{n + 1} + g^{n + 1} \circ \beta^{n + 1} )
\circ d_B^n \circ
( s^n + \alpha^n \circ h^n )
= \delta^n + g^{n + 1} \circ d_C^n + d_A^n \circ h^n
$$
$h^n = -g^n$ かつ $d_{A[1]}^{n - 1} = -d_A^n$ なので (2) が従う。
\end{proof}





\section{複体の切断}
\label{homology-section-truncations}

\noindent
$\mathcal{A}$ をアーベル圏とし、$A_\bullet$ を鎖複体とする。複体
$A_\bullet$ を {\it 切断} する方法はいくつかある。
\begin{enumerate}
\item {\it 「素朴な」切断 $\sigma_{\leq n}$} は、規則
$(\sigma_{\leq n} A_\bullet)_i = 0$（$i > n$ のとき）および
$(\sigma_{\leq n} A_\bullet)_i = A_i$（$i \leq n$ のとき）によって定められる
部分複体 $\sigma_{\leq n} A_\bullet$ である。図示すると
$$
\xymatrix{
\sigma_{\leq n}A_\bullet \ar[d]  &
\ldots \ar[r] &
0 \ar[r] \ar[d] &
A_n \ar[r] \ar[d] &
A_{n - 1} \ar[r] \ar[d] &
\ldots \\
A_\bullet  &
\ldots \ar[r] &
A_{n + 1} \ar[r] &
A_n \ar[r] &
A_{n - 1} \ar[r] &
\ldots
}
$$
となる。性質
$\sigma_{\leq n}A_\bullet / \sigma_{\leq n - 1}A_\bullet = A_n[-n]$ に注意する。
\item {\it 「素朴な」切断 $\sigma_{\geq n}$} は、規則
$(\sigma_{\geq n} A_\bullet)_i = A_i$（$i \geq n$ のとき）および
$(\sigma_{\geq n} A_\bullet)_i = 0$（$i < n$ のとき）によって定められる商複体
$\sigma_{\geq n} A_\bullet$ である。図示すると
$$
\xymatrix{
A_\bullet \ar[d]  &
\ldots \ar[r] &
A_{n + 1} \ar[r] \ar[d] &
A_n \ar[r] \ar[d] &
A_{n - 1} \ar[r] \ar[d] &
\ldots \\
\sigma_{\geq n}A_\bullet  &
\ldots \ar[r] &
A_{n + 1} \ar[r] &
A_n \ar[r] &
0 \ar[r] &
\ldots
}
$$
となる。複体の写像
$\sigma_{\geq n}A_\bullet \to \sigma_{\geq n + 1}A_\bullet$ は全射であり、
その核は $A_n[-n]$ である。
\item {\it 標準切断} $\tau_{\geq n}A_\bullet$ は図式
$$
\xymatrix{
\tau_{\geq n}A_\bullet \ar[d]  &
\ldots \ar[r] &
A_{n + 1} \ar[r] \ar[d] &
\Ker(d_n) \ar[r] \ar[d] &
0 \ar[r] \ar[d] &
\ldots \\
A_\bullet  &
\ldots \ar[r] &
A_{n + 1} \ar[r] &
A_n \ar[r] &
A_{n - 1} \ar[r] &
\ldots
}
$$
によって定められる。これらの複体は性質
$$
H_i(\tau_{\geq n}A_\bullet) =
\left\{
\begin{matrix}
H_i(A_\bullet) & \text{if} & i \geq n \\
0 & \text{if} & i < n
\end{matrix}
\right.
$$
をもつ。
\item {\it 標準切断} $\tau_{\leq n}A_\bullet$ は図式
$$
\xymatrix{
A_\bullet \ar[d]  &
\ldots \ar[r] &
A_{n + 1} \ar[r] \ar[d] &
A_n \ar[r] \ar[d] &
A_{n - 1} \ar[r] \ar[d] &
\ldots \\
\tau_{\leq n}A_\bullet  &
\ldots \ar[r] &
0 \ar[r] &
\Coker(d_{n + 1}) \ar[r] &
A_{n - 1} \ar[r] &
\ldots
}
$$
によって定められる。これらの複体は性質
$$
H_i(\tau_{\leq n}A_\bullet) =
\left\{
\begin{matrix}
H_i(A_\bullet) & \text{if} & i \leq n \\
0 & \text{if} & i > n
\end{matrix}
\right.
$$
をもつ。
\end{enumerate}

\noindent
$\mathcal{A}$ をアーベル圏とし、$A^\bullet$ を余鎖複体とする。複体
$A^\bullet$ を切断する方法は四つある。
\begin{enumerate}
\item {\it 「素朴な」切断 $\sigma_{\geq n}$} は、規則
$(\sigma_{\geq n} A^\bullet)^i = 0$（$i < n$ のとき）および
$(\sigma_{\geq n} A^\bullet)^i = A^i$（$i \geq n$ のとき）によって定められる
部分複体 $\sigma_{\geq n} A^\bullet$ である。図示すると
$$
\xymatrix{
\sigma_{\geq n}A^\bullet \ar[d]  &
\ldots \ar[r] &
0 \ar[r] \ar[d] &
A^n \ar[r] \ar[d] &
A^{n + 1} \ar[r] \ar[d] &
\ldots \\
A^\bullet  &
\ldots \ar[r] &
A^{n - 1} \ar[r] &
A^n \ar[r] &
A^{n + 1} \ar[r] &
\ldots
}
$$
となる。性質
$\sigma_{\geq n}A^\bullet / \sigma_{\geq n + 1}A^\bullet
= A^n[-n]$ に注意する。
\item {\it 「素朴な」切断 $\sigma_{\leq n}$} は、規則
$(\sigma_{\leq n} A^\bullet)^i = 0$（$i > n$ のとき）および
$(\sigma_{\leq n} A^\bullet)^i = A^i$（$i \leq n$ のとき）によって定められる
商複体 $\sigma_{\leq n} A^\bullet$ である。図示すると
$$
\xymatrix{
A^\bullet \ar[d]  &
\ldots \ar[r] &
A^{n - 1} \ar[r] \ar[d] &
A^n \ar[r] \ar[d] &
A^{n + 1} \ar[r] \ar[d] &
\ldots \\
\sigma_{\leq n}A^\bullet &
\ldots \ar[r] &
A^{n - 1} \ar[r] &
A^n \ar[r] &
0 \ar[r] &
\ldots \\
}
$$
となる。複体の写像
$\sigma_{\leq n}A^\bullet \to \sigma_{\leq n - 1}A^\bullet$ は全射であり、
その核は $A^n[-n]$ である。
\item {\it 標準切断} $\tau_{\leq n}A^\bullet$ は図式
$$
\xymatrix{
\tau_{\leq n}A^\bullet \ar[d]  &
\ldots \ar[r] &
A^{n - 1} \ar[r] \ar[d] &
\Ker(d^n) \ar[r] \ar[d] &
0 \ar[r] \ar[d] &
\ldots \\
A^\bullet  &
\ldots \ar[r] &
A^{n - 1} \ar[r] &
A^n \ar[r] &
A^{n + 1} \ar[r] &
\ldots
}
$$
によって定められる。これらの複体は性質
$$
H^i(\tau_{\leq n}A^\bullet) =
\left\{
\begin{matrix}
H^i(A^\bullet) & \text{if} & i \leq n \\
0 & \text{if} & i > n
\end{matrix}
\right.
$$
をもつ。
\item {\it 標準切断} $\tau_{\geq n}A^\bullet$ は図式
$$
\xymatrix{
A^\bullet \ar[d] &
\ldots \ar[r] &
A^{n - 1} \ar[r] \ar[d] &
A^n \ar[r] \ar[d] &
A^{n + 1} \ar[r] \ar[d] &
\ldots \\
\tau_{\geq n}A^\bullet &
\ldots \ar[r] &
0 \ar[r] &
\Coker(d^{n - 1}) \ar[r] &
A^{n + 1} \ar[r] &
\ldots
}
$$
によって定められる。これらの複体は性質
$$
H^i(\tau_{\geq n}A^\bullet) =
\left\{
\begin{matrix}
0 & \text{if} & i < n \\
H^i(A^\bullet) & \text{if} & i \geq n
\end{matrix}
\right.
$$
をもつ。
\end{enumerate}

\section{次数付き対象}
\label{homology-section-graded}

\noindent
次の定義を置く。

\begin{definition}
\label{homology-definition-graded}
$\mathcal{A}$ を加法圏とする。$\mathcal{A}$ の{\it 次数付き対象の圏}
$\text{Gr}(\mathcal{A})$ とは、次の対象と射をもつ圏である。
\begin{enumerate}
\item 対象 $A = (A^i)$ は、$\mathcal{A}$ の対象 $A^i$ を
$i \in \mathbf{Z}$ で添字付けた族であり、
\item 射 $f : A = (A^i) \to B = (B^i)$ は、$\mathcal{A}$ の射
$f^i : A^i \to B^i$ の族である。
\end{enumerate}
\end{definition}

\noindent
$\mathcal{A}$ が可算直和をもつとき、$\text{Gr}(\mathcal{A})$ の対象
$A = (A^i)$ に対象
$$
A = \bigoplus\nolimits_{i \in \mathbf{Z}} A^i
$$
を対応させ、$k^iA = A^i$ と置ける。この場合、$\text{Gr}(\mathcal{A})$ は、
$\mathcal{A}$ の対象 $A$ と、$\mathbf{Z}$ で添字付けられた直和因子による直和分解
$$
A = \bigoplus\nolimits_{i \in \mathbf{Z}} k^iA
$$
からなる対 $(A, k)$ の圏と同値である。このような対象の射
$(A, k) \to (B, k)$ は、すべての $i \in \mathbf{Z}$ に対して
$\varphi(k^iA) \subset k^iB$ を満たす $\mathcal{A}$ の射
$\varphi : A \to B$ によって与えられる。加法圏 $\mathcal{A}$ が可算直和をもつ
場合には、以後とくに断ることなくこの同値を用いる。

\medskip\noindent
しかし、本書の定義では、加法圏またはアーベル圏が必ずしもすべての
（可算）直和をもつとは限らない。この場合にも上の定義は意味をもつ。
例えば、$\mathcal{A} = \text{Vect}_k$ を体 $k$ 上の有限次元ベクトル空間の圏と
すると、$\text{Gr}(\text{Vect}_k)$ は、各次数成分が有限次元であるような、
指定された次数付けをもつベクトル空間の圏であって、指定された次数付けをもつ
有限次元ベクトル空間の圏ではない。

\begin{lemma}
\label{homology-lemma-graded}
$\mathcal{A}$ をアーベル圏とする。次数付き対象の圏
$\text{Gr}(\mathcal{A})$ はアーベル圏である。
\end{lemma}

\begin{proof}
$f : A = (A^i) \to B = (B^i)$ を、$\mathcal{A}$ の射
$f^i : A^i \to B^i$ によって与えられる $\mathcal{A}$ の次数付き対象の射とする。
このとき、圏 $\text{Gr}(\mathcal{A})$ において
$\Ker(f) = (\Ker(f^i))$ および $\Coker(f) = (\Coker(f^i))$ である。
$\mathcal{A}$ では $\Im = \Coim$ が成り立つので、
$\text{Gr}(\mathcal{A})$ でも同じことが成り立つ。
\end{proof}

\begin{remark}[注意]
\label{homology-remark-direct-sums-not-exact}
可算直和をもつアーベル圏 $\mathcal{A}$ であって、可算直和が完全でないものが
存在する。一例は、$\mathbf{R}$ 上のアーベル群の層の圏の反対圏である。
実際、$\mathbf{R}$ 上のアーベル群の層の圏は可算直積をもつが、可算直積は
完全ではない。そのような圏に対しては、上で述べた関手
$\text{Gr}(\mathcal{A}) \to \mathcal{A}$、
$(A^i) \mapsto \bigoplus A^i$ は完全ではない。
それでも $\text{Gr}(\mathcal{A})$ が $\mathcal{A}$ の次数付き対象 $(A, k)$ の圏と
同値であることは変わらないが、次数付き対象の射
$\varphi : (A, k) \to (B, k)$ の次数付き対象の圏における核は、直和分解を備えた
$\Ker(\varphi)$ とは等しくなく、むしろ写像 $k^iA \to k^iB$ の核の直和である。
\end{remark}

\begin{definition}
\label{homology-definition-graded-shift}
$\mathcal{A}$ を加法圏とする。$A = (A^i)$ が次数付き対象ならば、その第 $k$ {\it シフト}
$A[k]$ とは、$A[k]^i = A^{k + i}$ で与えられる次数付き対象である。
\end{definition}

\noindent
$A$ と $B$ が $\mathcal{A}$ の次数付き対象ならば、
\begin{equation}
\label{homology-equation-hom-into-shift}
\Hom_{\text{Gr}(\mathcal{A})}(A, B[k]) =
\Hom_{\text{Gr}(\mathcal{A})}(A[-k], B)
\end{equation}
であり、この群の元を{\it 次数 $k$ の斉次な}次数付き対象の写像と呼ぶことがある。

\medskip\noindent
任意の集合 $G$ に対し、$\mathcal{A}$ の $G$ 次数付き対象の圏を、対象が
$G$ の元でパラメータ付けられた対象の族 $A = (A^g)_{g \in G}$ である圏として
定義できる。射 $f : A \to B$ は、$g$ が $G$ の元を走るときの写像
$f^g : A^g \to B^g$ の族として定義する。$G$ がアーベル群ならば、
規則 $(A[g])^{g_0} = A^{g + g_0}$ によって $G$ 次数付き対象の圏上の
シフト関手 $[g]$ を曖昧さなく定義できる。この種の構成の特別な場合が
$G = \mathbf{Z} \times \mathbf{Z}$ である。この場合、圏の対象を
$\mathcal{A}$ の{\it 二重次数付き}対象という。二重次数付き対象 $A$ の
$(p, q)$ 成分は通常 $A^{p, q}$ と表す。
$(a, b) \in \mathbf{Z} \times \mathbf{Z}$ に対して、$A[(a, b)]$ の代わりに
$A[a, b]$ と書く。射 $A \to A[a, b]$ を{\it 双次数 $(a, b)$ の写像}と
呼ぶことがある。

\section{加法モノイド圏}
\label{homology-section-monoidal}

\noindent
圏上のモノイド構造と加法構造との相互作用について、いくつか述べる。

\begin{definition}
\label{homology-definition-additive-monoidal}
{\it 加法モノイド圏}とは、モノイド構造 $\otimes, \phi$
（圏論、定義 \ref{categories-definition-monoidal-category}）を備えた加法圏
$\mathcal{A}$ であって、$\otimes$ が各変数について加法関手となるものをいう。
\end{definition}

\begin{lemma}
\label{homology-lemma-additive-dual}
$\mathcal{A}$ を加法モノイド圏とする。$Y_i$（$i = 1, 2$）が
$X_i$（$i = 1, 2$）の左双対ならば、$Y_1 \oplus Y_2$ は
$X_1 \oplus X_2$ の左双対である。
\end{lemma}

\begin{proof}
随伴の一意性と、圏論、注意 \ref{categories-remark-left-dual-adjoint} から従う。
\end{proof}

\begin{lemma}
\label{homology-lemma-Karoubian-dual}
カルービ圏である加法モノイド圏では、左双対をもつ対象の任意の直和因子も
左双対をもつ。
\end{lemma}

\begin{proof}
以下では圏論、補題 \ref{categories-lemma-left-dual} を断りなく用いる。
$X$ を、左双対 $Y$ をもつ対象とする。このとき
$$
\Hom(X, X) = \Hom(\mathbf{1}, X \otimes Y) = \Hom(Y, Y)
$$
である。$a : X \to X$ が $b : Y \to Y$ に対応するとき、$b$ は、
$$
\Hom(Z' \otimes X, Z) = \Hom(Z', Z \otimes Y)
$$
において $a$ を前合成することが $1 \otimes b$ を後合成することと一致するような、
$Y$ の一意な自己準同型である。したがって全単射
$\Hom(X, X) \to \Hom(Y, Y)$、$a \mapsto b$ は、代数 $\Hom(X, X)$ の
反対代数から代数 $\Hom(Y, Y)$ への同型である。特に、
$X = X_1 \oplus X_2$ ならば、対応する射影子 $e_1, e_2$ は
$\Hom(Y, Y)$ の冪等元へ写される。$Y = Y_1 \oplus Y_2$ を、対応する
$Y$ の直和分解（節 \ref{homology-section-karoubian}）とすると、全単射
$\Hom(Z' \otimes X, Z) = \Hom(Z', Z \otimes Y)$ のもとで、
$i = 1, 2$ に対し部分群として関手的に
$\Hom(Z' \otimes X_i, Z) = \Hom(Z', Z \otimes Y_i)$ が成り立つ。
したがって、圏論、注意 \ref{categories-remark-left-dual-adjoint} の議論により、
$Y_i$ は $X_i$ の左双対である。
\end{proof}

\begin{example}
\label{homology-example-graded-vector-spaces}
$F$ を体とする。$\mathcal{C}$ を次数付き $F$-ベクトル空間の圏とする。
次数付きベクトル空間 $V$ と $W$ に対し、$V \otimes W$ を、その次数 $n$ の部分が
$$
(V \otimes W)^n = \bigoplus\nolimits_{n = p + q} V^p \otimes_F W^q
$$
である次数付き $F$-ベクトル空間とする。第三の次数付きベクトル空間 $U$ に対し、
結合制約 $\phi : U \otimes (V \otimes W) \to (U \otimes V) \otimes W$ として、
ベクトル空間の「通常の」同型
$$
U^p \otimes_F (V^q \otimes_F W^r) \to (U^p \otimes_F V^q) \otimes_F W^r
$$
を用いる。単位対象には、次数 $0$ では $F$、他の次数では零である次数付き
$F$-ベクトル空間 $\mathbf{1}$ を用いる。$\mathcal{C}$ 上には、$\mathcal{C}$ を
対称モノイド圏にする可換制約が二つあり、一方は符号を伴い、他方は伴わない。
通常は符号を伴う方を
用いる。具体的には、$V$ と $W$ が次数付き $F$-ベクトル空間ならば、次数 $n$ で
$$
V^p \otimes_F W^q \to W^q \otimes_F V^p,\quad
v \otimes w \mapsto (-1)^{pq} w \otimes v
$$
を用いる同型 $\psi : V \otimes W \to W \otimes V$ を採用する。
これが機能することの確認は省略する。
\end{example}

\begin{lemma}
\label{homology-lemma-left-dual-graded-vector-spaces}
$F$ を体とする。$\mathcal{C}$ を、例 \ref{homology-example-graded-vector-spaces} のように
モノイド圏とみなした次数付き $F$-ベクトル空間の圏とする。
$\mathcal{C}$ の $V$ が左双対 $W$ をもつならば、
$\sum_n \dim_F V^n < \infty$ であり、写像 $\epsilon$ は非退化なペアリング
$W^{-n} \times V^n \to F$ を定める。
\end{lemma}

\begin{proof}
単位対象としてとる。% P02-E0009: source line 4074 is the incomplete fragment ``As unit we take''.
圏論、定義 \ref{categories-definition-dual} により、写像
$$
\eta : \mathbf{1} \to V \otimes W\quad
\epsilon : W \otimes V \to \mathbf{1}
$$
をもつ。$\mathbf{1} = F$ は次数 $0$ に置かれているので、$\epsilon$ を補題の
主張にある対 $W^{-n} \times V^n \to F$ の列とみなせる。
すべての $n$ に対して $V^n$ の基底 $\{e_{n, i}\}_{i \in I_n}$ を選ぶ。
ほとんどすべてが零である元 $w_{-n, i} \in W^{-n}$ を用いて
$$
\eta(1) = \sum e_{n, i} \otimes w_{-n, i}
$$
と書く！ $(\epsilon \otimes 1) \circ (1 \otimes \eta)$ が $W$ 上の恒等写像で
あるという条件は、
$$
\sum\nolimits_{n, i} \epsilon(w, e_{n, i})w_{-n, i} = w
$$
を意味する。したがって、$W$ は、有限個の非零ベクトル $w_{-n, i}$ によって
次数付きベクトル空間として生成される。
$(1 \otimes \epsilon) \circ (\eta \otimes 1)$ が $V$ の恒等写像であるという条件は、
$$
\sum\nolimits_{n, i} e_{n, i}\ \epsilon(w_{-n, i}, v) = v
$$
を意味する。特に $v = e_{n, i}$ と置くと、
$\epsilon(w_{-n, i}, e_{n, i'}) = \delta_{ii'}$ を得る。
ゆえに補題の主張が成り立ち、$\{w_{-n, i}\}_{i \in I_n}$ は、選んだ $V^n$ の
基底に双対な $W^{-n}$ の基底である。
\end{proof}

\section{二重複体と付随する全複体}
\label{homology-section-double-complexes}

\noindent
二重複体とそれに付随する全複体について述べる。

\begin{definition}
\label{homology-definition-double-complex}
$\mathcal{A}$ を加法圏とする。$\mathcal{A}$ における{\it 二重複体}とは、
系 $(\{A^{p, q}, d_1^{p, q}, d_2^{p, q}\}_{p, q\in \mathbf{Z}})$ であって、
各 $A^{p, q}$ が $\mathcal{A}$ の対象であり、
$d_1^{p, q} : A^{p, q} \to A^{p + 1, q}$ および
$d_2^{p, q} : A^{p, q} \to A^{p, q + 1}$ が $\mathcal{A}$ の射で、
次の規則を満たすものをいう。
\begin{enumerate}
\item $d_1^{p + 1, q} \circ d_1^{p, q} = 0$
\item $d_2^{p, q + 1} \circ d_2^{p, q} = 0$
\item $d_1^{p, q + 1} \circ d_2^{p, q} = d_2^{p + 1, q} \circ d_1^{p, q}$
\end{enumerate}
ここで、これらはすべての $p, q \in \mathbf{Z}$ に対して成り立つ。
\end{definition}

\noindent
これは単に定義の余鎖版である。すなわち、各 $A^{p, \bullet}$ は余鎖複体であり、
各 $d_1^{p, \bullet}$ は複体の射 $A^{p, \bullet} \to A^{p + 1, \bullet}$ であって、
複体の射として $d_1^{p + 1, \bullet} \circ d_1^{p, \bullet} = 0$ を満たす。
言い換えれば、二重複体は複体からなる複体とみなせる。したがって図式
$$
\xymatrix{
\ldots &
\ldots &
\ldots &
\ldots \\
\ldots \ar[r] &
A^{p, q + 1} \ar[r]^{d_1^{p, q + 1}} \ar[u] &
A^{p + 1, q + 1} \ar[r] \ar[u] &
\ldots \\
\ldots \ar[r] &
A^{p, q} \ar[r]^{d_1^{p, q}} \ar[u]^{d_2^{p, q}} &
A^{p + 1, q} \ar[r] \ar[u]_{d_2^{p + 1, q}} &
\ldots \\
\ldots &
\ldots \ar[u] &
\ldots \ar[u] &
\ldots
}
$$
の各正方形は可換である。注意：文献では、「双複体」または「二重複体」が
この図式の各正方形を反可換にするという別の定義にも出会う。

\begin{example}
\label{homology-example-double-complex-as-tensor-product-of}
$\mathcal{A}$、$\mathcal{B}$、$\mathcal{C}$ を加法圏とする。
$$
\otimes : \mathcal{A} \times \mathcal{B} \longrightarrow \mathcal{C},
\quad
(X, Y) \longmapsto X \otimes Y
$$
が射について双線形な関手であると仮定する。$\mathcal{A} \times \mathcal{B}$ の
定義については、圏論、定義 \ref{categories-definition-product-category} を参照せよ。
$\mathcal{A}$ の複体 $X^\bullet$ と $\mathcal{B}$ の複体 $Y^\bullet$ が与えられると、
$\mathcal{C}$ における二重複体
$$
K^{\bullet, \bullet} = X^\bullet \otimes Y^\bullet
$$
を得る。ここで第一微分 $K^{p, q} \to K^{p + 1, q}$ は、射
$X^p \to X^{p + 1}$ と $Y^q$ 上の恒等射から誘導される射
$X^p \otimes Y^q \to X^{p + 1} \otimes Y^q$ である。第二微分も同様である。
\end{example}

\begin{definition}
\label{homology-definition-associated-simple-complex}
$\mathcal{A}$ を加法圏とする。$A^{\bullet, \bullet}$ を二重複体とする。
{\it 付随単複体}といい $sA^\bullet$ と表すもの、すなわち、しばしば
{\it 付随全複体}ともいい $\text{Tot}(A^{\bullet, \bullet})$ と表すものは、
$$
sA^n = \text{Tot}^n(A^{\bullet, \bullet}) =
\bigoplus\nolimits_{n = p + q} A^{p, q}
$$
（これが存在する場合）と、その微分
$$
d_{sA^\bullet}^n = d_{\text{Tot}(A^{\bullet, \bullet})}^n =
\sum\nolimits_{n = p + q} (d_1^{p, q} + (-1)^p d_2^{p, q})
$$
によって与えられる。
\end{definition}

\noindent
$\mathcal{A}$ に可算直和が存在するか、または各 $n$ に対して非零な
$A^{p, n - p}$ が高々有限個ならば、$\text{Tot}(A^{\bullet, \bullet})$ は存在する。
この定義は添字 $(p, q)$ について{\it 対称でない}ことに注意せよ。

\begin{remark}
\label{homology-remark-triple-complex}
$\mathcal{A}$ を加法圏とし、$A^{\bullet, \bullet, \bullet}$ を三重複体とする。
付随全複体は、項
$$
\text{Tot}^n(A^{\bullet, \bullet, \bullet}) =
\bigoplus\nolimits_{p + q + r = n} A^{p, q, r}
$$
と微分
$$
d^n_{\text{Tot}(A^{\bullet, \bullet, \bullet})} =
\sum\nolimits_{p + q + r = n}
d_1^{p, q, r} + (-1)^pd_2^{p, q, r} + (-1)^{p + q}d_3^{p, q, r}
$$
をもつ複体である。この定義により、簡単な計算から、付随全複体は
$$
\text{Tot}(A^{\bullet, \bullet, \bullet}) =
\text{Tot}(\text{Tot}_{12}(A^{\bullet, \bullet, \bullet})) =
\text{Tot}(\text{Tot}_{23}(A^{\bullet, \bullet, \bullet}))
$$
に等しいことが分かる。言い換えれば、最初に第一変数と第二変数をまとめ、
それらの和を最後の変数とまとめてもよいし、最初に最後の二変数をまとめ、
第一変数を最後の二変数の和とまとめてもよい。
\end{remark}

\begin{remark}
\label{homology-remark-shift-double-complex}
$\mathcal{A}$ を加法圏とする。$A^{\bullet, \bullet}$ を、微分 $d_1^{p, q}$ と
$d_2^{p, q}$ をもつ二重複体とする。$A^{\bullet, \bullet}[a, b]$ によって、
$$
(A^{\bullet, \bullet}[a, b])^{p, q} = A^{p + a, q + b}
$$
および微分
$$
d_{A^{\bullet, \bullet}[a, b], 1}^{p, q} = (-1)^a d_1^{p + a, q + b}
\quad\text{and}\quad
d_{A^{\bullet, \bullet}[a, b], 2}^{p, q} = (-1)^b d_2^{p + a, q + b}
$$
をもつ二重複体を表す。この状況では、次数 $n$ において写像
$$
\xymatrix{
(\text{Tot}(A^{\bullet, \bullet})[a + b])^n =
\bigoplus_{p + q = n + a + b} A^{p, q}
\ar[d]^{\epsilon(p, q, a, b)\text{id}_{A^{p, q}}} \\
\text{Tot}(A^{\bullet, \bullet}[a, b])^n =
\bigoplus_{p' + q' = n} A^{p' + a, q' + b}
}
$$
によって与えられる、よく定義された同型
$$
\gamma :
\text{Tot}(A^{\bullet, \bullet})[a + b]
\longrightarrow
\text{Tot}(A^{\bullet, \bullet}[a, b])
$$
が存在する。ここで $\epsilon(p, q, a, b)$ はある符号である。もちろん、
$p = p' + a$ かつ $q = q' + b$ のとき、直和因子 $A^{p, q}$ は直和因子
$A^{p' + a, q' + b}$ へ写る。これらの符号に対する条件を求めるには、
始域では
$$
d|_{A^{p, q}} = (-1)^{a + b}\left(d_1^{p, q} + (-1)^pd_2^{p, q}\right)
$$
である一方、終域では
$$
d|_{A^{p' + a, q' + b}} =
(-1)^ad_1^{p' + a, q' + b} + (-1)^{p'}(-1)^bd_2^{p' + a, q' + b}
$$
であることに注意する。したがって条件は
$$
(-1)^a \epsilon(p, q, a, b) = \epsilon(p + 1, q, a, b)(-1)^{a + b}
\Leftrightarrow
\epsilon(p + 1, q, a, b) = (-1)^b \epsilon(p, q, a, b)
$$
および
$$
(-1)^{p' + b}\epsilon(p, q, a, b) =
\epsilon(p, q + 1, a, b) (-1)^{a + b + p}
\Leftrightarrow
\epsilon(p, q, a, b) = \epsilon(p, q + 1, a, b)
$$
である。ゆえに $\epsilon(p, q, a, b) = (-1)^{pb}$ と選ぶ。
\end{remark}

\begin{remark}
\label{homology-remark-double-complex-complex-of-complexes-first}
$\mathcal{A}$ を可算直和をもつ加法圏とする。二重複体の圏を
$\text{DoubleComp}(\mathcal{A})$ と表す。
$\text{DoubleComp}(\mathcal{A})$ の対象 $A^{\bullet, \bullet}$ は、次のように
複体からなる複体とみなせる。
$$
\ldots \to A^{\bullet, -1} \to A^{\bullet, 0} \to A^{\bullet, 1} \to \ldots
$$
添字の役割を入れ替えた変種については、注意
\ref{homology-remark-double-complex-complex-of-complexes-second} を参照せよ。
この注意では、付随全複体をとる操作が、本章でここまでに調べた複体上の
すべての構造と両立することを示す。

\medskip\noindent
第一に、上のように複体からなる複体とみなした二重複体上のシフト関手は、
注意 \ref{homology-remark-shift-double-complex} で定義した関手 $[0, 1]$ である。
注意 \ref{homology-remark-shift-double-complex} により、関手
$$
\text{Tot} : \text{DoubleComp}(\mathcal{A}) \to \text{Comp}(\mathcal{A})
$$
はシフト関手と両立する。すなわち、関手的同型
$\gamma : \text{Tot}(A^{\bullet, \bullet})[1] \to
\text{Tot}(A^{\bullet, \bullet}[0, 1])$ が存在する。

\medskip\noindent
第二に、
$$
f, g : A^{\bullet, \bullet} \to B^{\bullet, \bullet}
$$
について、$f$ と $g$ を上のように複体からなる複体の射とみなしたときに
これらがホモトピックならば、
$$
\text{Tot}(f), \text{Tot}(g) :
\text{Tot}(A^{\bullet, \bullet}) \to \text{Tot}(B^{\bullet, \bullet})
$$
はホモトピックな複体写像である。実際、$h = (h^q)$ を $f$ と $g$ の間の
ホモトピーとする。$h^q$ の次数 $p$ の成分を
$h^{p, q} : A^{p, q} \to B^{p, q - 1}$ と表すと、これは
$$
f^{p, q} - g^{p, q} = d_2^{p, q - 1} \circ h^{p, q} +
h^{p, q + 1} \circ d_2^{p, q}
$$
を意味する。$h^q : A^{\bullet, q} \to B^{\bullet, q - 1}$ が複体の写像で
あるという事実は、
$$
d_1^{p, q - 1} \circ h^{p, q} = h^{p + 1, q} \circ d_1^{p, q}
$$
を意味する。$h' = ((h')^n)$ を、写像
$(h')^n : \text{Tot}^n(A^{\bullet, \bullet}) \to
\text{Tot}^{n - 1}(B^{\bullet, \bullet})$ が $p + q = n$ の直和因子
$A^{p, q}$ 上で $(-1)^ph^{p, q}$ を用いるようなホモトピーとして定義する。
すると、
$$
d_{\text{Tot}(B^{\bullet, \bullet})} \circ h' +
h' \circ d_{\text{Tot}(A^{\bullet, \bullet})}
$$
を直和因子 $A^{p, q}$ に制限したものは
$$
d_1^{p, q - 1} \circ (-1)^p h^{p, q} +
(-1)^p d_2^{p, q - 1} \circ (-1)^p h^{p, q} +
(-1)^{p + 1} h^{p + 1, q} \circ d_1^{p, q} +
(-1)^p h^{p, q + 1} \circ (-1)^p d_2^{p, q}
$$
に等しく、上の等式によりこれは $f^{p, q} - g^{p, q}$ となる。
これで第二の両立性が証明された。

\medskip\noindent
第三に、前段で $f = g$ であると仮定する。このとき、上の対応
$h \leadsto h'$ は補題 \ref{homology-lemma-homotopy-shift-cochain} の同一視と両立する。
より正確には、この補題を通じて $h$ を複体からなる複体の射
$A^{\bullet, \bullet} \to B^{\bullet, \bullet}[0, -1]$ とみなすと、
$$
\text{Tot}(A^{\bullet, \bullet})
\xrightarrow{\text{Tot}(h)}
\text{Tot}(B^{\bullet, \bullet}[0, -1])
\xrightarrow{\gamma^{-1}}
\text{Tot}(B^{\bullet, \bullet})[-1]
$$
は、この補題によって複体の射とみなした $h'$ に等しい。ここで $\gamma$ は
注意 \ref{homology-remark-shift-double-complex} の同一視である。この第三の点の確認は直ちにできる。

\medskip\noindent
第四に、
$$
0 \to A^{\bullet, \bullet} \to B^{\bullet, \bullet} \to
C^{\bullet, \bullet} \to 0
$$
を二重複体からなる複体とし、二重複体を上のように複体からなる複体とみなしたとき、
補題 \ref{homology-lemma-ses-termwise-split-cochain} のような分裂
$s^q : C^{\bullet, q} \to B^{\bullet, q}$ および
$\pi^q : B^{\bullet, q} \to A^{\bullet, q}$ が与えられていると仮定する。
一方で、補題 \ref{homology-lemma-ses-termwise-split-cochain} の手順により写像
$$
\delta : C^{\bullet, \bullet} \longrightarrow A^{\bullet, \bullet}[0, 1]
$$
が得られる。他方で $\text{Tot}$ をとると、項ごとに分裂する（下を参照）複体
$$
0 \to \text{Tot}(A^{\bullet, \bullet}) \to
\text{Tot}(B^{\bullet, \bullet}) \to
\text{Tot}(C^{\bullet, \bullet}) \to 0
$$
が得られる。したがって、補題 \ref{homology-lemma-ses-termwise-split-cochain} および
\ref{homology-lemma-ses-termwise-split-homotopy-cochain} により、これはホモトピーを除いて
よく定義された射
$$
\delta' :
\text{Tot}(C^{\bullet, \bullet})
\longrightarrow
\text{Tot}(A^{\bullet, \bullet})[1]
$$
を伴う。主張は、これらの写像が次の意味で一致するということである。すなわち、
$$
\text{Tot}(C^{\bullet, \bullet})
\xrightarrow{\text{Tot}(\delta)}
\text{Tot}(A^{\bullet, \bullet}[0, 1]) \xrightarrow{\gamma^{-1}}
\text{Tot}(A^{\bullet, \bullet})[1]
$$
は $\delta'$ に等しい。ここで $\gamma$ は注意
\ref{homology-remark-shift-double-complex} のものである。これを見るため、$s^q$ と $\pi^q$ の
成分を $s^{p, q} : C^{p, q} \to B^{\bullet, q}$ および
$\pi^{p, q} : B^{p, q} \to A^{p, q}$ と表す。
% P02-E0010: source line 4435 types the component codomain as B^{\bullet,q}, not B^{p,q}.
分裂
$(s')^n : \text{Tot}^n(C^{\bullet, \bullet}) \to
\text{Tot}^n(B^{\bullet, \bullet})$ および
$(\pi')^n : \text{Tot}^n(B^{\bullet, \bullet}) \to
\text{Tot}^n(A^{\bullet, \bullet})$ として、$p + q = n$ のときに成分が
$s^{p, q}$ および $\pi^{p, q}$ である写像を用いる。ここで
$$
(\delta')^n =
(\pi')^{n + 1} \circ d_{\text{Tot}(B^{\bullet, \bullet})}^n \circ (s')^n :
\text{Tot}^n(C^{\bullet, \bullet}) \to
\text{Tot}^{n + 1}(A^{\bullet, \bullet})
$$
であった。この直和因子 $C^{p, q}$ への制限は
$$
\pi^{p + 1, q} \circ
d_1^{p, q} \circ
s^{p, q} +
\pi^{p, q + 1} \circ
(-1)^p d_2^{p, q} \circ
s^{p, q} =
\pi^{p, q + 1} \circ
(-1)^p d_2^{p, q} \circ
s^{p, q}
$$
に等しい。この等式が成り立つのは、$s^q$ が（$d_1$ を微分とする）複体の射で
あり、かつ $s$ と $\pi$ が各双次数における $B$ の直和分解に対応するため
$\pi^{p + 1, q} \circ s^{p + 1, q} = 0$ だからである。他方、$\delta$ に対しては
$$
\delta^q =  \pi^q \circ d_2 \circ s^q :
C^{\bullet, q} \to A^{\bullet, q + 1}
$$
であり、その直和因子 $C^{p, q}$ への制限は
$\pi^{p, q + 1} \circ d_2^{p, q} \circ s^{p, q}$ に等しい。
符号の差は同型 $\gamma$ における符号 $(-1)^p$ によってちょうど打ち消され、
第四の主張が証明された。
\end{remark}

\begin{remark}
\label{homology-remark-double-complex-complex-of-complexes-second}
$\mathcal{A}$ を可算直和をもつ加法圏とする。二重複体の圏を
$\text{DoubleComp}(\mathcal{A})$ と表す。
$\text{DoubleComp}(\mathcal{A})$ の対象 $A^{\bullet, \bullet}$ は、次のように
複体からなる複体とみなせる。
$$
\ldots \to A^{-1, \bullet} \to A^{0, \bullet} \to A^{1, \bullet} \to \ldots
$$
添字の役割を入れ替えた変種については、注意
\ref{homology-remark-double-complex-complex-of-complexes-first} を参照せよ。
この注意では、付随全複体をとる操作が、本章でここまでに調べた複体上の
すべての構造と両立することを示す。

\medskip\noindent
第一に、上のように複体からなる複体とみなした二重複体上のシフト関手は、
注意 \ref{homology-remark-shift-double-complex} で定義した関手 $[1, 0]$ である。
注意 \ref{homology-remark-shift-double-complex} により、関手
$$
\text{Tot} : \text{DoubleComp}(\mathcal{A}) \to \text{Comp}(\mathcal{A})
$$
はシフト関手と両立する。すなわち、関手的同型
$\gamma : \text{Tot}(A^{\bullet, \bullet})[1] \to
\text{Tot}(A^{\bullet, \bullet}[1, 0])$ が存在する。

\medskip\noindent
第二に、
$$
f, g : A^{\bullet, \bullet} \to B^{\bullet, \bullet}
$$
について、$f$ と $g$ を上のように複体からなる複体の射とみなしたときに
これらがホモトピックならば、
$$
\text{Tot}(f), \text{Tot}(g) :
\text{Tot}(A^{\bullet, \bullet}) \to \text{Tot}(B^{\bullet, \bullet})
$$
はホモトピックな複体写像である。実際、$h = (h^p)$ を $f$ と $g$ の間の
ホモトピーとする。$h^q$ の次数 $p$ の成分を
$h^{p, q} : A^{p, q} \to B^{p - 1, q}$ と表すと、これは
% P02-E0011: source line 4518 says degree p of h^q although the family is h=(h^p).
$$
f^{p, q} - g^{p, q} = d_1^{p - 1, q} \circ h^{p, q} +
h^{p + 1, q} \circ d_1^{p, q}
$$
を意味する。$h^p : A^{p, \bullet} \to B^{p - 1, \bullet}$ が複体の写像で
あるという事実は、
$$
d_2^{p - 1, q} \circ h^{p, q} = h^{p, q + 1} \circ d_2^{p, q}
$$
を意味する。$h' = ((h')^n)$ を、写像
$(h')^n : \text{Tot}^n(A^{\bullet, \bullet}) \to
\text{Tot}^{n - 1}(B^{\bullet, \bullet})$ が $p + q = n$ の直和因子
$A^{p, q}$ 上で $h^{p, q}$ を用いるようなホモトピーとして定義する。
すると、
$$
d_{\text{Tot}(B^{\bullet, \bullet})} \circ h' +
h' \circ d_{\text{Tot}(A^{\bullet, \bullet})}
$$
を直和因子 $A^{p, q}$ に制限したものは
$$
d_1^{p - 1, q} \circ h^{p, q} +
(-1)^{p - 1} d_2^{p - 1, q} \circ h^{p, q} +
h^{p + 1, q} \circ d_1^{p, q} +
h^{p, q + 1} \circ (-1)^p d_2^{p, q}
$$
に等しく、上の等式によりこれは $f^{p, q} - g^{p, q}$ となる。
これで第二の両立性が証明された。

\medskip\noindent
第三に、前段で $f = g$ であると仮定する。このとき、上の対応
$h \leadsto h'$ は補題 \ref{homology-lemma-homotopy-shift-cochain} の同一視と両立する。
より正確には、この補題を通じて $h$ を複体からなる複体の射
$A^{\bullet, \bullet} \to B^{\bullet, \bullet}[-1, 0]$ とみなすと、
$$
\text{Tot}(A^{\bullet, \bullet})
\xrightarrow{\text{Tot}(h)}
\text{Tot}(B^{\bullet, \bullet}[-1, 0])
\xrightarrow{\gamma^{-1}}
\text{Tot}(B^{\bullet, \bullet})[-1]
$$
は、この補題によって複体の射とみなした $h'$ に等しい。ここで $\gamma$ は
注意 \ref{homology-remark-shift-double-complex} の同一視である。この第三の点の確認は直ちにできる。

\medskip\noindent
第四に、
$$
0 \to A^{\bullet, \bullet} \to B^{\bullet, \bullet} \to
C^{\bullet, \bullet} \to 0
$$
を二重複体からなる複体とし、二重複体を上のように複体からなる複体とみなしたとき、
補題 \ref{homology-lemma-ses-termwise-split-cochain} のような分裂
$s^p : C^{p, \bullet} \to B^{p, \bullet}$ および
$\pi^p : B^{p, \bullet} \to A^{p, \bullet}$ が与えられていると仮定する。
一方で、補題 \ref{homology-lemma-ses-termwise-split-cochain} の手順により写像
$$
\delta : C^{\bullet, \bullet} \longrightarrow A^{\bullet, \bullet}[0, 1]
$$
% P02-E0012: source line 4581 uses [0,1] although the second-index-switched construction and lines 4603/4648 require [1,0].
が得られる。他方で $\text{Tot}$ をとると、項ごとに分裂する（下を参照）複体
$$
0 \to \text{Tot}(A^{\bullet, \bullet}) \to
\text{Tot}(B^{\bullet, \bullet}) \to
\text{Tot}(C^{\bullet, \bullet}) \to 0
$$
が得られる。したがって、補題 \ref{homology-lemma-ses-termwise-split-cochain} および
\ref{homology-lemma-ses-termwise-split-homotopy-cochain} により、これはホモトピーを除いて
よく定義された射
$$
\delta' :
\text{Tot}(C^{\bullet, \bullet})
\longrightarrow
\text{Tot}(A^{\bullet, \bullet})[1]
$$
を伴う。主張は、これらの写像が次の意味で一致するということである。すなわち、
$$
\text{Tot}(C^{\bullet, \bullet})
\xrightarrow{\text{Tot}(\delta)}
\text{Tot}(A^{\bullet, \bullet}[1, 0]) \xrightarrow{\gamma^{-1}}
\text{Tot}(A^{\bullet, \bullet})[1]
$$
は $\delta'$ に等しい。ここで $\gamma$ は注意
\ref{homology-remark-shift-double-complex} のものである。これを見るため、$s^q$ と $\pi^q$ の
成分を $s^{p, q} : C^{p, q} \to B^{\bullet, q}$ および
$\pi^{p, q} : B^{p, q} \to A^{p, q}$ と表す。
% P02-E0013: source lines 4608-4609 retain s^q/pi^q and B^{\bullet,q} although this construction uses s^p/pi^p and bidegree components.
分裂
$(s')^n : \text{Tot}^n(C^{\bullet, \bullet}) \to
\text{Tot}^n(B^{\bullet, \bullet})$ および
$(\pi')^n : \text{Tot}^n(B^{\bullet, \bullet}) \to
\text{Tot}^n(A^{\bullet, \bullet})$ として、$p + q = n$ のときに成分が
$s^{p, q}$ および $\pi^{p, q}$ である写像を用いる。ここで
$$
(\delta')^n =
(\pi')^{n + 1} \circ d_{\text{Tot}(B^{\bullet, \bullet})}^n \circ (s')^n :
\text{Tot}^n(C^{\bullet, \bullet}) \to
\text{Tot}^{n + 1}(A^{\bullet, \bullet})
$$
であった。この直和因子 $C^{p, q}$ への制限は
$$
\pi^{p + 1, q} \circ
d_1^{p, q} \circ
s^{p, q} +
\pi^{p, q + 1} \circ
(-1)^p d_2^{p, q} \circ
s^{p, q} =
\pi^{p + 1, q} \circ
d_1^{p, q} \circ
s^{p, q}
$$
に等しい。この等式が成り立つのは、$s^p$ が（$d_2$ を微分とする）複体の射で
あり、かつ $s$ と $\pi$ が各双次数における $B$ の直和分解に対応するため
$\pi^{p, q + 1} \circ s^{p, q + 1} = 0$ だからである。他方、$\delta$ に対しては
$$
\delta^p =  \pi^p \circ d_1 \circ s^p :
C^{p, \bullet} \to A^{p + 1, \bullet}
$$
であり、その直和因子 $C^{p, q}$ への制限は
$\pi^{p + 1, q} \circ d_1^{p, q} \circ s^{p, q}$ に等しい。
したがって前と同じものが得られ、これは同型
$\gamma : \text{Tot}(A^{\bullet, \bullet})[1] \to
\text{Tot}(A^{\bullet, \bullet}[1, 0])$ が符号を介さずに定義されることとも一致する。
\end{remark}

\section{フィルトレーション}
\label{homology-section-filtrations}

\noindent
この内容についてのよい参考文献は \cite[Section 1]{HodgeII} である。
（アーベル圏に関する本書の規約とは異なることに注意せよ。）

\begin{definition}
\label{homology-definition-filtered}
$\mathcal{A}$ をアーベル圏とする。
\begin{enumerate}
\item 対象 $A$ 上の{\it 減少フィルトレーション} $F$ とは、$A$ の部分対象の族
$(F^nA)_{n \in \mathbf{Z}}$ であって、
$$
A \supset \ldots \supset F^nA \supset F^{n + 1}A \supset \ldots \supset 0
$$
を満たすものをいう。
\item $\mathcal{A}$ の{\it フィルトレーション付き対象}とは、$\mathcal{A}$ の対象
$A$ と $A$ 上の減少フィルトレーション $F$ からなる対 $(A, F)$ をいう。
\item {\it フィルトレーション付き対象の射 $(A, F) \to (B, F)$} とは、
すべての $i \in \mathbf{Z}$ に対して $\varphi(F^iA) \subset F^iB$ を満たす
$\mathcal{A}$ の射 $\varphi : A \to B$ によって与えられるものをいう。
\item フィルトレーション付き対象の圏を $\text{Fil}(\mathcal{A})$ と表す。
\item フィルトレーション付き対象 $(A, F)$ と部分対象 $X \subset A$ が与えられたとき、
$X$ 上の{\it 誘導フィルトレーション}とは、$F^nX = X \cap F^nA$ をもつ
フィルトレーションをいう。
\item フィルトレーション付き対象 $(A, F)$ と全射 $\pi : A \to Y$ が与えられたとき、
{\it 商フィルトレーション}とは、$F^nY = \pi(F^nA)$ をもつフィルトレーションをいう。
\item 対象 $A$ 上のフィルトレーション $F$ が{\it 有限}であるとは、
$F^nA = A$ および $F^mA = 0$ を満たす $n, m$ が存在することをいう。
\item フィルトレーション付き対象 $(A, F)$ に対し、すべての $F^iA$ に含まれる
$A$ の最大の部分対象が存在するとき、$\bigcap F^iA$ が存在するという。
すべての $F^iA$ を含む $A$ の最小の部分対象が存在するとき、
$\bigcup F^iA$ が存在するという。
\item フィルトレーション付き対象 $(A, F)$ 上のフィルトレーションは、
$\bigcap F^iA = 0$ ならば{\it 分離的}、$\bigcup F^iA = A$ ならば
{\it 網羅的}であるという。
\end{enumerate}
\end{definition}

\noindent
記号の濫用により、フィルトレーション付き対象の射
$f : (A, F) \to (B, F)$ が{\it 単射}であるとは、アーベル圏 $\mathcal{A}$ において
$f : A \to B$ が単射であることをいう。同様に、圏 $\mathcal{A}$ において
$f : A \to B$ が全射ならば $f$ は{\it 全射}であるという。単射（それぞれ全射）
であることは、$\text{Fil}(\mathcal{A})$ におけるモノ射（それぞれエピ射）である
ことと同値である。補題 \ref{homology-lemma-filtered} により、これは核（それぞれ余核）が
零であることとも同値である。

\begin{lemma}
\label{homology-lemma-filtered}
$\mathcal{A}$ をアーベル圏とする。フィルトレーション付き対象の圏
$\text{Fil}(\mathcal{A})$ は次の性質をもつ。
\begin{enumerate}
\item 加法圏である。
\item 零対象をもつ。
\item 核と余核、像と余像をもつ。
\item 一般にはアーベル圏ではない。
\end{enumerate}
\end{lemma}

\begin{proof}
$\text{Fil}(\mathcal{A})$ が加法的であり、その直和が
$(A, F) \oplus (B, F) = (A \oplus B, F)$、ここで
$F^p(A \oplus B) = F^pA \oplus F^pB$、によって与えられることは明らかである。
フィルトレーション付き対象の射 $f : (A, F) \to (B, F)$ の核は単射
$\Ker(f) \subset A$ であり、ここで $\Ker(f)$ には誘導フィルトレーションを入れる。
フィルトレーション付き対象の射 $f : A \to B$ の余核は全射
$B \to \Coker(f)$ であり、ここで $\Coker(f)$ には商フィルトレーションを入れる。
すべての核と余核が存在するので、すべての余像と像も存在する。最後の主張については
例 \ref{homology-example-not-abelian} を参照せよ。
\end{proof}

\begin{definition}
\label{homology-definition-strict}
$\mathcal{A}$ をアーベル圏とする。$\mathcal{A}$ のフィルトレーション付き対象の射
$f : A \to B$ が{\it 厳密}であるとは、すべての $i \in \mathbf{Z}$ に対して
$f(F^iA) = f(A) \cap F^iB$ が成り立つことをいう。
\end{definition}

\noindent
これは、すべての $i \in \mathbf{Z}$ に対して
$f^{-1}(F^iB) = F^iA + \Ker(f)$ を要求することとも同値である。
厳密射を次のように特徴付ける。

\begin{lemma}
\label{homology-lemma-characterize-strict-general}
$\mathcal{A}$ をアーベル圏とする。$f : A \to B$ を $\mathcal{A}$ の
フィルトレーション付き対象の射とする。次は同値である。
\begin{enumerate}
\item $f$ は厳密である。
\item 補題 \ref{homology-lemma-coim-im-map} の射 $\Coim(f) \to \Im(f)$ は同型である。
\end{enumerate}
\end{lemma}

\begin{proof}
$\Coim(f) \to \Im(f)$ は $\mathcal{A}$ の対象の同型であり、(2) はそれが
フィルトレーション付き対象の同型であることを意味する。補題
\ref{homology-lemma-filtered} の証明における核と余核の記述から、$\Coim(f)$ 上の
フィルトレーションは $A \to \Coim(f)$ からくる商フィルトレーションである。
同様に、$\Im(f)$ 上のフィルトレーションは単射 $\Im(f) \to B$ からくる
誘導フィルトレーションである。厳密性の定義は、商フィルトレーションが
誘導フィルトレーションであることにほかならない。
\end{proof}

\begin{lemma}
\label{homology-lemma-add-summand-strict-monomorphism}
$\mathcal{A}$ をアーベル圏とする。$f : A \to B$ をフィルトレーション付き対象の
厳密なモノ射とし、$g : A \to C$ をフィルトレーション付き対象の射とする。
このとき $f \oplus g : A \to B \oplus C$ は厳密なモノ射である。
\end{lemma}

\begin{proof}
定義から明らかである。
\end{proof}

\begin{lemma}
\label{homology-lemma-add-summand-strict-epimorphism}
$\mathcal{A}$ をアーベル圏とする。$f : B \to A$ をフィルトレーション付き対象の
厳密なエピ射とし、$g : C \to A$ をフィルトレーション付き対象の射とする。
このとき $f \oplus g : B \oplus C \to A$ は厳密なエピ射である。
\end{lemma}

\begin{proof}
定義から明らかである。
\end{proof}

\begin{lemma}
\label{homology-lemma-induced-and-quotient-strict}
$\mathcal{A}$ をアーベル圏とする。$(A, F)$、$(B, F)$ をフィルトレーション付き対象とし、
$u : A \to B$ をフィルトレーション付き対象の射とする。$u$ が単射ならば、
$u$ が厳密であるための必要十分条件は、$A$ 上のフィルトレーションが
誘導フィルトレーションであることである。$u$ が全射ならば、$u$ が厳密であるための
必要十分条件は、$B$ 上のフィルトレーションが商フィルトレーションであることである。
\end{lemma}

\begin{proof}
これは定義から直ちに従う。
\end{proof}

\begin{lemma}
\label{homology-lemma-composition-strict}
$\mathcal{A}$ をアーベル圏とする。$f : A \to B$、$g : B \to C$ を
フィルトレーション付き対象の厳密射とする。
\begin{enumerate}
\item 一般には合成 $g \circ f$ は厳密でない。
\item $g$ が単射ならば、$g \circ f$ は厳密である。
\item $f$ が全射ならば、$g \circ f$ は厳密である。
\end{enumerate}
\end{lemma}

\begin{proof}
$B$ を、基底 $e_1, e_2$ をもつ体 $k$ 上のベクトル空間とし、$n < 0$ では
$F^nB = B$、$F^0B = ke_1$、$n > 0$ では $F^nB = 0$ となるフィルトレーションを
入れる。ここで $A = k(e_1 + e_2)$ および $C = B/ke_2$ とし、$B$ から誘導される
フィルトレーション、すなわち $A \to B$ と $B \to C$ が厳密になるフィルトレーションを
入れる（補題 \ref{homology-lemma-induced-and-quotient-strict}）。すると、$n < 0$ では
$F^n(A) = A$、$n \geq 0$ では $F^n(A) = 0$ である。また、$n \leq 0$ では
$F^n(C) = C$、$n > 0$ では $F^n(C) = 0$ である。ゆえに非零な合成
$A \to C$ は厳密でない。

\medskip\noindent
$g$ が単射であると仮定する。このとき
\begin{align*}
g(f(F^pA)) & = g(f(A) \cap F^pB) \\
& = g(f(A)) \cap g(F^p(B)) \\
& = (g \circ f)(A) \cap (g(B) \cap F^pC) \\
& = (g \circ f)(A) \cap F^pC.
\end{align*}
第一の等号は $f$ が厳密だから、第二の等号は $g$ が単射だから、第三の等号は
$g$ が厳密だから、第四の等号は $(g \circ f)(A) \subset g(B)$ だから成り立つ。

\medskip\noindent
$f$ が全射であると仮定する。このとき
\begin{align*}
(g \circ f)^{-1}(F^iC) & = f^{-1}(F^iB + \Ker(g)) \\
& = f^{-1}(F^iB) + f^{-1}(\Ker(g)) \\
& = F^iA + \Ker(f) + \Ker(g \circ f) \\
& = F^iA + \Ker(g \circ f)
\end{align*}
第一の等号は $g$ が厳密だから、第二の等号は $f$ が全射だから、第三の等号は
$f$ が厳密だから、最後の等号は $\Ker(f) \subset \Ker(g \circ f)$ だから成り立つ。
\end{proof}

\noindent
次の補題は、フィルトレーション付き対象の部分対象には、その対象を余核として
書く仕方の選択に依存しない、よく定義されたフィルトレーションが入ることを述べる。

\begin{lemma}
\label{homology-lemma-filtration-subobject}
$\mathcal{A}$ をアーベル圏とする。$(A, F)$ を $\mathcal{A}$ の
フィルトレーション付き対象とし、$X \subset Y \subset A$ を $A$ の部分対象とする。
対象
$$
Y/X = \Ker(A/X \to A/Y)
$$
上では、$Y$ 上の誘導フィルトレーションからくる商フィルトレーションと、
$A/X$ 上の商フィルトレーションからくる誘導フィルトレーションとが一致する。
射 $X \to Y$、$X \to A$、$Y \to A$、$Y \to A/X$、$Y \to Y/X$、
$Y/X \to A/X$ はいずれも（誘導／商フィルトレーションに関して）厳密である。
\end{lemma}

\begin{proof}
$Y/X$ 上の商フィルトレーションは
$F^p(Y/X) = F^pY/(X \cap F^pY) = F^pY/F^pX$ によって与えられる。
これは $F^pY = Y \cap F^pA$ および $F^pX = X \cap F^pA$ だからである。
単射 $Y/X \to A/X$ からくる誘導フィルトレーションは
\begin{align*}
F^p(Y/X) & = Y/X \cap F^p(A/X) \\
& = Y/X \cap (F^pA + X)/X \\
& = (Y \cap F^pA)/(X \cap F^pA) \\
& = F^pY/F^pX.
\end{align*}
によって与えられる。これで補題の第一の主張を得る。他の場合の証明も同様である。
\end{proof}

\begin{lemma}
\label{homology-lemma-pushout-filtered}
$\mathcal{A}$ をアーベル圏とする。$A, B, C \in \text{Fil}(\mathcal{A})$ とし、
$f : A \to B$ および $g : A \to C$ を射とする。このとき
$\text{Fil}(\mathcal{A})$ に押し出し
$$
\xymatrix{
A \ar[r]_f \ar[d]_g & B \ar[d]^{g'} \\
C \ar[r]^{f'} & C \amalg_A B
}
$$
が存在する。$f$ が厳密ならば $f'$ も厳密である。
\end{lemma}

\begin{proof}
$C \amalg_A B$ を $\text{Fil}(\mathcal{A})$ における
$\Coker((g, -f) : A \to C \oplus B)$ と置く。この余核は補題
\ref{homology-lemma-filtered} により存在する。これは押し出しである。例
\ref{homology-example-fibre-product-pushouts} を参照せよ。$F^p(C \amalg_A B)$ は
$F^pC \oplus F^pB$ の像であることに注意する。したがって
$$
(f')^{-1}(F^p(C \amalg_A B)) = g(f^{-1}(F^pB)) + F^pC
$$
であり、最後の主張が従う。
\end{proof}

\begin{lemma}
\label{homology-lemma-fibre-product-filtered}
$\mathcal{A}$ をアーベル圏とする。$A, B, C \in \text{Fil}(\mathcal{A})$ とし、
$f : B \to A$ および $g : C \to A$ を射とする。このとき
$\text{Fil}(\mathcal{A})$ にファイバー積
$$
\xymatrix{
B \times_A C \ar[r]_{g'} \ar[d]_{f'} & B \ar[d]^f \\
C \ar[r]^g & A
}
$$
が存在する。$f$ が厳密ならば $f'$ も厳密である。
\end{lemma}

\begin{proof}
この補題は補題 \ref{homology-lemma-pushout-filtered} の双対である。
\end{proof}

\noindent
$\mathcal{A}$ をアーベル圏とし、$(A, F)$ を $\mathcal{A}$ の
フィルトレーション付き対象とする。$\mathcal{A}$ の対象 $F^pA/F^{p + 1}A$ を
$\text{gr}^p_F(A) = \text{gr}^p(A)$ と表す。これにより加法関手
$$
\text{gr}^p :
\text{Fil}(\mathcal{A})
\longrightarrow
\mathcal{A}, \quad
(A, F)
\longmapsto
\text{gr}^p(A).
$$
が定まる。節 \ref{homology-section-graded} で $\mathcal{A}$ の次数付き対象の圏
$\text{Gr}(\mathcal{A})$ を定義したことを思い出そう。
$\text{Fil}(\mathcal{A})$ の $(A, F)$ に対して
$$
\text{gr}(A) = \text{the graded object of }\mathcal{A}\text{ whose }
p\text{th graded piece is }\text{gr}^p(A)
$$
と置ける。また、$\mathcal{A}$ が可算直和をもつならば、単に
$$
\text{gr}(A) = \bigoplus \text{gr}^p(A)
$$
である。これにより加法関手
$$
\text{gr} :
\text{Fil}(\mathcal{A})
\longrightarrow
\text{Gr}(\mathcal{A}), \quad
(A, F) \longmapsto \text{gr}(A).
$$
が定まる。

\begin{lemma}
\label{homology-lemma-ses-gr}
$\mathcal{A}$ をアーベル圏とする。
\begin{enumerate}
\item $A$ をフィルトレーション付き対象、$X \subset A$ とする。このとき各 $p$ に対し、
$$
0 \to \text{gr}^p(X) \to \text{gr}^p(A) \to \text{gr}^p(A/X) \to 0
$$
は完全である（$X$ には誘導フィルトレーションを、$A/X$ には商フィルトレーションを入れる）。
\item $f : A \to B$ を $\mathcal{A}$ のフィルトレーション付き対象の射とする。
このとき各 $p$ に対し、列
$$
0 \to \text{gr}^p(\Ker(f)) \to \text{gr}^p(A) \to
\text{gr}^p(\Coim(f)) \to 0
$$
および
$$
0 \to \text{gr}^p(\Im(f)) \to \text{gr}^p(B) \to
\text{gr}^p(\Coker(f)) \to 0
$$
は完全である。
\end{enumerate}
\end{lemma}

\begin{proof}
$F^{p + 1}X = X \cap F^{p + 1}A$ なので、写像
$\text{gr}^p(X) \to \text{gr}^p(A)$ は単射である。双対的に、写像
$\text{gr}^p(A) \to \text{gr}^p(A/X)$ は全射である。
$F^pA/F^{p + 1}A \to A/X + F^{p + 1}A$ の核は明らかに
$F^{p + 1}A + X \cap F^pA/F^{p + 1}A = F^pX/F^{p + 1}X$ であり、
したがって中央で完全である。(2) の二つの短完全列は (1) の短完全列の特別な場合である。
\end{proof}

\begin{lemma}
\label{homology-lemma-characterize-strict}
$\mathcal{A}$ をアーベル圏とし、$f : A \to B$ を $\mathcal{A}$ の
有限フィルトレーション付き対象の射とする。次は同値である。
\begin{enumerate}
\item $f$ は厳密である。
\item 射 $\Coim(f) \to \Im(f)$ は同型である。
\item $\text{gr}(\Coim(f)) \to \text{gr}(\Im(f))$ は同型である。
\item 列 $\text{gr}(\Ker(f)) \to \text{gr}(A) \to \text{gr}(B)$ は完全である。
\item 列 $\text{gr}(A) \to \text{gr}(B) \to
\text{gr}(\Coker(f))$ は完全である。
\item 列
$$
0 \to
\text{gr}(\Ker(f)) \to
\text{gr}(A) \to
\text{gr}(B) \to
\text{gr}(\Coker(f)) \to 0
$$
は完全である。
\end{enumerate}
\end{lemma}

\begin{proof}
(1) と (2) の同値性は補題 \ref{homology-lemma-characterize-strict-general} である。
補題 \ref{homology-lemma-ses-gr} により、(4)、(5)、(6) は (3) を含意し、(3) は
(4)、(5)、(6) を含意する。したがって (3) が (2) を含意することを示せば十分である。
そこで、有限フィルトレーション付き対象の単射かつ全射な写像
$f : A \to B$ が同型 $\text{gr}(A) \to \text{gr}(B)$ を誘導するならば、
$f$ がフィルトレーション付き対象の同型を誘導することを示す必要がある。
言い換えれば、すべての $p$ に対して $f(F^pA) = F^pB$ であることを示す。
フィルトレーションは有限なので、$p$ に関する下降帰納法で証明できる。
$f(F^{p + 1}A) = F^{p + 1}B$ と仮定する。このとき可換図式
$$
\xymatrix{
0 \ar[r] &
F^{p + 1}A \ar[r] \ar[d]^f &
F^pA \ar[r] \ar[d]^f &
\text{gr}^p(A) \ar[r] \ar[d]^{\text{gr}^p(f)} &
0 \\
0 \ar[r] &
F^{p + 1}B \ar[r] &
F^pB \ar[r] &
\text{gr}^p(B) \ar[r] &
0
}
$$
と五補題から $f(F^pA) = F^pB$ が従う。
\end{proof}

\begin{lemma}
\label{homology-lemma-filtered-complex}
$\mathcal{A}$ をアーベル圏とする。$A \to B \to C$ を $\mathcal{A}$ の
フィルトレーション付き対象からなる複体とする。
$\alpha : A \to B$ と $\beta : B \to C$ がフィルトレーション付き対象の
厳密射であると仮定する。このとき
$\text{gr}(\Ker(\beta)/\Im(\alpha)) =
\Ker(\text{gr}(\beta))/\Im(\text{gr}(\alpha)))$ である。
% P02-E0014: source line 5082 contains an unmatched closing parenthesis.
\end{lemma}

\begin{proof}
これは補題 \ref{homology-lemma-ses-gr} と、補題
\ref{homology-lemma-characterize-strict-general} により
$\Coim(\alpha) \cong \Im(\alpha)$ および
$\Coim(\beta) \cong \Im(\beta)$ であることから形式的に従う。
\end{proof}

\begin{lemma}
\label{homology-lemma-filtered-acyclic}
$\mathcal{A}$ をアーベル圏とする。$A \to B \to C$ を $\mathcal{A}$ の
フィルトレーション付き対象からなる複体とする。$A, B, C$ は有限フィルトレーションを
もち、$\text{gr}(A) \to \text{gr}(B) \to \text{gr}(C)$ が完全であると仮定する。
このとき、
\begin{enumerate}
\item 各 $p \in \mathbf{Z}$ に対して、列
$\text{gr}^p(A) \to \text{gr}^p(B) \to \text{gr}^p(C)$ は完全である。
\item 各 $p \in \mathbf{Z}$ に対して、列
$F^p(A) \to F^p(B) \to F^p(C)$ は完全である。
\item 各 $p \in \mathbf{Z}$ に対して、列
$A/F^p(A) \to B/F^p(B) \to C/F^p(C)$ は完全である。
\item 写像 $A \to B$ および $B \to C$ は厳密である。
\item $A \to B \to C$ は（$\mathcal{A}$ における列として）完全である。
\end{enumerate}
\end{lemma}

\begin{proof}
(1) は定義から直ちに従う。(3) をフィルトレーションの長さに関する帰納法で証明する。
$A$、$B$、$C$ のそれぞれが非零な次数成分を一つしかもたないならば、
$\text{gr}(A) = A$ などであるから (3) は成り立つ。$F^nA, F^nB, F^nC$ の
少なくとも一つが非零となる最大の整数を $n$ とする。誘導フィルトレーションを入れて
$A' = A/F^nA$、$B' = B/F^nB$、$C' = C/F^nC$ と置く。
$\text{gr}(A) = F^nA \oplus \text{gr}(A')$ であり、$B$ と $C$ についても同様である。
帰納法の仮定を $A' \to B' \to C'$ に適用すると、$p \geq n$ に対して
$A/F^p(A) \to B/F^p(B) \to C/F^p(C)$ が完全であることが従う。
$p = n + 1$ についても同じ結論、すなわち $A \to B \to C$ が完全であることを
示すため、可換図式
$$
\xymatrix{
0 \ar[r] & F^nA \ar[r] \ar[d] & A \ar[r] \ar[d] & A' \ar[r] \ar[d] & 0 \\
0 \ar[r] & F^nB \ar[r] \ar[d] & B \ar[r] \ar[d] & B' \ar[r] \ar[d] & 0 \\
0 \ar[r] & F^nC \ar[r] & C \ar[r] & C' \ar[r] & 0
}
$$
を用いる。その各行は $\mathcal{A}$ の対象の短完全列である。(2) の証明は双対的である。
もちろん (5) は (2) から従う。

\medskip\noindent
(4) を証明するため、与えられた射を $f : A \to B$ および $g : B \to C$ と表す。
(2) により $f(F^p(A)) = \Ker(F^p(B) \to F^p(C))$ であり、(5) により
$f(A) = \Ker(g)$ である。したがって
$f(F^p(A)) =  \Ker(F^p(B) \to F^p(C)) =
\Ker(g) \cap F^p(B) = f(A) \cap F^p(B)$ であり、$f$ が厳密であることが分かる。
$g$ が厳密であることの証明はこれと双対的である。
\end{proof}

\section{スペクトル系列}
\label{homology-section-spectral-sequence}

\noindent
スペクトル系列についてのよい解説は \cite{Eisenbud} にある。
\cite{McCleary}、\cite{Lang} なども参照せよ。

\begin{definition}
\label{homology-definition-spectral-sequence}
$\mathcal{A}$ をアーベル圏とする。
\begin{enumerate}
\item $\mathcal{A}$ における{\it スペクトル系列}とは、系
$(E_r, d_r)_{r \geq 1}$ であって、各 $E_r$ が $\mathcal{A}$ の対象、
各 $d_r : E_r \to E_r$ が $d_r \circ d_r = 0$ を満たす射であり、
$r \geq 1$ に対して $E_{r + 1} = \Ker(d_r)/\Im(d_r)$ となるものをいう。
\item {\it スペクトル系列の射}
$f : (E_r, d_r)_{r \geq 1} \to (E'_r, d'_r)_{r \geq 1}$ とは、
$f_r \circ d_r = d'_r \circ f_r$ を満たす射の族 $f_r : E_r \to E'_r$ であって、
同一視 $E_{r + 1} = \Ker(d_r)/\Im(d_r)$ および
$E'_{r + 1} = \Ker(d'_r)/\Im(d'_r)$ を通じて $f_{r + 1}$ が $f_r$ から
誘導されるものをいう。
\end{enumerate}
\end{definition}

\noindent
ときにはこの定義を少し緩め、指定された同型
$E_{r + 1} \to \Ker(d_r)/\Im(d_r)$ をもつ対象 $E_{r + 1}$ を許す。
さらに、ある $r_0 \in \mathbf{Z}$ に対して、添字 $\geq r_0$ で上の定義の性質を
満たす系 $(E_r, d_r)_{r \geq r_0}$ が与えられることもある。単純に番号を付け替えれば
いずれにせよ定義に含まれるので、これもスペクトル系列と呼ぶ。実際、文献には
$r_0 = 0$ および $r_0 = -1$ の場合が現れる。

\medskip\noindent
スペクトル系列 $(E_r, d_r)_{r \geq 1}$ が与えられたとき、
$$
0 = B_1 \subset B_2 \subset \ldots \subset B_r \subset \ldots
\subset Z_r \subset \ldots \subset Z_2 \subset Z_1 = E_1
$$
を次の簡単な手順で定義する。$B_2 = \Im(d_1)$ および $Z_2 = \Ker(d_1)$ と置く。
すると明らかに $d_2 : Z_2/B_2 \to Z_2/B_2$ である。したがって、$B_2$ を含み、
$B_3/B_2$ が $d_2$ の像となるような $E_1$ の一意な部分対象として $B_3$ を定義できる。
同様に、$B_2$ を含み、$Z_3/B_2$ が $d_2$ の核となるような $E_1$ の一意な部分対象として
$Z_3$ を定義できる。以下同様である。特に、すべての $r \geq 1$ に対して
$$
E_r = Z_r/B_r
$$
である。スペクトル系列が $r = r_0$ から始まる場合も、$B_i$、$Z_i$ を
$E_{r_0}$ の部分対象として同様に構成できる。実際、文献では、$E_r$ から始めた
上のフィルトレーションを表すために記号
$$
0 = B_r(E_r) \subset B_{r + 1}(E_r) \subset B_{r + 2}(E_r) \subset \ldots
\subset Z_{r + 2}(E_r) \subset Z_{r + 1}(E_r) \subset Z_r(E_r) = E_r
$$
を用いることがある。

\begin{definition}
\label{homology-definition-limit-spectral-sequence}
$\mathcal{A}$ をアーベル圏とし、$(E_r, d_r)_{r \geq 1}$ をスペクトル系列とする。
\begin{enumerate}
\item $E_1$ の部分対象 $Z_{\infty} = \bigcap Z_r$ および
$B_{\infty} = \bigcup B_r$ が存在するならば、スペクトル系列の
{\it 極限}\footnote{この記法は普遍的に採用されているわけではない。一部の文献では、
$0 = B_1 \subset B_2 \subset \ldots \subset B_\infty \subset Z_\infty
\subset \ldots \subset Z_2 \subset Z_1 = E_1$ を満たす $E_1$ の部分対象の追加の対
$Z_\infty$ と $B_\infty$ が、スペクトル系列をなすデータの一部である！}
を対象 $E_{\infty} = Z_{\infty}/B_{\infty}$ と定義する。
\item 微分 $d_r, d_{r + 1}, \ldots$ がすべて零であるとき、スペクトル系列は
{\it $E_r$ で退化する}という。
\end{enumerate}
\end{definition}

\noindent
スペクトル系列が $E_r$ で退化するならば
$E_r = E_{r + 1} = \ldots = E_{\infty}$ である（もちろん極限は存在する）。
また、本書で出会うほとんどすべてのアーベル圏は可算和と可算共通部分をもつ。

\begin{remark}[変種]
\label{homology-remark-allow-translation-functors}
スペクトル系列の項が、例えば次数付けや二重次数付けのような付加構造をもつことが
しばしばある。これに対応するため（また、ある種の技術的問題を回避するため）、
次の概念を導入する。$\mathcal{A}$ をアーベル圏とする。
$(T_r)_{r \geq 1}$ を{\it 並進}関手または{\it シフト}関手の列とする。すなわち、
$T_r : \mathcal{A} \to \mathcal{A}$ は圏の同型である。この設定での
{\it スペクトル系列}とは、系 $(E_r, d_r)_{r \geq 1}$ であって、
各 $E_r$ が $\mathcal{A}$ の対象、各 $d_r : E_r \to T_rE_r$ が
$T_rd_r \circ d_r = 0$ を満たす射となるものをいう。したがって
$$
\xymatrix{
\ldots \ar[r] &
T_r^{-1}E_r \ar[r]^-{T_r^{-1}d_r} &
E_r \ar[r]^-{d_r} &
T_rE_r \ar[r]^{T_r d_r} &
T_r^2E_r \ar[r] & \ldots
}
$$
は複体であり、$r \geq 1$ に対して
$E_{r + 1} = \Ker(d_r)/\Im(T_r^{-1}d_r)$ である。この設定における
{\it スペクトル系列の射}の意味は明らかである。この設定でも
$$
0 = B_1 \subset B_2 \subset \ldots \subset B_r \subset \ldots
\subset Z_r \subset \ldots \subset Z_2 \subset Z_1 = E_1
$$
および、上と同様に（存在するならば）$Z_\infty$ と $B_\infty$ を定義できる。
\end{remark}

\section{スペクトル系列：完全対}
\label{homology-section-exact-couple}

\begin{definition}
\label{homology-definition-exact-couple}
$\mathcal{A}$ をアーベル圏とする。
\begin{enumerate}
\item {\it 完全対}とはデータ $(A, E, \alpha, f, g)$ であって、$A$、$E$ が
$\mathcal{A}$ の対象であり、$\alpha$、$f$、$g$ が図式
$$
\xymatrix{
A \ar[rr]_{\alpha} & & A \ar[ld]^g \\
& E \ar[lu]^f &
}
$$
のような射で、各矢印の核が直前の矢印の像であるという性質をもつものをいう。
すなわち、$\Ker(\alpha) = \Im(f)$、$\Ker(f) = \Im(g)$、
$\Ker(g) = \Im(\alpha)$ である。
\item {\it 完全対の射}
$t : (A, E, \alpha, f, g) \to (A', E', \alpha', f', g')$ とは、
$\alpha' \circ t_A = t_A \circ \alpha$、$f' \circ t_E = t_A \circ f$、
$g' \circ t_A = t_E \circ g$ を満たす射 $t_A : A \to A'$ および
$t_E : E \to E'$ によって与えられるものをいう。
\end{enumerate}
\end{definition}

\begin{lemma}
\label{homology-lemma-derived-exact-couple}
$(A, E, \alpha, f, g)$ をアーベル圏 $\mathcal{A}$ における完全対とする。
次のように置く。
\begin{enumerate}
\item $d = g \circ f : E \to E$ とし、したがって $d \circ d = 0$ とする。
\item $E' = \Ker(d)/\Im(d)$。
\item $A' = \Im(\alpha)$。
\item $\alpha$ から誘導される $\alpha' : A' \to A'$。
\item $f$ から誘導される $f' : E' \to A'$。
\item 「$g \circ \alpha^{-1}$」から誘導される $g' : A' \to E'$。
\end{enumerate}
このとき次が成り立つ。
\begin{enumerate}
\item $\Ker(d) = f^{-1}(\Ker(g)) = f^{-1}(\Im(\alpha))$。
\item $\Im(d) = g(\Im(f)) = g(\Ker(\alpha))$。
\item $(A', E', \alpha', f', g')$ は完全対である。
\end{enumerate}
\end{lemma}

\begin{proof}
省略する。
\end{proof}

\noindent
したがって、完全対 $(A, E, \alpha, f, g)$ が与えられると、
$E_1 = E$、$d_1 = d$、$E_2 = E'$、$d_2 = d' = g' \circ f'$、
$E_3 = E''$、$d_3 = d'' = g'' \circ f''$ などと置くことにより、
スペクトル系列を得ることは明らかである。

\begin{definition}
\label{homology-definition-spectral-sequence-associated-exact-couple}
$\mathcal{A}$ をアーベル圏とし、$(A, E, \alpha, f, g)$ を完全対とする。
{\it 完全対に付随するスペクトル系列}とは、
$E_1 = E$、$d_1 = d$、$E_2 = E'$、$d_2 = d' = g' \circ f'$、
$E_3 = E''$、$d_3 = d'' = g'' \circ f''$ などをもつスペクトル系列
$(E_r, d_r)_{r \geq 1}$ をいう。
\end{definition}

\begin{lemma}
\label{homology-lemma-spectral-sequence-associated-exact-couple}
$\mathcal{A}$ をアーベル圏とし、$(A, E, \alpha, f, g)$ を完全対とする。
$(E_r, d_r)_{r \geq 1}$ をこの完全対に付随するスペクトル系列とする。
この場合、
$$
0 = B_1 \subset \ldots \subset
B_{r + 1} = g(\Ker(\alpha^r))
\subset \ldots \subset
Z_{r + 1} = f^{-1}(\Im(\alpha^r))
\subset \ldots \subset Z_1 = E
$$
であり、写像 $d_{r + 1} : E_{r + 1} \to E_{r + 1}$ は次の規則で記述される。
$\mathcal{A}$ の任意の（試験）対象 $T$ と、$f \circ x = \alpha^r \circ y$ を満たす
任意の元 $x : T \to Z_{r + 1}$ および $y : T \to A$ に対して
$$
d_{r + 1} \circ \overline{x} = \overline{g \circ y}
$$
である。ここで $\overline{x} : T \to E_{r + 1}$ は誘導される射である。
\end{lemma}

\begin{proof}
省略する。
\end{proof}

\noindent
補題の状況では、$\bigcup \Ker(\alpha^r)$ と $\bigcap \Im(\alpha^r)$ が存在するならば、
明らかに
$$
B_\infty = g\left(\bigcup\nolimits_r \Ker(\alpha^r)\right)
\subset
Z_\infty = f^{-1}\left(\bigcap\nolimits_r \Im(\alpha^r)\right)
$$
である。これから極限 $E_\infty = Z_\infty / B_\infty$ が得られる。定義
\ref{homology-definition-limit-spectral-sequence} を参照せよ。

\begin{remark}[変種]
\label{homology-remark-shifted-exact-couple}
$\mathcal{A}$ をアーベル圏とする。$S, T : \mathcal{A} \to \mathcal{A}$ をシフト関手、
すなわち圏の同型とする。$A \in \Ob(\mathcal{A})$ および $n \in \mathbf{Z}$ に対し、
$n$ 回の合成を $S^nA$ および $T^nA$ と表す。この状況での{\it 完全対}とは、
データ $(A, E, \alpha, f, g)$ であって、$A$、$E$ が $\mathcal{A}$ の対象、
$\alpha : A \to T^{-1}A$、$f : E \to A$、$g : A \to SE$ が射であり、
$$
\xymatrix{
TE \ar[r]^-{Tf} &
TA \ar[r]^-{T\alpha} &
A \ar[r]^-{g} &
SE \ar[r]^{Sf} & SA
}
$$
が完全複体となるものをいう。これを次のように図示しよう。
$$
\xymatrix{
TA \ar[rrr]_{T\alpha} & & &
A \ar[ld]^g \ar[rrr]_\alpha & & &
T^{-1}A \ar[ld]^{T^{-1}g} \\
& TE \ar[lu]^{Tf} \ar@{..}[r] & SE & &
E \ar[lu]^f \ar@{..}[r] & T^{-1}SE
}
$$
$d = g \circ f : E \to SE$ と置く。このとき
$f \circ S^{-1}g = 0$ なので、
$d \circ S^{-1}d = g \circ f \circ S^{-1}g \circ S^{-1}f = 0$ である。
$E' = \Ker(d)/\Im(S^{-1}d)$、$A' = \Im(T\alpha)$ と置く。
$\alpha$ から誘導される $\alpha' : A' \to T^{-1}A'$ をとる。
$f(\Ker(d)) \subset \Ker(g) = \Im(T\alpha)$ なので、$f$ から誘導される
$f' : E' \to A'$ をとれる。最後に、
「$Tg \circ (T\alpha)^{-1}$」から誘導される $g' : A' \to TSE'$ をとる
\footnote{これは、$TSE' = \Ker(TSd)/\Im(Td)$、
$Tg(\Ker(T\alpha)) = Tg(\Im(Tf)) = \Im(T(d))$、および
$TS(d)(\Im(Tg)) = \Im(TSg \circ TSf \circ Tg) = 0$ だから可能である。}。

\medskip\noindent
上とまったく同様に、次を得る。
\begin{enumerate}
\item $\Ker(d) = f^{-1}(\Ker(g)) = f^{-1}(\Im(T\alpha))$。
\item $\Im(d) = g(\Im(f)) = g(\Ker(\alpha))$。
\item $(A', E', \alpha', f', g')$ はシフト関手 $TS$ および $T$ に関する完全対である。
\end{enumerate}
注意 \ref{homology-remark-allow-translation-functors} のようなスペクトル系列を得る。
これは $E_1 = E$、$E_2 = E'$ などをもち、すべての $r \geq 1$ に対して
$d_r : E_r \to T^{r - 1}SE_r$ である。補題
\ref{homology-lemma-spectral-sequence-associated-exact-couple} により、この状況では
$$
SB_{r + 1} =
g(\Ker(T^{-r + 1}\alpha \circ \ldots \circ T^{-1}\alpha \circ \alpha))
$$
および
$$
Z_{r + 1} = f^{-1}(\Im(T\alpha \circ T^2\alpha \circ \ldots \circ T^r\alpha))
$$
である。写像 $d_{r + 1}$ の記述は補題のものと同様である。
（このような完全対からスペクトル系列が得られることを証明するには、
これらの明示的記述を使う方が容易かもしれない。）
\end{remark}

\section{スペクトル系列：微分対象}
\label{homology-section-differential-object}

\begin{definition}
\label{homology-definition-differential-object}
$\mathcal{A}$ をアーベル圏とする。$\mathcal{A}$ の{\it 微分対象}とは、
$\mathcal{A}$ の対象 $A$ と $d \circ d = 0$ を満たす自己写像 $d$ からなる
対 $(A, d)$ をいう。{\it 微分対象の射} $(A, d) \to (B, d)$ とは、
$d \circ \alpha = \alpha \circ d$ を満たす射 $\alpha : A \to B$ によって
与えられるものをいう。
\end{definition}

\begin{lemma}
\label{homology-lemma-differential-objects-abelian}
\begin{slogan}
アーベル圏の微分対象の圏は、それ自身アーベル圏である。
\end{slogan}
$\mathcal{A}$ をアーベル圏とする。$\mathcal{A}$ の微分対象の圏はアーベル圏である。
\end{lemma}

\begin{proof}
省略する。
\end{proof}

\begin{definition}
\label{homology-definition-differential-object-homology}
微分対象 $(A, d)$ に対して、
$$
H(A, d) = \Ker(d)/\Im(d)
$$
をその{\it ホモロジー}と呼ぶ。
\end{definition}

\begin{lemma}
\label{homology-lemma-differential-objects-ses}
$\mathcal{A}$ をアーベル圏とする。
$0 \to (A, d) \to (B, d) \to (C, d) \to 0$ を微分対象の短完全列とする。
このとき完全ホモロジー列
$$
\ldots \to H(C, d) \to H(A, d) \to H(B, d) \to H(C, d) \to \ldots
$$
を得る。
\end{lemma}

\begin{proof}
補題 \ref{homology-lemma-long-exact-sequence-cochain} を、複体の短完全列
$$
\begin{matrix}
0 & \to & A & \to & B & \to & C & \to & 0 \\
& & \downarrow & & \downarrow & & \downarrow \\
0 & \to & A & \to & B & \to & C & \to & 0 \\
& & \downarrow & & \downarrow & & \downarrow \\
0 & \to & A & \to & B & \to & C & \to & 0
\end{matrix}
$$
に適用する。ここで縦の矢印は $d$ である。
\end{proof}

\noindent
スペクトル系列の重要な例に移る。$\mathcal{A}$ をアーベル圏とし、
$(A, d)$ を $\mathcal{A}$ の微分対象とする。
$\alpha : (A, d) \to (A, d)$ をこの微分対象の自己準同型とする。
$\alpha$ が単射であると仮定すると、微分対象の短完全列
$$
0 \to (A, d) \to (A, d) \to (A/\alpha A, d) \to 0
$$
を得る。補題 \ref{homology-lemma-differential-objects-ses} により完全対
$$
\xymatrix{
H(A, d) \ar[rr]_{\overline{\alpha}} & & H(A, d) \ar[ld]^g \\
& H(A/\alpha A, d) \ar[lu]^f &
}
$$
を得る。ここで $g$ は標準写像、$f$ は蛇の補題で定義される写像である。
節 \ref{homology-section-exact-couple} により、付随するスペクトル系列を得る。
この場合 $E_1 = H(A/\alpha A, d)$ なので、$E_0 = A/\alpha A$ および
$d_0 = d$ と定義するのが自然である。言い換えれば、$r = 0$ から
スペクトル系列を始める。節 \ref{homology-section-spectral-sequence} の規約に従い、部分対象の列
$$
0 = B_0 \subset \ldots \subset B_r \subset \ldots
\subset Z_r \subset \ldots \subset Z_0 = E_0
$$
を $E_r = Z_r/B_r$ となるように定義する。補題
\ref{homology-lemma-spectral-sequence-associated-exact-couple} から、$r \geq 1$ に対して
次が従う。
\begin{enumerate}
\item $B_r$ は、自然な写像 $A \to A/\alpha A$ による
$(\alpha^{r - 1})^{-1}(d A)$ の像である。
\item $Z_r$ は、自然な写像 $A \to A/\alpha A$ による
$d^{-1}(\alpha^r A)$ の像である。
\item $d_r : E_r \to E_r$ は次のように与えられる。元 $z \in Z_r$ に対し、
$d(z) = \alpha^r(y)$ を満たす元 $y \in A$ を選ぶ。このとき
$d_r(z + B_r + \alpha A) = y + B_r + \alpha A$ である。
\end{enumerate}
注意：必ずしも $\alpha A \subset (\alpha^{r - 1})^{-1}(dA)$ でも
$\alpha A \subset d^{-1}(\alpha^r A)$ でもない。
一方、$(\alpha^{r - 1})^{-1}(dA) \subset d^{-1}(\alpha^r A)$ は成り立つ。
また、
$$
E_r
=
\frac{d^{-1}(\alpha^r A) + \alpha A}{(\alpha^{r - 1})^{-1}(dA) + \alpha A}.
$$
である。(1)--(3) がスペクトル系列を与えることを直接確認するのは難しくない。

\begin{definition}
\label{homology-definition-differential-object-selfmap}
$\mathcal{A}$ をアーベル圏とし、$(A, d)$ を $\mathcal{A}$ の微分対象とする。
$\alpha : A \to A$ を $d$ と可換な $A$ の単射自己写像とする。
{\it $(A, d, \alpha)$ に付随するスペクトル系列}とは、上で記述した
スペクトル系列 $(E_r, d_r)_{r \geq 0}$ をいう。
\end{definition}

\begin{remark}[変種]
\label{homology-remark-differential-object-selfmap}
$\mathcal{A}$ をアーベル圏とし、$S, T : \mathcal{A} \to \mathcal{A}$ を
シフト関手、すなわち圏の同型とする。関手として $TS = ST$ であると仮定する。
$\mathcal{A}$ の対象 $A$ と $d \circ S^{-1}d = 0$ を満たす射
$d : A \to SA$ からなる対 $(A, d)$ を考える。これらの対象の圏はアーベル圏である。
$H(A, d) = \Ker(d)/\Im(S^{-1}d)$ と定義すると、
$H(SA, Sd) = SH(A, d)$（標準同型）である。短完全列
$$
0 \to (A, d) \to (B, d) \to (C, d) \to 0
$$
が与えられると、長完全ホモロジー列
$$
\ldots \to S^{-1}H(C, d) \to
H(A, d) \to H(B, d) \to H(C, d) \to SH(A, d) \to \ldots
$$
を得る（境界写像におけるシフトに注意せよ）。$ST = TS$ なので、関手 $T$ は
$T(A, d) = (TA, Td)$ と置くことにより対の上のシフト関手を定める。
次に、$\alpha : (A, d) \to T^{-1}(A, d)$ を余核 $(Q, d)$ をもつ単射とする。
すると注意 \ref{homology-remark-shifted-exact-couple} のような、シフト関手 $TS$ および
$T$ に関する完全対
$$
(H(A, d), S^{-1}H(Q, d), \overline{\alpha}, f, g)
$$
を得る。ここで $\overline{\alpha} : H(A, d) \to T^{-1}H(A, d)$ は $\alpha$ から誘導され、
写像 $f : S^{-1}H(Q, d) \to H(A, d)$ は境界写像、
$g : H(A, d) \to TH(Q, d) = TS(S^{-1}H(Q, d))$ は商写像
$A \to TQ$ から誘導される。したがって、$E_1 = S^{-1}H(Q, d)$ と微分
$d_r : E_r \to T^rSE_r$ をもつ上記のスペクトル系列を得る。上と同じく
$E_0 = S^{-1}Q$ と置き、$d_0 : E_0 \to SE_0$ を
$S^{-1}d : S^{-1}Q \to Q$ によって定める。規約に従って
$B_r \subset Z_r \subset E_0$ と定義すると、$r \geq 1$ に対して次が成り立つ。
\begin{enumerate}
\item $SB_r$ は、
$$
(T^{-r + 1}\alpha \circ \ldots \circ T^{-1}\alpha)^{-1}
\Im(T^{-r}S^{-1}d)
$$
の自然な写像 $T^{-1}A \to Q$ による像である。
\item $Z_r$ は、
$$
(S^{-1}T^{-1}d)^{-1}
\Im(\alpha \circ \ldots \circ T^{r - 1}\alpha)
$$
の自然な写像 $S^{-1}T^{-1}A \to S^{-1}Q$ による像である。
\end{enumerate}
微分は次のように記述できる。$x \in Z_r$ ならば、$x$ に写る
$x' \in S^{-1}T^{-1}A$ を選ぶ。このとき、ある $y \in T^{r - 1}A$ に対して
$S^{-1}T^{-1}d(x')$ は $(\alpha \circ \ldots \circ T^{r - 1}\alpha)(y)$ である。
すると $d_r(x) \in T^rSE_r$ は、$T^rSB_r$ を法とする
$T^rSE_0 = T^rQ$ における $y$ の像の類によって表される。
\end{remark}

\section{スペクトル系列：フィルトレーション付き微分対象}
\label{homology-section-filtered-differential}

\noindent
フィルトレーション付き微分対象から出発してスペクトル系列を構成できる。

\begin{definition}
\label{homology-definition-filtered-differential}
$\mathcal{A}$ をアーベル圏とする。{\it フィルトレーション付き微分対象}
$(K, F, d)$ とは、$\mathcal{A}$ のフィルトレーション付き対象 $(K, F)$ に、
二乗が零となる自己準同型 $d : (K, F) \to (K, F)$、すなわち
$d \circ d = 0$ を満たすものを備えたものをいう。
\end{definition}

\noindent
このような対象に付随するスペクトル系列を記述するため、ひとまず
$\mathcal{A}$ は可算直和をもち、可算直和が完全であるアーベル圏だと仮定する
（これは自動的ではない。注意 \ref{homology-remark-direct-sums-not-exact} を参照せよ）。
$(K, F, d)$ を $\mathcal{A}$ のフィルトレーション付き微分対象とする。
各 $F^nK$ はそれ自身微分対象であることに注意する。
対象 $A = \bigoplus F^nK$ を考え、各直和因子上で $d$ を用いることにより
その上に微分 $d$ を入れる。このとき $(A, d)$ は $\mathcal{A}$ の微分対象であり、
次数付けを備える。包含 $F^nK \to F^{n - 1}K$ によって与えられる写像
$$
\alpha : A \to A
$$
を考える。これは明らかに微分対象の単射
$\alpha : (A, d) \to (A, d)$ である。したがって定義
\ref{homology-definition-differential-object-selfmap} によりスペクトル系列を得る。
これを{\it フィルトレーション付き微分対象 $(K, F, d)$ に付随するスペクトル系列}
と呼ぶ。

\medskip\noindent
このスペクトル系列の項を求めよう。まず、$A/\alpha A = \text{gr}(K)$ であり、
その微分は $d = \text{gr}(d)$ である。したがって
$$
E_0 = \text{gr}(K), \quad d_0 = \text{gr}(d).
$$
である。ゆえに次数付き微分対象 $\text{gr}(K)$ のホモロジーが次の項となる。
$$
E_1 = H(\text{gr}(K), \text{gr}(d)).
$$
さらに、$E_0$ は $\mathcal{A}$ の次数付き対象であり、$d_0$ はその次数付けと
両立する。したがって明らかに $E_1$ も次数付き対象である。しかし微分 $d_1$ は
この次数付けを保たず、次数を $1$ だけシフトすることが分かる。

\medskip\noindent
これを正確に述べるため、
$$
Z_r^p =
\frac{F^pK \cap d^{-1}(F^{p + r}K) + F^{p + 1}K}{F^{p + 1}K}
$$
および
$$
B_r^p =
\frac{F^pK \cap d(F^{p - r + 1}K) + F^{p + 1}K}{F^{p + 1}K}.
$$
と定義する。この記法はきわめて自然ではあるが、文献の大部分の記法とは異なるようである。
もっとも、文献にも一貫した記法の選択はないようなので、問題ではないかもしれない。
この選択により、$B_r \subset E_0$、それぞれ $Z_r \subset E_0$（節
\ref{homology-section-differential-object} で定義したもの）は、$\bigoplus_p B_r^p$、
それぞれ $\bigoplus_p Z_r^p$ に等しい。したがって、$r \geq 0$ および
$p \in \mathbf{Z}$ に対して
$$
E_r^p = Z_r^p/B_r^p
$$
と定義すれば、$E_r = \bigoplus_p E_r^p$ である。微分
$d_r^p : E_r^p \to E_r^{p + r}$ を規則
$$
z + F^{p + 1}K
\longmapsto
dz + F^{p + r + 1}K
$$
によって定義できる。ここで $z \in F^pK \cap d^{-1}(F^{p + r}K)$ である。

\begin{lemma}
\label{homology-lemma-spectral-sequence-filtered-differential}
$\mathcal{A}$ をアーベル圏とし、$(K, F, d)$ を $\mathcal{A}$ の
フィルトレーション付き微分対象とする。$(K, F, d)$ に付随する
$\text{Gr}(\mathcal{A})$ におけるスペクトル系列 $(E_r, d_r)_{r \geq 0}$ であって、
すべての $r$ に対して $d_r : E_r \to E_r[r]$ となり、次数成分 $E_r^p$ と写像
$d_r^p : E_r^p \to E_r^{p + r}$ が上で与えたものとなるものが存在する。
さらに、$E_0^p = \text{gr}^p K$、$d_0^p = \text{gr}^p(d)$、
$E_1^p = H(\text{gr}^pK, \text{gr}^p(d))$ である。
\end{lemma}

\begin{proof}
$\mathcal{A}$ が可算直和をもち、可算直和が完全ならば、これは上の議論から従う。
一般の場合は次のように進めるが、読者にはこの証明を飛ばすことを強く勧める。
$\text{Gr}(\mathcal{A})$ の対象 $A = (F^{p + 1}K)$、すなわち
$F^{p + 1}K$ を次数 $p$ に置いたものを考える（後で番号が正しくなるように
奇妙なシフトを入れている）。各成分上で $d$ を用いて微分 $d$ を入れる。
すると $(A, d)$ は $\text{Gr}(\mathcal{A})$ の微分対象である。次数 $p$ において
包含 $F^{p + 1}A \to F^pA$ によって与えられる写像
% P02-E0015: source line 5773 applies F to A; the construction defined A^p=F^{p+1}K and appears to require F^{p+1}K -> F^pK.
$$
\alpha : A \to A[-1]
$$
を考える。これは明らかに微分対象の単射
$\alpha : (A, d) \to (A, d)[-1]$ である。したがって、$S = \text{id}$、
$T = [1]$ として注意 \ref{homology-remark-differential-object-selfmap} を適用できる。
$\text{Gr}(\mathcal{A})$ における対応するスペクトル系列
$(E_r, d_r)_{r \geq 0}$ が、求めるスペクトル系列である。定義を少し展開しよう。
まず $E_r = (E_r^p)$ は $\text{Gr}(\mathcal{A})$ の対象である。
次に、$T^rS = [r]$ なので $d_r : E_r \to E_r[r]$ であり、これは
$d_r^p : E_r^p \to E_r^{p + r}$ を意味する。

\medskip\noindent
次数成分の記述を確かめるには、上と同様に議論する。まず
$E_0 = \Coker(\alpha : A \to A[-1])$ であり、上の番号付けの選択により
$E_0^p = \text{gr}^pK$ を得る。第一微分は
$d_0^p = \text{gr}^pd : E_0^p \to E_0^p$ によって与えられる。
次に、注意 \ref{homology-remark-differential-object-selfmap} における境界 $B_r$ と
コサイクル $Z_r$ の記述は、上で与えた $Z_r^p$ と $B_r^p$ の公式へ直ちに移される。
\end{proof}

\begin{lemma}
\label{homology-lemma-spectral-sequence-filtered-differential-d1}
$\mathcal{A}$ をアーベル圏とし、$(K, F, d)$ を $\mathcal{A}$ の
フィルトレーション付き微分対象とする。$(K, F, d)$ に付随するスペクトル系列
$(E_r, d_r)_{r \geq 0}$ における
$$
d_1^p :
E_1^p = H(\text{gr}^pK)
\longrightarrow
H(\text{gr}^{p + 1}K) = E_1^{p + 1}
$$
は、微分対象の短完全列
$$
0 \to \text{gr}^{p + 1}K \to F^pK/F^{p + 2}K \to \text{gr}^pK \to 0.
$$
に付随するホモロジーの境界写像に等しい。
\end{lemma}

\begin{proof}
これは補題 \ref{homology-lemma-spectral-sequence-filtered-differential} の直前に与えた
微分 $d_1^p$ の公式から明らかである。
\end{proof}

\begin{definition}
\label{homology-definition-filtration-cohomology-filtered-differential}
$\mathcal{A}$ をアーベル圏とし、$(K, F, d)$ を $\mathcal{A}$ の
フィルトレーション付き微分対象とする。$H(K, d)$ 上の{\it 誘導フィルトレーション}
とは、$F^pH(K, d) = \Im(H(F^pK, d) \to H(K, d))$ によって定義される
フィルトレーションをいう。
\end{definition}

\noindent
これを具体的に書けば、
$$
F^pH(K, d) =
\frac{\Ker(d) \cap F^pK + \Im(d)}{\Im(d)}
$$
であり、したがって
$$
\text{gr}^p H(K) =
\frac{\Ker(d) \cap F^pK + \Im(d)}{\Ker(d) \cap F^{p + 1}K + \Im(d)} =
\frac{\Ker(d) \cap F^pK}{\Ker(d) \cap F^{p + 1}K + \Im(d) \cap F^pK}
$$
である。

\begin{lemma}
\label{homology-lemma-compute-filtered-cohomology}
$\mathcal{A}$ をアーベル圏とし、$(K, F, d)$ を $\mathcal{A}$ の
フィルトレーション付き微分対象とする。$Z_\infty^p$ と $B_\infty^p$ が存在する
ならば（証明を参照）、
\begin{enumerate}
\item 極限 $E_\infty$ が存在して次数付きであり、その次数 $p$ 成分は
$E_\infty^p = Z_\infty^p/B_\infty^p$ である。
\item $K$ のホモロジーの付随次数付き対象 $\text{gr}(H(K))$ は、
次数付き極限対象 $E_\infty$ の次数付き部分商である。
\end{enumerate}
\end{lemma}

\begin{proof}
定義 \ref{homology-definition-limit-spectral-sequence} の対象 $Z_\infty$、$B_\infty$、
および極限 $E_\infty = Z_\infty/B_\infty$ は、補題
\ref{homology-lemma-spectral-sequence-filtered-differential} の証明における
スペクトル系列の構成により $\text{Gr}(\mathcal{A})$ の対象である。
$Z_r^p$、それぞれ $B_r^p$ は $Z_r$、それぞれ $B_r$ の第 $p$ 次数成分なので、
$$
Z_\infty^p = \bigcap\nolimits_r Z_r^p =
\frac{\bigcap_r (F^pK \cap d^{-1}(F^{p + r}K) + F^{p + 1}K)}{F^{p + 1}K}
$$
および
$$
B_\infty^p = \bigcup\nolimits_r B_r^p =
\frac{\bigcup_r (F^pK \cap d(F^{p - r + 1}K) + F^{p + 1}K)}{F^{p + 1}K}.
$$
が存在すると仮定すれば、$Z_\infty$ と $B_\infty$ は、次数 $p$ 成分が
$Z_\infty^p$ と $B_\infty^p$ であるものとして存在する（次数付き部分対象の
和集合と共通部分に関する初等的議論から従う）。したがって
$$
E_\infty^p =
\frac{\bigcap_r (F^pK \cap d^{-1}(F^{p + r}K) + F^{p + 1}K)}
{\bigcup_r (F^pK \cap d(F^{p - r + 1}K) + F^{p + 1}K)}.
$$
であり、ここで分子と分母は存在する。また、
\begin{equation}
\label{homology-equation-on-top}
\Ker(d) \cap F^pK + F^{p + 1}K
\subset
\bigcap\nolimits_r \left(F^pK \cap d^{-1}(F^{p + r}K) + F^{p + 1}K\right)
\end{equation}
および
\begin{equation}
\label{homology-equation-at-bottom}
\bigcup\nolimits_r \left(F^pK \cap d(F^{p - r + 1}K) + F^{p + 1}K\right)
\subset
\Im(d) \cap F^pK + F^{p + 1}K.
\end{equation}
である。したがって $E_\infty^p$ の部分商として
$$
\frac{\Ker(d) \cap F^pK + F^{p + 1}K}{\Im(d) \cap F^pK + F^{p + 1}K} =
\frac{\Ker(d) \cap F^pK}{\Im(d) \cap F^pK + \Ker(d) \cap F^{p + 1}K}
$$
を得る。定義 \ref{homology-definition-filtration-cohomology-filtered-differential} に続く議論で
与えた $\text{gr}^pH(K)$ の公式と比較すれば結論を得る。
\end{proof}

\begin{definition}
\label{homology-definition-filtered-differential-ss-converges}
$\mathcal{A}$ をアーベル圏とし、$(K, F, d)$ を $\mathcal{A}$ の
フィルトレーション付き微分対象とする。$(K, F, d)$ に付随するスペクトル系列が、
\begin{enumerate}
\item 補題 \ref{homology-lemma-compute-filtered-cohomology} を通じて
$\text{gr}H(K) = E_{\infty}$ となるとき、{\it $H(K)$ に弱収束する}という。
\item $H(K)$ に弱収束し、かつ $\bigcap F^pH(K) = 0$ および
$\bigcup F^pH(K) = H(K)$ であるとき、{\it $H(K)$ に終着する}という。
\end{enumerate}
\end{definition}

\noindent
残念ながら、文献においてこれらの概念に関する一貫した用語を見つけるのは難しいようである。

\begin{lemma}
\label{homology-lemma-filtered-differential-ss-converges}
$\mathcal{A}$ をアーベル圏とし、$(K, F, d)$ を $\mathcal{A}$ の
フィルトレーション付き微分対象とする。付随するスペクトル系列について、
\begin{enumerate}
\item すべての $p \in \mathbf{Z}$ に対して式
(\ref{homology-equation-at-bottom}) および (\ref{homology-equation-on-top}) が等号となることと、
$H(K)$ に弱収束することとは同値である。
\item $H(K)$ に終着することと、$H(K)$ に弱収束し、かつ
$\bigcap_p (\Ker(d) \cap F^pK + \Im(d)) = \Im(d)$ および
$\bigcup_p (\Ker(d) \cap F^pK + \Im(d)) = \Ker(d)$ が成り立つこととは同値である。
\end{enumerate}
\end{lemma}

\begin{proof}
上の議論から直ちに従う。
\end{proof}

\section{スペクトル系列：フィルトレーション付き複体}
\label{homology-section-filtered-complex}

\begin{definition}
\label{homology-definition-filtered-complex}
$\mathcal{A}$ をアーベル圏とする。$\mathcal{A}$ の{\it フィルトレーション付き複体
$K^\bullet$}とは、$\text{Fil}(\mathcal{A})$ の複体をいう（定義
\ref{homology-definition-filtered} を参照せよ）。
\end{definition}

\noindent
対象上のフィルトレーションを $F$ と表す。したがって $F^pK^n$ は、複体の
第 $n$ 項のフィルトレーションの第 $p$ 段階を表す。各 $F^pK^\bullet$ は
$\mathcal{A}$ の複体であることに注意する。したがって、フィルトレーション付き複体を
$\mathcal{A}$ の複体からなる（アーベル）圏におけるフィルトレーション付き対象として
定義することもできた。特に $\text{gr} K^\bullet$ は、$\mathcal{A}$ の複体の圏の
次数付き対象である。

\medskip\noindent
このような対象に付随するスペクトル系列を記述するため、ひとまず
$\mathcal{A}$ は可算直和をもち、可算直和が完全であるアーベル圏だと仮定する
（これは自動的ではない。注意 \ref{homology-remark-direct-sums-not-exact} を参照せよ）。
$K$ の微分を $d$ と表す。次数付けを忘れれば、$\bigoplus K^n$ を $\mathcal{A}$ の
フィルトレーション付き微分対象とみなせる。したがって節
\ref{homology-section-filtered-differential} によりスペクトル系列
$(E_r, d_r)_{r \geq 0}$ を得る。この節ではこのスペクトル系列の項を求め、
$K$ の次数付けからくる付加構造を各項に入れる。

\medskip\noindent
まず、$E_0^p = \text{gr}^p K^\bullet$ は複体であり、したがって次数付きである。
ゆえに $E_0$ は自然に二重次数付きである。通常、二重次数付け
$$
E_0 = \bigoplus\nolimits_{p, q} E_0^{p, q},
\quad
E_0^{p, q} = \text{gr}^p K^{p + q}
$$
を用いる。$p + q$ は（コ）ホモロジー類の{\it 全次数}と考える。
また $p$ を{\it フィルトレーション次数}、$q$ を{\it 補次数}と呼ぶ。
微分 $d_0$ はこの二重次数付けと次のように両立する。
$$
d_0  = \bigoplus d_0^{p, q},
\quad
d_0^{p, q} : E_0^{p, q} \to E_0^{p, q + 1}.
$$
すなわち、$d_0^p$ は複体 $\text{gr}^p K^\bullet$ 上の微分そのものである
（これは $\text{gr}^pE_0$ として少しシフトされて現れる）。

\medskip\noindent
さらに進むため、節 \ref{homology-section-filtered-differential} で導入した対象
$B_r^p$ と $Z_r^p$ を次数付き対象として同定し、微分の対応する分解を求める。
これは完全に直接的に行えるが、ここでも本書の記法が他所の記法とは異なることを
読者に注意しておく。次のように定義する。
$$
Z_r^{p, q} =
\frac{F^pK^{p + q} \cap d^{-1}(F^{p + r}K^{p + q + 1}) + F^{p + 1}K^{p + q}}
{F^{p + 1}K^{p + q}}
$$
および
$$
B_r^{p, q} =
\frac{F^pK^{p + q} \cap d(F^{p - r + 1}K^{p + q - 1}) + F^{p + 1}K^{p + q}}
{F^{p + 1}K^{p + q}}
$$
とし、もちろん $E_r^{p, q} = Z_r^{p, q}/B_r^{p, q}$ とする。
これらの定義により、$Z_r^p = \bigoplus_q Z_r^{p, q}$、
$B_r^p = \bigoplus_q B_r^{p, q}$、$E_r^p = \bigoplus_q E_r^{p, q}$ であることは
完全に明らかである。さらに、
$$
0 \subset \ldots \subset B_r^{p, q} \subset
\ldots
\subset Z_r^{p, q} \subset \ldots \subset E_0^{p, q}
$$
である。また、写像 $d_r^p$ は写像
$$
d_r^{p, q} : E_r^{p, q} \longrightarrow E_r^{p + r, q - r + 1},
\quad
z + F^{p + 1}K^{p + q}
\mapsto
dz + F^{p + r + 1}K^{p + q + 1}
$$
の直和に分解する。ここで
$z \in F^pK^{p + q} \cap d^{-1}(F^{p + r}K^{p + q + 1})$ である。

\begin{lemma}
\label{homology-lemma-spectral-sequence-filtered-complex}
$\mathcal{A}$ をアーベル圏とし、$(K^\bullet, F)$ を $\mathcal{A}$ の
フィルトレーション付き複体とする。$(K^\bullet, F)$ に付随する、
$\mathcal{A}$ の二重次数付き対象の圏におけるスペクトル系列
$(E_r, d_r)_{r \geq 0}$ であって、$d_r$ が双次数 $(r, - r + 1)$ をもち、
$E_r$ の二重次数成分 $E_r^{p, q}$ と写像
$d_r^{p, q} : E_r^{p, q} \to E_r^{p + r, q - r + 1}$ が上で与えたものとなるものが
存在する。さらに、$E_0^{p, q} = \text{gr}^p(K^{p + q})$、
$d_0^{p, q} = \text{gr}^p(d^{p + q})$、
$E_1^{p, q} = H^{p + q}(\text{gr}^p(K^\bullet))$ である。
\end{lemma}

\begin{proof}
$\mathcal{A}$ が可算直和をもち、可算直和が完全ならば、これは上の議論から従う。
一般の場合は次のように進めるが、読者にはこの証明を飛ばすことを強く勧める。
$\mathcal{A}$ の二重次数付き対象 $A = (F^{p + 1}K^{p + 1 + q})$、すなわち
$F^{p + 1}K^{p + 1 + q}$ を次数 $(p, q)$ に置いたものを考える
（後で番号が正しくなるように奇妙なシフトを入れている）。各成分上で $d$ を用い、
微分 $d : A \to A[0, 1]$ を入れる。すると $(A, d)$ は微分二重次数付き対象である。
次数 $(p, q)$ において包含
$F^{p + 1}K^{p + 1 + q} \to F^pK^{p + 1 + q}$ によって与えられる写像
$$
\alpha : A \to A[-1, 1]
$$
を考える。これは微分対象の単射
$\alpha : (A, d) \to (A, d)[-1, 1]$ である。したがって、
$S = [0, 1]$、$T = [1, -1]$ として注意
\ref{homology-remark-differential-object-selfmap} を適用できる。二重次数付き対象の対応する
スペクトル系列 $(E_r, d_r)_{r \geq 0}$ が、求めるスペクトル系列である。
定義を少し展開しよう。まず $E_r = (E_r^{p, q})$ である。次に、
$T^rS = [r, -r + 1]$ なので $d_r : E_r \to E_r[r, -r + 1]$ であり、これは
$d_r^{p, q} : E_r^{p, q} \to E_r^{p + r, q - r + 1}$ を意味する。

\medskip\noindent
次数成分の記述を確かめるには、上と同様に議論する。まず
$$
E_0 = \Coker(\alpha : A \to A[-1, 1])[0, -1] =
\Coker(\alpha[0, -1] : A[0, -1] \to A[-1, 0])
$$
であり、上の番号付けの選択により
$$
E_0^{p, q} = \Coker(F^{p + 1}K^{p + q} \to F^pK^{p + q}) = \text{gr}^pK^{p + q}
$$
を得る。第一微分は
$d_0^{p, q} = \text{gr}^pd^{p + q} : E_0^{p, q} \to E_0^{p, q + 1}$ によって
与えられる。次に、注意 \ref{homology-remark-differential-object-selfmap} における境界 $B_r$ と
コサイクル $Z_r$ の記述は、上で与えた $Z_r^{p, q}$ と $B_r^{p, q}$ の公式へ
直ちに移される。
\end{proof}

\begin{lemma}
\label{homology-lemma-spectral-sequence-filtered-complex-d1}
$\mathcal{A}$ をアーベル圏とし、$(K^\bullet, F)$ を $\mathcal{A}$ の
フィルトレーション付き複体とする。$\mathcal{A}$ は可算直和をもつと仮定する。
$(E_r, d_r)_{r \geq 0}$ を $(K^\bullet, F)$ に付随するスペクトル系列とする。
\begin{enumerate}
\item 写像
$$
d_1^{p, q} :
E_1^{p, q} = H^{p + q}(\text{gr}^p(K^\bullet))
\longrightarrow
E_1^{p + 1, q} = H^{p + q + 1}(\text{gr}^{p + 1}(K^\bullet))
$$
は、複体の短完全列
$$
0 \to \text{gr}^{p + 1}(K^\bullet) \to
F^pK^\bullet/F^{p + 2}K^\bullet \to \text{gr}^p(K^\bullet) \to 0.
$$
に付随するコホモロジーの境界写像に等しい。
\item すべての $p \in \mathbf{Z}$ に対して $d(F^pK) \subset F^{p + 1}K$ であると仮定する。
すると $d$ は $\text{gr}^p(K^\bullet)$ 上に零微分を誘導し、したがって
$E_1^{p, q} = \text{gr}^p(K^\bullet)^{p + q}$ である。さらに、この場合
$$
d_1^{p, q} :
E_1^{p, q} = \text{gr}^p(K^\bullet)^{p + q}
\longrightarrow
E_1^{p + 1, q} = \text{gr}^{p + 1}(K^\bullet)^{p + q + 1}
$$
は $d$ から誘導される射である。
\end{enumerate}
\end{lemma}

\begin{proof}
これは補題 \ref{homology-lemma-spectral-sequence-filtered-complex} の直前に与えた微分
$d_1^{p, q}$ の公式から明らかである。
\end{proof}

\begin{lemma}
\label{homology-lemma-spectral-sequence-filtered-complex-functorial}
$\mathcal{A}$ をアーベル圏とし、
$\alpha : (K^\bullet, F) \to (L^\bullet, F)$ を $\mathcal{A}$ の
フィルトレーション付き複体の射とする。$(E_r(K), d_r)_{r \geq 0}$、それぞれ
$(E_r(L), d_r)_{r \geq 0}$ を、$(K^\bullet, F)$、それぞれ
$(L^\bullet, F)$ に付随するスペクトル系列とする。射 $\alpha$ は、二重次数付けと
両立するスペクトル系列の標準射
$\{\alpha_r : E_r(K) \to E_r(L)\}_{r \geq 0}$ を誘導する。
\end{lemma}

\begin{proof}
スペクトル系列の各項の明示的表示から明らかである。
\end{proof}

\begin{definition}
\label{homology-definition-filtration-cohomology-filtered-complex}
$\mathcal{A}$ をアーベル圏とし、$(K^\bullet, F)$ を $\mathcal{A}$ の
フィルトレーション付き複体とする。$H^n(K^\bullet)$ 上の
{\it 誘導フィルトレーション}とは、
$F^pH^n(K^\bullet) = \Im(H^n(F^pK^\bullet) \to H^n(K^\bullet))$ によって
定義されるフィルトレーションをいう。
\end{definition}

\noindent
これを具体的に書けば、
\begin{equation}
\label{homology-equation-filtration-cohomology}
F^pH^n(K^\bullet, d) =
\frac{\Ker(d) \cap F^pK^n + \Im(d) \cap K^n}{\Im(d) \cap K^n}
\end{equation}
であり、したがって
\begin{equation}
\label{homology-equation-graded-cohomology}
\text{gr}^p H^n(K^\bullet) =
\frac{\Ker(d) \cap F^pK^n}{\Ker(d) \cap F^{p + 1}K^n + \Im(d) \cap F^pK^n}
\end{equation}
である（中間の一段階は省略した）。

\begin{lemma}
\label{homology-lemma-compute-cohomology-filtered-complex}
$\mathcal{A}$ をアーベル圏とし、$(K^\bullet, F)$ を $\mathcal{A}$ の
フィルトレーション付き複体とする。$Z_\infty^{p, q}$ と $B_\infty^{p, q}$ が
存在するならば（証明を参照）、
\begin{enumerate}
\item 極限 $E_\infty$ が存在して二重次数付き対象であり、その双次数 $(p, q)$ 成分は
$E_\infty^{p, q} = Z_\infty^{p, q}/B_\infty^{p, q}$ である。
\item $K^\bullet$ の第 $n$ コホモロジー対象の第 $p$ 次数成分
$\text{gr}^pH^n(K^\bullet)$ は $E_\infty^{p, n - p}$ の部分商である。
\end{enumerate}
\end{lemma}

\begin{proof}
定義 \ref{homology-definition-limit-spectral-sequence} の対象 $Z_\infty$、$B_\infty$、
および極限 $E_\infty = Z_\infty/B_\infty$ は、補題
\ref{homology-lemma-spectral-sequence-filtered-complex} におけるスペクトル系列の構成により、
$\mathcal{A}$ の二重次数付き対象である。$Z_r^{p, q}$、それぞれ $B_r^{p, q}$ は
$Z_r$、それぞれ $B_r$ の双次数 $(p, q)$ 成分なので、
$$
Z_\infty^{p, q} = \bigcap\nolimits_r Z_r^{p, q} =
\bigcap\nolimits_r
\frac{F^pK^{p + q} \cap d^{-1}(F^{p + r}K^{p + q + 1}) + F^{p + 1}K^{p + q}}
{F^{p + 1}K^{p + q}}
$$
および
$$
B_\infty^{p, q} = \bigcup\nolimits_r B_r^{p, q} =
\bigcup\nolimits_r
\frac{F^pK^{p + q} \cap d(F^{p - r + 1}K^{p + q - 1}) + F^{p + 1}K^{p + q}}
{F^{p + 1}K^{p + q}}
$$
が存在すると仮定すれば、$Z_\infty$ と $B_\infty$ は双次数 $(p, q)$ 成分が
$Z_\infty^{p, q}$ と $B_\infty^{p, q}$ であるものとして存在する
（二重次数付き対象の和集合と共通部分に関する初等的議論から従う）。したがって
$$
E_\infty^{p, q} =
\frac{\bigcap_r (F^pK^{p + q} \cap d^{-1}(F^{p + r}K^{p + q + 1})
+ F^{p + 1}K^{p + q})}
{\bigcup_r (F^pK^{p + q} \cap d(F^{p - r + 1}K^{p + q - 1})
+ F^{p + 1}K^{p + q})}.
$$
であり、ここで分子と分母は存在する。$n = p + q$ と置くと、
\begin{equation}
\label{homology-equation-on-top-bigraded}
\Ker(d) \cap F^pK^{n} + F^{p + 1}K^{n}
\subset
\bigcap\nolimits_r
\left(
F^pK^{n} \cap d^{-1}(F^{p + r}K^{n + 1}) + F^{p + 1}K^{n}
\right)
\end{equation}
および
\begin{equation}
\label{homology-equation-at-bottom-bigraded}
\bigcup\nolimits_r
\left(
F^pK^{n} \cap d(F^{p - r + 1}K^{n - 1}) + F^{p + 1}K^{n}
\right)
\subset
\Im(d) \cap F^pK^{n} + F^{p + 1}K^{n}.
\end{equation}
である。したがって $E_\infty^{p, q}$ の部分商として
$$
\frac{\Ker(d) \cap F^pK^{n} + F^{p + 1}K^n}
{\Im(d) \cap F^pK^{n} + F^{p + 1}K^{n}} =
\frac{\Ker(d) \cap F^pK^n}{\Im(d) \cap F^pK^n + \Ker(d) \cap F^{p + 1}K^n}
$$
を得る。式 (\ref{homology-equation-graded-cohomology}) と比較すれば結論を得る。
\end{proof}

\begin{definition}
\label{homology-definition-bounded-ss}
$\mathcal{A}$ をアーベル圏とし、$(E_r, d_r)_{r \geq r_0}$ を、
$d_r$ が双次数 $(r, -r + 1)$ をもつ $\mathcal{A}$ の二重次数付き対象の
スペクトル系列とする。このようなスペクトル系列が、
\begin{enumerate}
\item すべての $p, q \in \mathbf{Z}$ に対して、$r \geq b$ ならば写像
$d_r^{p, q} : E_r^{p, q} \to E_r^{p + r, q - r + 1}$ が零となるような
$b = b(p, q)$ が存在するとき、{\it 正則}であるという。
\item すべての $p, q \in \mathbf{Z}$ に対して、$r \geq b$ ならば写像
$d_r^{p - r, q + r - 1} : E_r^{p - r, q + r - 1} \to E_r^{p, q}$ が零となるような
$b = b(p, q)$ が存在するとき、{\it 余正則}であるという。
\item すべての $n$ に対して、非零な $E_{r_0}^{p, n - p}$ が有限個しかないとき、
{\it 有界}であるという。
\item すべての $n$ に対して、$p \geq b$ ならば $E_{r_0}^{p, n - p} = 0$ となるような
$b = b(n)$ が存在するとき、{\it 下に有界}であるという。
\item すべての $n$ に対して、$p \leq b$ ならば $E_{r_0}^{p, n - p} = 0$ となるような
$b = b(n)$ が存在するとき、{\it 上に有界}であるという。
\end{enumerate}
\end{definition}

\noindent
下に有界であるとは、直線 $p + q = n$（傾き $-1$）上で $E_r^{p, q}$ を見ると、
$(p, q)$ が右下へ動くにつれて零になることを意味する。上で述べたように、
文献にはこれらの概念に関する一貫した用語がない。

\begin{lemma}
\label{homology-lemma-relate-boundedness}
定義 \ref{homology-definition-bounded-ss} の状況を考える。節
\ref{homology-section-spectral-sequence} のように定義される $Z_r, B_r$ の
$(p, q)$ 次数成分を $Z_r^{p, q}, B_r^{p, q} \subset E_{r_0}^{p, q}$ とする。
\begin{enumerate}
\item スペクトル系列が正則であるための必要十分条件は、すべての $p, q$ に対して
$Z_r^{p, q} = Z_{r + 1}^{p, q} = \ldots$ となる $r = r(p, q)$ が存在することである。
\item スペクトル系列が余正則であるための必要十分条件は、すべての $p, q$ に対して
$B_r^{p, q} = B_{r + 1}^{p, q} = \ldots$ となる $r = r(p, q)$ が存在することである。
\item スペクトル系列が有界であるための必要十分条件は、下にも上にも有界であることである。
\item スペクトル系列が下に有界ならば正則である。
\item スペクトル系列が上に有界ならば余正則である。
\end{enumerate}
\end{lemma}

\begin{proof}
省略する。ヒント：$E_r^{p, q} = 0$ ならば、すべての $r' \geq r$ に対して
$E_{r'}^{p, q} = 0$ である。
\end{proof}

\begin{definition}
\label{homology-definition-filtered-complex-ss-converges}
$\mathcal{A}$ をアーベル圏とし、$(K^\bullet, F)$ を $\mathcal{A}$ の
フィルトレーション付き複体とする。$(K^\bullet, F)$ に付随するスペクトル系列が、
\begin{enumerate}
\item すべての $p, n \in \mathbf{Z}$ に対して、補題
\ref{homology-lemma-compute-cohomology-filtered-complex} を通じて
$\text{gr}^pH^n(K^\bullet) = E_{\infty}^{p, n - p}$ となるとき、
{\it $H^*(K^\bullet)$ に弱収束する}という。
\item $H^*(K^\bullet)$ に弱収束し、すべての $n$ に対して
$\bigcap_p F^pH^n(K^\bullet) = 0$ および
$\bigcup_p F^p H^n(K^\bullet) = H^n(K^\bullet)$ であるとき、
{\it $H^*(K^\bullet)$ に終着する}という。
\item 正則であり、$H^*(K^\bullet)$ に終着し、かつ
$H^n(K^\bullet) = \lim_p H^n(K^\bullet)/F^pH^n(K^\bullet)$ であるとき、
{\it $H^*(K^\bullet)$ に収束する}という。
\end{enumerate}
\end{definition}

\noindent
弱収束、終着、または収束は、記法 $E_r^{p, q} \Rightarrow H^{p + q}(K^\bullet)$ で表す。
上で述べたように、文献にはこれらの概念に関する一貫した用語がない。

\begin{lemma}
\label{homology-lemma-filtered-complex-ss-converges}
$\mathcal{A}$ をアーベル圏とし、$(K^\bullet, F)$ を $\mathcal{A}$ の
フィルトレーション付き複体とする。付随するスペクトル系列について、
\begin{enumerate}
\item すべての $p, q \in \mathbf{Z}$ に対して式
(\ref{homology-equation-at-bottom-bigraded}) および (\ref{homology-equation-on-top-bigraded}) が
等号となることと、$H^*(K^\bullet)$ に弱収束することとは同値である。
\item $H^*(K)$ に終着することと、$H^*(K^\bullet)$ に弱収束し、かつ
$\bigcap_p (\Ker(d) \cap F^pK^n + \Im(d) \cap K^n) = \Im(d) \cap K^n$
および
$\bigcup_p (\Ker(d) \cap F^pK^n + \Im(d) \cap K^n) = \Ker(d) \cap K^n$
が成り立つこととは同値である。
\end{enumerate}
\end{lemma}

\begin{proof}
上の議論から直ちに従う。
\end{proof}

\begin{lemma}
\label{homology-lemma-biregular-ss-converges}
$\mathcal{A}$ をアーベル圏とし、$(K^\bullet, F)$ を $\mathcal{A}$ の
フィルトレーション付き複体とする。各 $K^n$ 上のフィルトレーションが有限であると
仮定する（定義 \ref{homology-definition-filtered} を参照せよ）。このとき、
\begin{enumerate}
\item $(K^\bullet, F)$ に付随するスペクトル系列は有界である。
\item 各 $H^n(K^\bullet)$ 上のフィルトレーションは有限である。
\item $(K^\bullet, F)$ に付随するスペクトル系列は $H^*(K^\bullet)$ に収束する。
\item $\mathcal{C} \subset \mathcal{A}$ が弱セール部分圏であり、ある $r$ に対して
すべての $p, q \in \mathbf{Z}$ について $E_r^{p, q} \in \mathcal{C}$ ならば、
$H^n(K^\bullet)$ は $\mathcal{C}$ に属する。
\end{enumerate}
\end{lemma}

\begin{proof}
(1) は $E_0^{p, n - p} = \text{gr}^p K^n$ から従う。(2) は式
(\ref{homology-equation-filtration-cohomology}) から明らかである。スペクトル系列が弱収束する
ことを証明するため、補題 \ref{homology-lemma-filtered-complex-ss-converges} を用いる。
$p, n \in \mathbf{Z}$ を固定する。$r \gg 0$ に対して $F^{p + r}K^n = 0$ なので、
(\ref{homology-equation-on-top-bigraded}) の右辺は
$F^pK^n \cap \Ker(d) + F^{p + 1}K^n$ に等しい。したがって
(\ref{homology-equation-on-top-bigraded}) は等号である。また $r \gg 0$ に対して
$F^{p - r + 1}K^{n - 1} = K^{n - 1}$ なので、
(\ref{homology-equation-at-bottom-bigraded}) の左辺は
$F^pK^n \cap \Im(d) + F^{p + 1}K^n$ に等しい。したがって
(\ref{homology-equation-at-bottom-bigraded}) は等号である。(2) により
$H^n(K^\bullet)$ 上のフィルトレーションは有限なので、終着性を得る。
収束性を得るには、スペクトル系列が正則であることを示す必要があるが、これは有界性から
従う（補題 \ref{homology-lemma-relate-boundedness}）。さらに
$H^n(K^\bullet) = \lim_p H^n(K^\bullet)/F^pH^n(K^\bullet)$ を示す必要があるが、
これは (2) で証明した $H^*(K^\bullet)$ 上のフィルトレーションの有限性から従う。

\medskip\noindent
(4) を証明する。ある $r \geq 0$ に対して、$\mathcal{A}$ のある弱セール部分圏
$\mathcal{C}$ について $E_r^{p, q} \in \mathcal{C}$ であると仮定する。
すると補題 \ref{homology-lemma-characterize-weak-serre-subcategory} により
$E_{r + 1}^{p, q}$ も $\mathcal{C}$ に属する。上で証明した有界性
（補題 \ref{homology-lemma-relate-boundedness} によりスペクトル系列が正則かつ余正則であることを
含意する）から、$p + q = n$ を満たすすべての $p, q$ に対して
$E_\infty^{p, q} = E_{r'}^{p, q}$ となる $r' \geq r$ を選べる。
したがって $H^n(K^\bullet)$ は、次数成分が $\mathcal{C}$ に属する有限
フィルトレーションをもつ $\mathcal{A}$ の対象である。補題
\ref{homology-lemma-characterize-weak-serre-subcategory} により、これは
$H^n(K^\bullet)$ が $\mathcal{C}$ に属することを含意する。
\end{proof}

\begin{lemma}
\label{homology-lemma-biregular-ss-relation-in-K0}
$\mathcal{A}$ をアーベル圏とし、$(K^\bullet, F)$ を $\mathcal{A}$ の
フィルトレーション付き複体とする。各 $K^n$ 上のフィルトレーションが有限であり
（定義 \ref{homology-definition-filtered} を参照せよ）、ある $r$ に対して非零な
$E_r^{p, q}$ が有限個しかないと仮定する。このとき非零な $H^n(K^\bullet)$ は
有限個しかなく、
$$
\sum (-1)^n[H^n(K^\bullet)] = \sum (-1)^{p + q} [E_r^{p, q}]
$$
が $K_0(\mathcal{A}')$ において成り立つ。ここで $\mathcal{A}'$ は対象
$E_r^{p, q}$ を含む $\mathcal{A}$ の最小の弱セール部分圏である。
\end{lemma}

\begin{proof}
$E_r$ の偶部分と奇部分、すなわち $p + q$ が偶数および奇数となる
$(p, q)$ 成分の直和を $E_r^{even}$ および $E_r^{odd}$ と表す。
微分 $d_r$ は写像 $\varphi : E_r^{even} \to E_r^{odd}$ および
$\psi : E_r^{odd} \to E_r^{even}$ を定め、どちらの順の合成も零となる。
したがって
\begin{align*}
[E_r^{even}] - [E_r^{odd}] & =
[\Ker(\varphi)] + [\Im(\varphi)] - [\Ker(\psi)] - [\Im(\psi)] \\
& =
[\Ker(\varphi)/\Im(\psi)] - [\Ker(\psi)/\Im(\varphi)] \\
& =
[E_{r + 1}^{even}] - [E_{r + 1}^{odd}]
\end{align*}
である。途中に現れる対象はすべて、対象 $E_r^{p, q}$ を含む最小のセール部分圏に
属することに注意する。この方法を続ければ $r$ を任意に増加させられる。
$E_r^{p, q}$ が非零となる対 $(p, q)$ は有限個しかなく、この性質は
$E_{r + 1}, E_{r + 2}, \ldots$ に受け継がれるので、$d_r = 0$ と仮定してよい。
この段階で、$H^n(K^\bullet)$ は次数成分がちょうど $E_r^{p, n - p}$ である
有限フィルトレーションをもつ（補題 \ref{homology-lemma-biregular-ss-converges}）。
したがって結果は明らかである。
\end{proof}

\noindent
次の補題は、本書の定義に対する一種の健全性確認である。すべての $n$ に対して
$$
H^n(F^pK^\bullet) = 0\text{ for }p \gg 0
\quad\text{and}\quad
H^n(F^pK^\bullet) = H^n(K^\bullet)\text{ for }p \ll 0,
$$
となるフィルトレーション付き複体があるならば、対応するスペクトル系列は当然
収束するはずではないだろうか。

\begin{lemma}
\label{homology-lemma-ss-converges-trivial}
$\mathcal{A}$ をアーベル圏とし、$(K^\bullet, F)$ を $\mathcal{A}$ の
フィルトレーション付き複体とする。次を仮定する。
\begin{enumerate}
\item すべての $n$ に対し、$p \geq p_0(n)$ ならば
$H^n(F^pK^\bullet) = 0$ となる $p_0(n)$ が存在する。
\item すべての $n$ に対し、$p \leq p_1(n)$ ならば
$H^n(F^pK^\bullet) \to H^n(K^\bullet)$ が同型となる $p_1(n)$ が存在する。
\end{enumerate}
このとき、
\begin{enumerate}
\item $(K^\bullet, F)$ に付随するスペクトル系列は有界である。
\item 各 $H^n(K^\bullet)$ 上のフィルトレーションは有限である。
\item $(K^\bullet, F)$ に付随するスペクトル系列は $H^*(K^\bullet)$ に収束する。
\end{enumerate}
\end{lemma}

\begin{proof}
$n$ を固定する。複体の短完全列
$$
0 \to F^{p + 1}K^\bullet \to F^pK^\bullet \to \text{gr}^pK^\bullet \to 0
$$
に付随する長完全コホモロジー列を用いると、
$p \geq \max(p_0(n), p_0(n + 1))$ および
$p < \min(p_1(n), p_1(n + 1))$ に対して $E_1^{p, n - p} = 0$ であることが分かる。
したがってスペクトル系列は有界である（定義 \ref{homology-definition-bounded-ss}）。
これで (1) が証明された。

\medskip\noindent
仮定と定義 \ref{homology-definition-filtration-cohomology-filtered-complex} から、
$H^n(K^\bullet)$ 上のフィルトレーションが有限であることは明らかである。
これで (2) が証明された。

\medskip\noindent
次に補題 \ref{homology-lemma-filtered-complex-ss-converges} を用いて、スペクトル系列が
$H^*(K^\bullet)$ に弱収束することを証明する。
(\ref{homology-equation-on-top-bigraded}) が等号となることを示そう。
実際、$p + r > p_0(n + 1)$ のとき、複体 $F^{p + r}K^\bullet$ は次数 $n + 1$ で
完全なので、写像
$$
d : F^pK^{n} \cap d^{-1}(F^{p + r}K^{n + 1}) \to F^{p + r}K^{n + 1}
$$
の像は $d : F^{p + r}K^n \to F^{p + r}K^{n + 1}$ の像に含まれる。
したがって
$F^pK^{n} \cap d^{-1}(F^{p + r}K^{n + 1}) =
d(F^{p + r}K^n) + \Ker(d) \cap F^pK^n$ である。ゆえに、このような $r$ に対して
$$
\Ker(d) \cap F^pK^{n} + F^{p + 1}K^{n} =
F^pK^{n} \cap d^{-1}(F^{p + r}K^{n + 1}) + F^{p + 1}K^{n}
$$
であり、所望の等号が証明される。(\ref{homology-equation-at-bottom-bigraded}) が等号となることを
示すため、$p - r + 1 < p_1(n - 1)$ のとき
$$
d(F^{p - r + 1}K^{n - 1}) = \Im(d) \cap F^{p - r + 1}K^n
$$
であることを用いる。これは写像 $F^{p - r + 1}K^\bullet \to K^\bullet$ が次数
$n - 1$ のコホモロジー上に同型を誘導するからである。したがって、このような $r$ に対して
$$
F^pK^{n} \cap d(F^{p - r + 1}K^{n - 1}) + F^{p + 1}K^{n} =
\Im(d) \cap F^pK^{n} + F^{p + 1}K^{n}
$$
となり、所望の等号が証明される。

\medskip\noindent
補題 \ref{homology-lemma-filtered-complex-ss-converges} を用いて、スペクトル系列が
$H^*(K^\bullet)$ に終着することを見るには、
$\bigcap_p (\Ker(d) \cap F^pK^n + \Im(d) \cap K^n) = \Im(d) \cap K^n$
および
$\bigcup_p (\Ker(d) \cap F^pK^n + \Im(d) \cap K^n) = \Ker(d) \cap K^n$
を示す必要がある。$p \geq p_0(n)$ ならば
$\Ker(d) \cap F^pK^n + \Im(d) \cap K^n = \Im(d) \cap K^n$ であり、
$p \leq p_1(n)$ ならば
$\Ker(d) \cap F^pK^n + \Im(d) \cap K^n = \Ker(d) \cap K^n$ である。
弱収束、終着性、有界性を合わせれば、(2) と (3) が成り立つことが分かる。
\end{proof}

\section{スペクトル系列：二重複体}
\label{homology-section-double-complex}

\noindent
$K^{\bullet, \bullet}$ を二重複体とする。Section \ref{homology-section-double-complexes}
を参照されたい。慣例として $H^p_I(K^{\bullet, \bullet})$ は、
項が $\Ker(d_1^{p, q})/\Im(d_1^{p - 1, q})$（$q$ を動かす）であり、
微分が $d_2$ から誘導される複体を表す。
このとき $H^q_{II}(H^p_I(K^{\bullet, \bullet}))$ は、その次数 $q$ の
コホモロジーを表す。また慣例として $H^q_{II}(K^{\bullet, \bullet})$ は、
項が $\Ker(d_2^{p, q})/\Im(d_2^{p, q - 1})$（$p$ を動かす）であり、
微分が $d_1$ から誘導される複体を表す。
このとき $H^p_I(H^q_{II}(K^{\bullet, \bullet}))$ は、その次数 $p$ の
コホモロジーを表す。これらのコホモロジー群は、付随する全複体、すなわち
単純複体上のフィルトレーションに対するスペクトル系列の項として現れることが
後で分かる。Definition \ref{homology-definition-associated-simple-complex} を参照されたい。

\medskip\noindent
二重複体 $K^{\bullet, \bullet}$ に付随する全複体
$\text{Tot}(K^{\bullet, \bullet})$ 上には、二つの自然なフィルトレーションがある。
すなわち、
$$
F_I^p(\text{Tot}^n(K^{\bullet, \bullet})) =
\bigoplus\nolimits_{i + j = n, \ i \geq p} K^{i, j}
\quad
\text{and}
\quad
F_{II}^p(\text{Tot}^n(K^{\bullet, \bullet})) =
\bigoplus\nolimits_{i + j = n, \ j \geq p} K^{i, j}.
$$
と定める。$(\text{Tot}(K^{\bullet, \bullet}), F_I)$ および
$(\text{Tot}(K^{\bullet, \bullet}), F_{II})$ がフィルトレーション付き複体で
あることは直ちに確かめられる。Section \ref{homology-section-filtered-complex} により、
二つのスペクトル系列を得る。慣例として、フィルトレーション $F_I$ に付随する
スペクトル系列を $({}'E_r, {}'d_r)_{r \geq 0}$ と書き、フィルトレーション
$F_{II}$ に付随するスペクトル系列を $({}''E_r, {}''d_r)_{r \geq 0}$ と書く。
これらのスペクトル系列は次のように記述される。

\begin{lemma}
\label{homology-lemma-ss-double-complex}
$\mathcal{A}$ をアーベル圏とし、$K^{\bullet, \bullet}$ を二重複体とする。
$K^{\bullet, \bullet}$ に付随するスペクトル系列の項は次のとおりである。
\begin{enumerate}
\item ${}'E_0^{p, q} = K^{p, q}$ であり、
${}'d_0^{p, q} = (-1)^p d_2^{p, q} : K^{p, q} \to K^{p, q + 1}$,
\item ${}''E_0^{p, q} = K^{q, p}$ であり、
${}''d_0^{p, q} = d_1^{q, p} : K^{q, p} \to K^{q + 1, p}$,
\item ${}'E_1^{p, q} = H^q(K^{p, \bullet})$ であり、
${}'d_1^{p, q} = H^q(d_1^{p, \bullet})$,
\item ${}''E_1^{p, q} = H^q(K^{\bullet, p})$ であり、
${}''d_1^{p, q} = (-1)^q H^q(d_2^{\bullet, p})$,
\item ${}'E_2^{p, q} = H^p_I(H^q_{II}(K^{\bullet, \bullet}))$,
\item ${}''E_2^{p, q} = H^p_{II}(H^q_I(K^{\bullet, \bullet}))$.
\end{enumerate}
\end{lemma}

\begin{proof}
省略する。
\end{proof}

\noindent
これらのスペクトル系列は $H^n(\text{Tot}(K^{\bullet, \bullet}))$ 上に
二つのフィルトレーションを定める。これらを $F_I$ および $F_{II}$ と書く。

\begin{definition}
\label{homology-definition-ss-double-complex-converge}
$\mathcal{A}$ をアーベル圏とし、$K^{\bullet, \bullet}$ を二重複体とする。
Definition \ref{homology-definition-filtered-complex-ss-converges} が適用されるとき、
スペクトル系列 $({}'E_r, {}'d_r)_{r \geq 0}$ は
{\it $H^n(\text{Tot}(K^{\bullet, \bullet}))$ に弱収束する}、
{\it $H^n(\text{Tot}(K^{\bullet, \bullet}))$ に終着する}、または
{\it $H^n(\text{Tot}(K^{\bullet, \bullet}))$ に収束する}
という。同様に、Definition \ref{homology-definition-filtered-complex-ss-converges} が
適用されるとき、スペクトル系列 $({}''E_r, {}''d_r)_{r \geq 0}$ は
{\it $H^n(\text{Tot}(K^{\bullet, \bullet}))$ に弱収束する}、
{\it $H^n(\text{Tot}(K^{\bullet, \bullet}))$ に終着する}、または
{\it $H^n(\text{Tot}(K^{\bullet, \bullet}))$ に収束する}
という。
\end{definition}

\noindent
上で述べたように、文献ではこれらの概念について用語が統一されていない。
定義の状況で、すべての $n$ に対して
$$
\text{gr}_{F_I}(H^n(\text{Tot}(K^{\bullet, \bullet}))) =
\oplus_{p + q = n} {}'E_\infty^{p, q}
$$
が Lemma \ref{homology-lemma-compute-cohomology-filtered-complex} の標準的比較によって
成り立つならば、第1のスペクトル系列は弱収束する。
同様に、すべての $n$ に対して
$$
\text{gr}_{F_{II}}(H^n(\text{Tot}(K^{\bullet, \bullet}))) =
\oplus_{p + q = n} {}''E_\infty^{p, q}
$$
が Lemma \ref{homology-lemma-compute-cohomology-filtered-complex} の標準的比較によって
成り立つならば、第2のスペクトル系列 $({}''E_r, {}''d_r)_{r \geq 0}$ は
弱収束する。

\begin{lemma}
\label{homology-lemma-first-quadrant-ss}
$\mathcal{A}$ をアーベル圏とし、$K^{\bullet, \bullet}$ を二重複体とする。
各 $n \in \mathbf{Z}$ に対し、$p + q = n$ を満たす非零の
$K^{p, q}$ は有限個しかないと仮定する。このとき、
\begin{enumerate}
\item $K^{\bullet, \bullet}$ に付随する二つのスペクトル系列は有界である。
\item 各 $H^n(\text{Tot}(K^{\bullet, \bullet}))$ 上のフィルトレーション
$F_I$, $F_{II}$ は有限である。
\item スペクトル系列 $({}'E_r, {}'d_r)_{r \geq 0}$ および
$({}''E_r, {}''d_r)_{r \geq 0}$ は
$H^*(\text{Tot}(K^{\bullet, \bullet}))$ に収束する。
\item $\mathcal{C} \subset \mathcal{A}$ を弱セール部分圏とし、ある $r$ に対して
すべての $p, q \in \mathbf{Z}$ について ${}'E_r^{p, q} \in \mathcal{C}$
であるとする。このとき $H^n(\text{Tot}(K^{\bullet, \bullet}))$ は
$\mathcal{C}$ に属する。$({}''E_r, {}''d_r)_{r \geq 0}$ についても同様である。
\end{enumerate}
\end{lemma}

\begin{proof}
Lemma \ref{homology-lemma-biregular-ss-converges} から直ちに従う。
\end{proof}

\noindent
スペクトル系列の最初の応用を与える。

\begin{lemma}
\label{homology-lemma-double-complex-gives-resolution}
$\mathcal{A}$ をアーベル圏とする。$K^\bullet$ を複体とし、
$A^{\bullet, \bullet}$ を二重複体とする。
$\alpha^p : K^p \to A^{p, 0}$ を射とする。次を仮定する。
\begin{enumerate}
\item 各 $n \in \mathbf{Z}$ に対し、$p + q = n$ を満たす非零の
$A^{p, q}$ は有限個しかない。
\item $q < 0$ ならば $A^{p, q} = 0$ である。
\item 射 $\alpha^p$ は複体の射
$\alpha : K^\bullet \to A^{\bullet, 0}$ を与える。
\item 複体 $A^{p, \bullet}$ はすべての次数 $q \not = 0$ で完全であり、
射 $K^p \to A^{p, 0}$ は同型 $K^p \to \Ker(d_2^{p, 0})$ を誘導する。
\end{enumerate}
このとき $\alpha$ は複体の擬同型
$$
K^\bullet \longrightarrow \text{Tot}(A^{\bullet, \bullet})
$$
を誘導する。
さらに、この補題には、第1変数 $p$ の代わりに第2変数 $q$ を用いる変形がある。
\end{lemma}

\begin{proof}
この写像は、射
$K^n \to A^{n, 0} \to \text{Tot}^n(A^{\bullet, \bullet})$
によって与えられる写像にほかならず、これらが複体の射を定めることは
容易に分かる。二重複体 $A^{\bullet, \bullet}$ に付随するスペクトル系列
$({}'E_r, {}'d_r)_{r \geq 0}$ を考える。
Lemma \ref{homology-lemma-first-quadrant-ss} により、このスペクトル系列は収束し、
$H^n(\text{Tot}(A^{\bullet, \bullet}))$ 上に誘導されるフィルトレーションは
各 $n$ に対して有限である。
Lemma \ref{homology-lemma-ss-double-complex} と仮定 (4) により、
$q = 0$ でない限り ${}'E_1^{p, q} = 0$ であり、
$'E_1^{p, 0} = K^p$ で、その微分 ${}'d_1^{p, 0}$ は
$d_K^p$ と同一視される。したがって ${}'E_2^{p, 0} = H^p(K^\bullet)$ であり、
それ以外では零である。次数の理由から、これは明らかに
${}'d_2^{p, q} = {}'d_3^{p, q} = \ldots = 0$ を含意する。
ゆえに $H^n(\text{Tot}(A^{\bullet, \bullet})) = H^n(K^\bullet)$ と結論する。
この同一視が上で導入した複体の射
$K^\bullet \to \text{Tot}(A^{\bullet, \bullet})$ によって与えられることの
確認は省略する。
\end{proof}

\begin{lemma}
\label{homology-lemma-homotopy-complex-complexes}
$\mathcal{A}$ をアーベル圏とし、$M^\bullet$ を $\mathcal{A}$ の複体とする。
$$
a :
M^\bullet[0]
\longrightarrow
\left(A^{0, \bullet} \to A^{1, \bullet} \to A^{2, \bullet} \to \ldots \right)
$$
を $\mathcal{A}$ の複体の複体の圏におけるホモトピー同値とする。
このとき、$M^\bullet \to A^{0, \bullet}$ から誘導される写像
$\alpha : M^\bullet \to \text{Tot}(A^{\bullet, \bullet})$ は
ホモトピー同値である。
\end{lemma}

\begin{proof}
複体の複体は二重複体と同じものであるから、主張には意味がある。
仮定は、複体の複体の圏において $a \circ b$ と $b \circ a$ が恒等写像に
ホモトピックとなるような写像
$$
b :
\left(A^{0, \bullet} \to A^{1, \bullet} \to A^{2, \bullet} \to \ldots \right)
\longrightarrow
M^\bullet[0]
$$
が存在することを意味する。したがって $b \circ a$ は
$M^\bullet[0]$ の恒等写像である（次数 $0$ に項が一つしかないため）。
また、$b$ は写像 $b^0 : A^{0, \bullet} \to M^\bullet$ と、それ以外の
すべての次数における零写像によって与えられることに注意する。
ゆえに $b$ は写像
$\beta : \text{Tot}(A^{\bullet, \bullet}) \to M^\bullet$ を誘導し、
$\beta \circ \alpha$ は $M^\bullet$ 上の恒等写像である。
最後に、写像 $\alpha \circ \beta$ が恒等写像にホモトピックであることを
示さなければならない。そのため、仮定により存在する複体の写像
$h^n : A^{n, \bullet} \to A^{n - 1, \bullet}$ を、
$a \circ b - \text{id} = d_1 \circ h + h \circ d_1$ となるように選ぶ。
ここで $d_1 : A^{n, \bullet} \to A^{n + 1, \bullet}$ は
複体の複体の微分である。また、すべての $n$ に対し、
複体 $A^{n, \bullet}$ の微分も $d_2$ と書く。
$h^n$ の成分を $h^{n, m} : A^{n, m} \to A^{n - 1, m}$ とする。
すると、各直和成分 $A^{n, m}$ 上で $h^{n, m}$ によって与えられる写像
$$
h' : \text{Tot}(A^{\bullet, \bullet})^k =
\bigoplus\nolimits_{n + m = k} A^{n, m}
\to
\bigoplus\nolimits_{n + m = k - 1} A^{n, m} =
\text{Tot}(A^{\bullet, \bullet})^{k - 1}
$$
を考えられる。このとき、写像
$$
d_{\text{Tot}(A^{\bullet, \bullet})} \circ h' +
h' \circ d_{\text{Tot}(A^{\bullet, \bullet})}
$$
を直和成分 $A^{n, m}$ に制限したものは
$$
d_1^{n - 1, m} \circ h^{n, m} +
(-1)^{n - 1} d_2^{n - 1, m} \circ h^{n, m} +
h^{n + 1, m} \circ d_1^{n, m} + h^{n, m + 1} \circ (-1)^nd_2^{n, m}
$$
に等しい。$h^n$ は複体の写像であるから、
$(-1)^{n - 1} d_2^{n - 1, m} \circ h^{n, m}$ と
$h^{n, m + 1} \circ (-1)^nd_2^{n, m}$ は相殺する。
残る二項は $a \circ b - \text{id} = d_1 \circ h + h \circ d_1$
により $(\alpha \circ \beta)|_{A^{n, m}} - \text{id}_{A^{n, m}}$ を与える。
これで証明が完了する。
\end{proof}

\section{アーベル群の二重複体}
\label{homology-section-double-complexes-abelian-groups}

\noindent
この節では、一般のアーベル圏に対する類似の結果が（まだ）得られていない、
アーベル群の二重複体に関するいくつかの結果を述べる。
基礎となるアーベル圏がアーベル群の圏またはこれに類する圏
（たとえば環上の加群の圏）である場合を除き、これらの補題を用いないよう
注意されたい。いくつかの議論は、紙片に「ジグザグ」を描かなければ
追いにくいだろう。Algebra, Lemma \ref{algebra-lemma-no-spectral-sequence}
の証明と比較されたい。

\begin{lemma}
\label{homology-lemma-right-resolution-gives-qis}
$M^\bullet$ をアーベル群の複体とする。
$$
0 \to M^\bullet \to A_0^\bullet \to A_1^\bullet \to A_2^\bullet \to \ldots
$$
をアーベル群の複体の完全な複体とする。
$A^{p, q} = A_p^q$ と置いて二重複体を得る。
このとき、$M^\bullet \to A_0^\bullet$ から誘導される写像
$M^\bullet \to \text{Tot}(A^{\bullet, \bullet})$ は擬同型である。
\end{lemma}

\begin{proof}
$q < t$ ならば $A_0^q = 0$ となる $t \in \mathbf{Z}$ が存在する場合、
これは Lemma \ref{homology-lemma-double-complex-gives-resolution} から直ちに従う
（同補題の最後の主張にあるように $p$ と $q$ を入れ替える）。
一方、すべての $t \in \mathbf{Z}$ に対し、素朴な切断からなる複体
$$
0 \to
\sigma_{\geq t}M^\bullet \to
\sigma_{\geq t}A_0^\bullet \to
\sigma_{\geq t}A_1^\bullet \to
\sigma_{\geq t}A_2^\bullet \to \ldots
$$
がある。対応する二重複体を $A(t)^{\bullet, \bullet}$ と書く。
$H^n(\text{Tot}(A^{\bullet, \bullet}))$ の任意の元 $\xi$ は、ある $t$ に対する
$H^n(\text{Tot}(A(t)^{\bullet, \bullet}))$ の元の像である
（コホモロジー類の具体的な代表元を見ればよい）。
したがって $\xi$ は $H^n(\sigma_{\geq t}M^\bullet)$ の像に入る。
ゆえに写像
$H^n(M^\bullet) \to H^n(\text{Tot}(A^{\bullet, \bullet}))$ は全射である。
また、すべての $t$ に対して写像
$H^n(\sigma_{\geq t}M^\bullet) \to H^n(\text{Tot}(A(t)^{\bullet, \bullet}))$
が単射であることと同様の議論により、この写像は単射である。
\end{proof}

\begin{lemma}
\label{homology-lemma-good-resolution-gives-qis}
$M^\bullet$ をアーベル群の複体とする。
$$
\ldots \to A_2^\bullet \to A_1^\bullet \to A_0^\bullet \to M^\bullet \to 0
$$
をアーベル群の複体の完全な複体とし、すべての $p \in \mathbf{Z}$ に対して
複体
$$
\ldots \to \Ker(d_{A_2^\bullet}^p) \to \Ker(d_{A_1^\bullet}^p)
\to \Ker(d_{A_0^\bullet}^p) \to \Ker(d_{M^\bullet}^p) \to 0
$$
も完全であるとする。$A^{p, q} = A_{-p}^q$ と置いて二重複体を得る。
このとき、$A_0^\bullet \to M^\bullet$ から誘導される
$\text{Tot}(A^{\bullet, \bullet}) \to M^\bullet$ は擬同型である。
\end{lemma}

\begin{proof}
短完全列
$0 \to \Ker(d^p_{A_n^\bullet}) \to A_n^p \to \Im(d^p_{A_n^\bullet}) \to 0$
と仮定から、すべての $p \in \mathbf{Z}$ に対して
$$
\ldots \to \Im(d_{A_2^\bullet}^p) \to \Im(d_{A_1^\bullet}^p)
\to \Im(d_{A_0^\bullet}^p) \to \Im(d_{M^\bullet}^p) \to 0
$$
が完全であることが分かる。完全列
$0 \to \Im(d^{p - 1}_{A_n^\bullet}) \to \Ker(d^p_{A_n^\bullet})
\to H^p(A_n^\bullet) \to 0$
について同じ議論を繰り返すと、すべての $p \in \mathbf{Z}$ に対して
$$
\ldots \to H^p(A_2^\bullet) \to H^p(A_1^\bullet)
\to H^p(A_0^\bullet) \to H^p(M^\bullet) \to 0
$$
が完全であることが分かる。

\medskip\noindent
$T^\bullet = \text{Tot}(A^{\bullet, \bullet})$ と書く。
$H^0(T^\bullet) \to H^0(M^\bullet)$ が同型であることを示す。
他の次数についても同じ議論が成り立つ。
$x \in \Ker(\text{d}_{T^\bullet}^0)$ を、ある元
$\xi \in H^0(T^\bullet)$ の代表元とする。
$x = \sum_{i = n, \ldots, 0} x_i$ と書く。ただし $x_i \in A_i^i$ である。
$n > 0$ と仮定する。このとき $x_n$ は $d_{A_n^\bullet}^n$ の核に属し、
$H^n(A_{n - 1}^\bullet)$ において零に写る。実際、その像は符号を除いて
$x_{n - 1}$ の境界となる元へ写るからである。
証明の第1段落により、
$x_n \bmod \Im(d^{n - 1}_{A_n^\bullet})$ は
$H^n(A_{n + 1}^\bullet) \to H^n(A_n^\bullet)$ の像に属する。
したがって $x$ を境界だけ変えて、$x_n$ が境界である状況にできる。
$x$ をもう一度変えれば、$x_n = 0$ と仮定してよいことが分かる。
帰納法により、すべてのコホモロジー類 $\xi$ はコサイクル $x = x_0$
によって表される。最後に、核の完全性に関する条件から、二つのコサイクル
$x_0$ と $x_0'$ がコホモロガスであることと、それらの
$H^0(M^\bullet)$ における像が同じであることとは同値である。
\end{proof}

\begin{lemma}
\label{homology-lemma-good-right-resolution-gives-qis}
$M^\bullet$ をアーベル群の複体とする。
$$
0 \to M^\bullet \to A_0^\bullet \to A_1^\bullet \to A_2^\bullet \to \ldots
$$
をアーベル群の複体の完全な複体とし、すべての $p \in \mathbf{Z}$ に対して
複体
$$
0 \to
\Coker(d_{M^\bullet}^p) \to
\Coker(d_{A_0^\bullet}^p) \to
\Coker(d_{A_1^\bullet}^p) \to
\Coker(d_{A_2^\bullet}^p) \to \ldots
$$
も完全であるとする。$A^{p, q} = A_p^q$ と置いて二重複体を得る。
$\text{Tot}_\pi(A^{\bullet, \bullet})$ を、この二重複体に付随する
積全複体とする（証明を参照）。このとき、$M^\bullet \to A_0^\bullet$
から誘導される写像
$M^\bullet \to \text{Tot}_\pi(A^{\bullet, \bullet})$ は擬同型である。
\end{lemma}

\begin{proof}
$T^\bullet = \text{Tot}_\pi(A^{\bullet, \bullet})$ と略記し、
$$
T^n = \prod\nolimits_{p + q = n} A^{p, q} =
\prod\nolimits_{p + q = n} A_p^q
\quad\text{with}\quad
\text{d}_{T^\bullet}^n =
\prod\nolimits_{n = p + q} (f_p^q + (-1)^pd_{A_p^\bullet}^q)
$$
と定める。ここで $f_p^\bullet : A_p^\bullet \to A_{p + 1}^\bullet$ は、
補題における複体の写像である。

\medskip\noindent
$H^0(M^\bullet) \to H^0(T^\bullet)$ が同型であることを示す。
他の次数についても同じ議論が成り立つ。
$x \in \Ker(\text{d}_{T^\bullet}^0)$ を
$\xi \in H^0(T^\bullet)$ の代表元とする。
$x = (x_i)$ と書く。ただし $x_i \in A_i^{-i}$ である。
$x_0$ は $\Coker(A_1^{-1} \to A_1^0)$ において零に写ることに注意する。
したがって、ある $m_0 \in M^0$ および $y \in A_0^{-1}$ に対して
$x_0 = m_0 + d_{A_0^\bullet}^{-1}(y)$ である。
$\text{d}_{T^\bullet}(x) = 0$ から
$\text{d}_{A_0^\bullet}(x_0) = 0$ なので、
$d_{M^\bullet}(m_0) = 0$ である。
ゆえに、$\xi$ を $H^0(M^\bullet) \to H^0(T^\bullet)$ の像に属する元だけ
変えることにより、$x_0$ が $\Im(d^{-1}_{A_0^\bullet})$ に属すると
仮定してよい。

\medskip\noindent
$x_0 \in \Im(d^{-1}_{A_0^\bullet})$ と仮定する。この場合 $\xi = 0$
であることを主張する。これを示すため、$n$ に関する帰納法により、
$y_i \in A_i^{-i - 1}$ を満たす元 $y_0, y_1, \ldots, y_n$ で、
$x_0 = \text{d}_{A_0}^{-1}(y_0)$ および
$j = 1, \ldots, n$ に対して
$x_j = f_{j - 1}^{-j}(y_{j - 1}) + (-1)^j d^{-j - 1}_{A_j^\bullet}(y_j)$
となるものを求める。$n = 0$ の場合は明らかである。
帰納段階を証明する。$y_0, \ldots, y_{n - 1}$ が得られたと仮定する。
このとき $w_n = x_n - f_{n - 1}^{-n}(y_{n - 1})$ は
$d^{-n}_{A_n^\bullet}$ の核に属し、$H^{-n}(A_{n + 1}^\bullet)$ において
零に写る（その像は符号を除いて $x_{n + 1}$ の境界となる元へ写るからである）。
Lemma \ref{homology-lemma-good-resolution-gives-qis} の証明とまったく同様に、
補題の仮定から、すべての $p \in \mathbf{Z}$ に対して
$$
0 \to
H^p(M^\bullet) \to
H^p(A_0^\bullet) \to
H^p(A_1^\bullet) \to
H^p(A_2^\bullet) \to \ldots
$$
が完全であることが従う。したがって $y_{n - 1}$ を
$\Ker(d^{n - 1}_{A_{n - 1}^\bullet})$ の元だけ変えることにより、
$w_n$ が $H^{-n}(A_n^\bullet)$ において零に写ると仮定してよい。
これは所望の $y_n$ が存在することを意味する。この手続きは
$y_0, \ldots, y_{n - 2}$ を変えないことに注意する。
したがってこの手続きを無限に続けると、
$d_{T^\bullet}(y) = \xi$ を満たす $T^{n - 1}$ の元
$y = (y_i)$ が得られる。
% P02-E0016: source line 7003 says y lies in T^{n-1}; the constructed components y_i have total degree -1, so T^{-1} appears required.
これにより $H^0(M^\bullet) \to H^0(T^\bullet)$ が全射であることが分かる。

\medskip\noindent
$m_0 \in \Ker(d^0_{M^\bullet})$ が $H^0(T^\bullet)$ において零に写るとする。
その像は $y = (y_i) \in T^{-1}$ に微分を作用させたものとする。
このとき $y_0 \in A_0^{-1}$ は $\Coker(d^{-2}_{A_1^\bullet})$ において
零に写る。仮定により、これは
$y_0 \bmod \Im(d^{-2}_{A_0^\bullet})$ が、ある $z \in M^{-1}$ の像であることを
意味する。したがって $m_0 = d^{-1}_{M^\bullet}(z)$ である。
これで単射性が証明され、証明が完了する。
\end{proof}

\begin{lemma}
\label{homology-lemma-resolution-gives-qis}
$M^\bullet$ をアーベル群の複体とする。
$$
\ldots \to A_2^\bullet \to A_1^\bullet \to A_0^\bullet \to M^\bullet \to 0
$$
をアーベル群の複体の完全な複体とする。
$A^{p, q} = A_{-p}^q$ と置いて二重複体を得る。
$\text{Tot}_\pi(A^{\bullet, \bullet})$ を、この二重複体に付随する
積全複体とする（証明を参照）。
このとき、$A_0^\bullet \to M^\bullet$ から誘導される写像
$\text{Tot}_\pi(A^{\bullet, \bullet}) \to M^\bullet$ は擬同型である。
\end{lemma}

\begin{proof}
$T^\bullet = \text{Tot}_\pi(A^{\bullet, \bullet})$ と略記し、
$$
T^n = \prod\nolimits_{p + q = n} A^{p, q} =
\prod\nolimits_{p + q = n} A_{-p}^q
\quad\text{with}\quad
\text{d}_{T^\bullet}^n =
\prod\nolimits_{n = p + q} (f_{-p}^q + (-1)^pd_{A_{-p}^\bullet}^q)
$$
と定める。ここで $f_p^\bullet : A_p^\bullet \to A_{p - 1}^\bullet$ は、
補題における複体の写像である。
$M^\bullet$ が零複体であるとき $T^\bullet$ が非輪状であることを示す。
次の技巧により、これで十分である。$B_n^\bullet = A_{n + 1}^\bullet$
および $B_0^\bullet = M^\bullet$ と置く。このとき、補題と同様の完全列
$$
\ldots \to B_2^\bullet \to B_1^\bullet \to B_0^\bullet \to 0 \to 0
$$
がある。$S^\bullet = \text{Tot}_\pi(B^{\bullet, \bullet})$ と置く。
すると、複体の明らかな短完全列
$$
0 \to M^\bullet \to S^\bullet \to T^\bullet[1] \to 0
$$
があり、長完全コホモロジー列から結論を得る。いくつかの細部は省略する。

\medskip\noindent
$M^\bullet = 0$ と仮定する。$H^0(T^\bullet) = 0$ を示す。
他の次数についても同じ議論が成り立つ。
$x =(x_n) \in \Ker(d_{T^\bullet})$ を、ある
$\xi \in H^0(T^\bullet)$ に写る元とし、$x_n \in A^{-n, n} = A_n^n$ とする。
$M^0 = 0$ なので、ある $y_0 \in A_1^0$ に対して
$x_0 = f_1^0(y_0)$ である。
$x$ はコサイクルなので、
$x_1 - d^0_{A_1^\bullet}(y_0) = f_2^1(y_1)$ である。
実際、その左辺は $f_1^1$ により零に写る。ここで $y_1 \in A_2^1$ である。
同様に帰納法を続けると、$d_{T^\bullet}(y) = x$ を満たす
$y = (y_i) \in T^{-1}$ が得られ、所望の結論が従う。
\end{proof}

\section{入射対象}
\label{homology-section-injectives}

\begin{definition}
\label{homology-definition-injective}
$\mathcal{A}$ をアーベル圏とする。対象 $J \in \Ob(\mathcal{A})$ が
{\it 入射対象}であるとは、任意の単射 $A \hookrightarrow B$ と任意の射
$A \to J$ に対し、次の図式を可換にする射 $B \to J$ が存在することをいう。
$$
\xymatrix{
A \ar[r] \ar[d] & B \ar@{-->}[ld] \\
J &
}
$$
\end{definition}

\noindent
入射対象の標準的な特徴づけを述べる。

\begin{lemma}
\label{homology-lemma-characterize-injectives}
$\mathcal{A}$ をアーベル圏とし、$I$ を $\mathcal{A}$ の対象とする。
次は同値である。
\begin{enumerate}
\item 対象 $I$ は入射対象である。
\item 関手 $B \mapsto \Hom_\mathcal{A}(B, I)$ は完全である。
\item $\mathcal{A}$ の任意の短完全列
$$
0 \to I \to A \to B \to 0
$$
は分裂する。
\item すべての $B \in \Ob(\mathcal{A})$ に対して
$\Ext_\mathcal{A}(B, I) = 0$ である。
\end{enumerate}
\end{lemma}

\begin{proof}
省略する。
\end{proof}

\begin{lemma}
\label{homology-lemma-product-injectives}
$\mathcal{A}$ をアーベル圏とする。
$I_\omega$, $\omega \in \Omega$ を $\mathcal{A}$ の入射対象の族とする。
$\prod_{\omega \in \Omega} I_\omega$ が存在するならば、これは入射対象である。
\end{lemma}

\begin{proof}
省略する。
\end{proof}

\begin{definition}
\label{homology-definition-enough-injectives}
$\mathcal{A}$ をアーベル圏とする。すべての対象 $A$ に対し、ある入射対象
$J$ への単射 $A \to J$ が存在するとき、$\mathcal{A}$ は
{\it 十分な入射対象をもつ}という。
\end{definition}

\begin{definition}
\label{homology-definition-functorial-injective-embedding}
$\mathcal{A}$ をアーベル圏とする。次を満たす関手
$$
J : \mathcal{A} \longrightarrow \text{Arrows}(\mathcal{A})
$$
が存在するとき、$\mathcal{A}$ は{\it 関手的入射埋め込みをもつ}という。
\begin{enumerate}
\item $s \circ J = \text{id}_\mathcal{A}$,
\item 任意の対象 $A \in \Ob(\mathcal{A})$ に対して射 $J(A)$ は単射である。
\item 任意の対象 $A \in \Ob(\mathcal{A})$ に対して対象 $t(J(A))$ は
$\mathcal{A}$ の入射対象である。
\end{enumerate}
このような関手を $A \mapsto (A \to J(A))$ と書く。
\end{definition}

\section{射影対象}
\label{homology-section-projectives}

\begin{definition}
\label{homology-definition-projective}
$\mathcal{A}$ をアーベル圏とする。対象 $P \in \Ob(\mathcal{A})$ が
{\it 射影対象}であるとは、任意の全射 $A \rightarrow B$ と任意の射
$P \to B$ に対し、次の図式を可換にする射 $P \to A$ が存在することをいう。
$$
\xymatrix{
A \ar[r] & B \\
P \ar@{-->}[u] \ar[ru] &
}
$$
\end{definition}

\noindent
射影対象の標準的な特徴づけを述べる。

\begin{lemma}
\label{homology-lemma-characterize-projectives}
$\mathcal{A}$ をアーベル圏とし、$P$ を $\mathcal{A}$ の対象とする。
次は同値である。
\begin{enumerate}
\item 対象 $P$ は射影対象である。
\item 関手 $B \mapsto \Hom_\mathcal{A}(P, B)$ は完全である。
\item $\mathcal{A}$ の任意の短完全列
$$
0 \to A \to B \to P \to 0
$$
は分裂する。
\item すべての $A \in \Ob(\mathcal{A})$ に対して
$\Ext_\mathcal{A}(P, A) = 0$ である。
\end{enumerate}
\end{lemma}

\begin{proof}
省略する。
\end{proof}

\begin{lemma}
\label{homology-lemma-coproduct-projectives}
$\mathcal{A}$ をアーベル圏とする。
$P_\omega$, $\omega \in \Omega$ を $\mathcal{A}$ の射影対象の族とする。
$\coprod_{\omega \in \Omega} P_\omega$ が存在するならば、これは射影対象である。
\end{lemma}

\begin{proof}
省略する。
\end{proof}

\begin{definition}
\label{homology-definition-enough-projectives}
$\mathcal{A}$ をアーベル圏とする。すべての対象 $A$ に対し、ある射影対象
$P$ からの全射 $P \to A$ が存在するとき、$\mathcal{A}$ は
{\it 十分な射影対象をもつ}という。
\end{definition}

\begin{definition}
\label{homology-definition-functorial-projective-surjections}
$\mathcal{A}$ をアーベル圏とする。次を満たす関手
$$
P : \mathcal{A} \longrightarrow \text{Arrows}(\mathcal{A})
$$
が存在するとき、$\mathcal{A}$ は{\it 関手的射影全射をもつ}という。
\begin{enumerate}
\item $t \circ P = \text{id}_\mathcal{A}$,
\item 任意の対象 $A \in \Ob(\mathcal{A})$ に対して射 $P(A)$ は全射である。
\item 任意の対象 $A \in \Ob(\mathcal{A})$ に対して対象 $s(P(A))$ は
$\mathcal{A}$ の射影対象である。
\end{enumerate}
このような関手を $A \mapsto (P(A) \to A)$ と書く。
\end{definition}

\section{入射対象と随伴関手}
\label{homology-section-adjoint}

\noindent
随伴関手と入射対象との関係に関するいくつかの補題を述べる。
Lemma \ref{homology-lemma-adjoint-get-abelian} も参照されたい。

\begin{lemma}
\label{homology-lemma-adjoint-preserve-injectives}
\begin{slogan}
完全な左随伴をもつ関手は入射対象を保つ
\end{slogan}
$\mathcal{A}$ と $\mathcal{B}$ をアーベル圏とする。
$u : \mathcal{A} \to \mathcal{B}$ および
$v : \mathcal{B} \to \mathcal{A}$ を加法関手とし、
$u$ は $v$ の右随伴であるとする。次の条件を考える。
\begin{enumerate}
\item[(a)] $v$ は単射を単射に移す。
\item[(b)] $v$ は完全である。
\item[(c)] $u$ は入射対象を入射対象に移す。
\end{enumerate}
このとき (a) $\Leftrightarrow$ (b) $\Rightarrow$ (c) である。
$\mathcal{A}$ が十分な入射対象をもつならば、三条件はすべて同値である。
\end{lemma}

\begin{proof}
$v$ は左随伴であるから右完全であることに注意する
（Categories, Lemma \ref{categories-lemma-exact-adjoint}）。
これと Lemma \ref{homology-lemma-exact-functor} を合わせれば、
(a) $\Leftrightarrow$ (b) である理由が分かる。

\medskip\noindent
(a) を仮定する。$I$ を $\mathcal{A}$ の入射対象とする。
$\varphi : N \to M$ を $\mathcal{B}$ の単射とし、
$\alpha : N \to uI$ を射とする。
随伴性により射 $\alpha : vN \to I$ を得て、仮定により
$v\varphi : vN \to vM$ は単射である。
$I$ は入射対象なので、$\alpha$ を延長する射 $\beta : vM \to I$ を得る。
再び随伴性により、これは $\alpha$ を延長する射
$\beta : M \to uI$ に対応する。したがって (c) が成り立つ。

\medskip\noindent
$\mathcal{A}$ が十分な入射対象をもち、(c) が成り立つと仮定する。
$f : B \to B'$ を $\mathcal{B}$ の単射とし、$A = \Ker(v(f))$ と置く。
$I$ を入射対象として単射 $g : A \to I$ を選ぶ（仮定により存在する）。
すると $g$ は $g' : v(B) \to I$ に延長され、随伴により射
$B \to u(I)$ を得る。$u(I)$ は入射対象なので、この射は
$h : B' \to u(I)$ に延長される。したがって随伴により、
$g'$ を延長する射 $k : v(B') \to I$ を得る。
ところがこのとき $k$ は $g$ を「延長」するので、$g$ が単射であることから
$A = 0$ でなければならない。ゆえに (a) が成り立つ。
\end{proof}

\begin{remark}
\label{homology-remark-need-left-exactness}
$R \to S$ を環準同型とする。$u : \text{Mod}_S \to \text{Mod}_R$ を
$u(N) = N_R$ とし、$v : \text{Mod}_R \to \text{Mod}_S$ を
$v(M) = M \otimes_R S$ とする。このとき $u$ は $v$ の右随伴であり、
$u$ は完全、$v$ は右完全である。しかし一般には、
Lemma \ref{homology-lemma-adjoint-preserve-injectives} の条件 (a), (b), (c) は
成り立たない。たとえば $R = \mathbf{Z}$ および
$S = \mathbf{Z}/p\mathbf{Z}$ とすると、入射的 $S$-加群
$\mathbf{Z}/p\mathbf{Z}$ は入射的 $\mathbf{Z}$-加群ではない。
実際、この補題は、すべての入射的 $S$-加群が $R$-加群として入射的で
あることと、$R \to S$ が平坦な環準同型であることとが同値であることを示す。
\end{remark}

\begin{lemma}
\label{homology-lemma-adjoint-enough-injectives}
$\mathcal{A}$ と $\mathcal{B}$ をアーベル圏とする。
$u : \mathcal{A} \to \mathcal{B}$ および
$v : \mathcal{B} \to \mathcal{A}$ を加法関手とする。次を仮定する。
\begin{enumerate}
\item $u$ は $v$ の右随伴である。
\item $v$ は単射を単射に移す。
\item $\mathcal{A}$ は十分な入射対象をもつ。
\item 任意の $B \in \Ob(\mathcal{B})$ に対し、$vB = 0$ ならば $B = 0$ である。
\end{enumerate}
このとき $\mathcal{B}$ は十分な入射対象をもつ。
\end{lemma}

\begin{proof}
$B \in \Ob(\mathcal{B})$ を選ぶ。
$I$ を $\mathcal{A}$ の入射対象として単射 $vB \to I$ を選ぶ。
Lemma \ref{homology-lemma-adjoint-preserve-injectives} と仮定により、
対応する写像 $B \to uI$ は $B$ から入射対象への単射である。
\end{proof}

\begin{remark}
\label{homology-remark-faithfulness-needed}
$\mathcal{A}$, $\mathcal{B}$, $u : \mathcal{A} \to \mathcal{B}$,
$v : \mathcal{B} \to \mathcal{A}$ を
Lemma \ref{homology-lemma-adjoint-enough-injectives} のとおりとする。
条件 (1) と (2) のもとでは、条件 (4) は $v$ が忠実であることと同値である。
さらに、条件 (4) は必要である。たとえば、関手 $u$ と $v$ がともに
零関手である場合を考えればよい。
\end{remark}

\begin{lemma}
\label{homology-lemma-adjoint-functorial-injectives}
$\mathcal{A}$ と $\mathcal{B}$ をアーベル圏とする。
$u : \mathcal{A} \to \mathcal{B}$ および
$v : \mathcal{B} \to \mathcal{A}$ を加法関手とする。次を仮定する。
\begin{enumerate}
\item $u$ は $v$ の右随伴である。
\item $v$ は単射を単射に移す。
\item $\mathcal{A}$ は十分な入射対象をもつ。
\item 任意の $B \in \Ob(\mathcal{B})$ に対し、$vB = 0$ ならば $B = 0$ である。
\item $\mathcal{A}$ は関手的入射埋め込みをもつ。
\end{enumerate}
このとき $\mathcal{B}$ は関手的入射埋め込みをもつ。
\end{lemma}

\begin{proof}
$A \mapsto (A \to J(A))$ を $\mathcal{A}$ 上の関手的入射埋め込みとする。
このとき $B \mapsto (B \to uJ(vB))$ は $\mathcal{B}$ 上の
関手的入射埋め込みである。
Lemma \ref{homology-lemma-adjoint-enough-injectives} の証明と比較されたい。
\end{proof}

\begin{lemma}
\label{homology-lemma-partially-defined-adjoint}
$\mathcal{A}$ と $\mathcal{B}$ をアーベル圏とし、
$u : \mathcal{A} \to \mathcal{B}$ を関手とする。
次を満たす部分集合 $\mathcal{P} \subset \Ob(\mathcal{B})$ が存在するとする。
\begin{enumerate}
\item $\mathcal{B}$ のすべての対象は $\mathcal{P}$ のある元の商である。
\item すべての $P \in \mathcal{P}$ に対し、$\mathcal{A}$ のある対象 $Q$ が
存在し、$A$ に関して関手的に
$\Hom_\mathcal{A}(Q, A) = \Hom_\mathcal{B}(P, u(A))$ となる。
\end{enumerate}
このとき $u$ の左随伴 $v$ が存在する。
\end{lemma}

\begin{proof}
米田の補題（Categories, Lemma \ref{categories-lemma-yoneda}）により、
$P$ に対応する $\mathcal{A}$ の対象 $Q$ は、式
$\Hom_\mathcal{A}(Q, A) = \Hom_\mathcal{B}(P, u(A))$ によって
一意な同型を除いて定まる。$Q = v(P)$ と書く。
$\Hom_\mathcal{A}(v(P), v(P))$ の $\text{id}_{v(P)}$ に対応する写像を
$i_P : P \to u(v(P))$ と書く。(2) の関手性により、この全単射は
$$
\Hom_\mathcal{A}(v(P), A) \to \Hom_\mathcal{B}(P, u(A)),\quad
\varphi \mapsto u(\varphi) \circ i_P
$$
で与えられる。任意の二元 $P_1, P_2 \in \mathcal{P}$ に対して標準写像
$$
\Hom_\mathcal{B}(P_2, P_1)
\to
\Hom_\mathcal{A}(v(P_2), v(P_1)),\quad
\varphi \mapsto v(\varphi)
$$
があり、これは
$\Hom_\mathcal{B}(P_2, u(v(P_1)))$ における等式
$u(v(\varphi)) \circ i_{P_2} = i_{P_1} \circ \varphi$
によって特徴づけられる。
$\varphi \mapsto v(\varphi)$ は合成と両立することに注意する。
これは特徴づけから直接分かる。したがって $P \mapsto v(P)$ は、
対象が $\mathcal{P}$ の元であるような $\mathcal{B}$ の充満部分圏からの
関手である。

\medskip\noindent
$\mathcal{B}$ の任意の対象 $B$ に対し、完全列
$$
P_2 \to P_1 \to B \to 0
$$
を選ぶ。これは仮定 (1) により可能である。
$v(B)$ を、完全列
$$
v(P_2) \to v(P_1) \to v(B) \to 0
$$
に入る $\mathcal{A}$ の対象として定める。このとき
\begin{align*}
\Hom_\mathcal{A}(v(B), A)
& =
\Ker(\Hom_\mathcal{A}(v(P_1), A) \to \Hom_\mathcal{A}(v(P_2), A)) \\
& =
\Ker(\Hom_\mathcal{B}(P_1, u(A)) \to \Hom_\mathcal{B}(P_2, u(A))) \\
& =
\Hom_\mathcal{B}(B, u(A))
\end{align*}
である。したがって $\mathcal{P} = \Ob(\mathcal{B})$ としてよいこと、
すなわち $v$ がすべての対象上で定義されることが分かる。
\end{proof}

\section{本質的に定値な系}
\label{homology-section-essentially-constant}

\noindent
この節では、加法圏に値をとる本質的に定値な系を論じる。

\begin{lemma}
\label{homology-lemma-essentially-constant-into-karoubian}
$\mathcal{I}$ を圏、$\mathcal{A}$ を前加法的カルービ圏とし、
$M : \mathcal{I} \to \mathcal{A}$ を図式とする。
\begin{enumerate}
\item $\mathcal{I}$ がフィルター圏であると仮定する。次は同値である。
\begin{enumerate}
\item $M$ は本質的に定値である。
\item $X = \colim M$ が存在し、余終フィルター部分圏
$\mathcal{I}' \subset \mathcal{I}$ と、各 $i' \in \Ob(\mathcal{I}')$
に対する直和分解 $M_{i'} = X_{i'} \oplus Z_{i'}$ が存在して、
$X_{i'}$ は $X$ に同型に写り、$\mathcal{I}'$ のある $i' \to i''$ に対して
$Z_{i'}$ は $M_{i''}$ において零に写る。
\end{enumerate}
\item $\mathcal{I}$ が余フィルター圏であると仮定する。次は同値である。
\begin{enumerate}
\item $M$ は本質的に定値である。
\item $X = \lim M$ が存在し、始的な余フィルター部分圏
$\mathcal{I}' \subset \mathcal{I}$ と、各 $i' \in \Ob(\mathcal{I}')$
に対する直和分解 $M_{i'} = X_{i'} \oplus Z_{i'}$ が存在して、
$X$ は $X_{i'}$ に同型に写り、$\mathcal{I}'$ のある $i'' \to i'$ に対して
$M_{i''} \to Z_{i'}$ は零である。
\end{enumerate}
\end{enumerate}
\end{lemma}

\begin{proof}
(1)(a)、すなわち $\mathcal{I}$ がフィルター圏で $M$ が本質的に定値で
あると仮定する。$X = \colim M_i$ と置く。
Categories, Definition
\ref{categories-definition-essentially-constant-diagram} のように
$i$ と $X \to M_i$ を選ぶ。$i$ を始点とする射の終点となる対象からなる
充満部分圏を $\mathcal{I}'$ とする。
$i' \in \Ob(\mathcal{I}')$ とし、射 $i \to i'$ を選ぶ。
このとき $X \to M_i \to M_{i'}$ と $M_{i'} \to X$ の合成は
$X$ 上の恒等写像である。$\mathcal{A}$ はカルービ圏なので、
直和因子分解 $M_{i'} = X_{i'} \oplus Z_{i'}$ を得る。ここで
$Z_{i'} = \Ker(M_{i'} \to X)$ であり、$X_{i'}$ は $X$ に同型に写る。
Categories, Definition
\ref{categories-definition-essentially-constant-diagram} のように、
$M_{i'} \to X \to M_i \to M_k$ が $M_{i'} \to M_k$ に等しくなる
$i \to k$ および $i' \to k$ を選ぶ。
すると $M_{i'} \to M_k$ は $Z_{i'}$ を零に写す。
したがって (1)(b) が成り立つ。

\medskip\noindent
(1)(b)、すなわち $\mathcal{I}$ がフィルター圏で、補題に述べた
$\mathcal{I}' \subset \mathcal{I}$ および各
$i' \in \Ob(\mathcal{I}')$ に対する直和分解
$M_{i'} = X_{i'} \oplus Z_{i'}$ があると仮定する。
$M$ が本質的に定値であることを見るには、
Categories, Lemma \ref{categories-lemma-cofinal-essentially-constant}
により $\mathcal{I}$ を $\mathcal{I}'$ で置き換えてよい。
任意の $i \in \Ob(\mathcal{I})$ を選び、同型 $X_i \to X$ の逆写像と
包含写像 $X_i \to M_i$ の合成を $X \to M_i$ と書く。
$j$ を別の対象とする。$Z_j \to M_k$ が零となる $j \to k$ を選ぶ。
$\mathcal{I}$ はフィルター圏なので、必要なら $k$ を大きくすることにより、
射 $i \to k$ も存在すると仮定してよい。
このとき $M_j \to X \to M_i \to M_k$ と $M_j \to M_k$ はともに
$Z_j$ を零に写す。したがって、直和因子 $Z_k$ を零に写す射
$M_k \to M_l$ を後合成すれば、
$M_j \to X \to M_i \to M_l$ と $M_j \to M_l$ は等しくなる。
すなわち $M$ は本質的に定値である。

\medskip\noindent
(2) の証明は双対的である。
\end{proof}

\begin{lemma}
\label{homology-lemma-direct-sum-from-product-colimit}
$\mathcal{I}$ を圏、$\mathcal{A}$ を加法的カルービ圏とする。
$F : \mathcal{I} \to \mathcal{A}$ および
$G : \mathcal{I} \to \mathcal{A}$ を関手とする。次は同値である。
\begin{enumerate}
\item $\colim_\mathcal{I} F \oplus G$ が存在する。
\item $\colim_\mathcal{I} F$ および $\colim_\mathcal{I} G$ が存在する。
\end{enumerate}
この場合
$\colim_\mathcal{I} F \oplus G =
\colim_\mathcal{I} F \oplus \colim_\mathcal{I} G$ である。
\end{lemma}

\begin{proof}
(1) を仮定し、$W = \colim_\mathcal{I} F \oplus G$ と置く。
$F$ への射影は、冪等な自然変換 $F \oplus G \to F \oplus G$ を定めることに
注意する。したがって Categories, Lemma
\ref{categories-lemma-functorial-colimit} により、冪等自己準同型
$W \to W$ を得る。$\mathcal{A}$ はカルービ圏なので、
Lemma \ref{homology-lemma-karoubian} により対応する直和分解 $W = X \oplus Y$ を得る。
直接的な議論（省略する）により、
$X = \colim_\mathcal{I} F$ および $Y = \colim_\mathcal{I} G$ である。
したがって (2) が成り立つ。(2) から (1) が従う証明は省略する。
\end{proof}

\begin{lemma}
\label{homology-lemma-direct-sum-from-product-essentially-constant}
$\mathcal{I}$ をフィルター圏、$\mathcal{A}$ を加法的カルービ圏とする。
$F : \mathcal{I} \to \mathcal{A}$ および
$G : \mathcal{I} \to \mathcal{A}$ を関手とする。次は同値である。
\begin{enumerate}
\item $F \oplus G : \mathcal{I} \to \mathcal{A}$ は本質的に定値である。
\item $F$ と $G$ は本質的に定値である。
\end{enumerate}
\end{lemma}

\begin{proof}
(1) を仮定する。特に $W = \colim_\mathcal{I} F \oplus G$ が存在するので、
Lemma \ref{homology-lemma-direct-sum-from-product-colimit} により
$W = X \oplus Y$ であり、$X = \colim_\mathcal{I} F$ および
$Y = \colim_\mathcal{I} G$ である。
たとえば Categories, Lemma
\ref{categories-lemma-characterize-essentially-constant-ind}
の特徴づけを用いる直接的な議論（省略する）により、
$F$ は値 $X$ をもつ本質的に定値な関手であり、
$G$ は値 $Y$ をもつ本質的に定値な関手であることが分かる。
したがって (2) が成り立つ。(2) から (1) が従う証明は省略する。
\end{proof}

\section{逆系}
\label{homology-section-inverse-systems}

\noindent
$\mathcal{C}$ を圏とする。Categories, Section
\ref{categories-section-posets-limits} において、前順序集合上の
（圏 $\mathcal{C}$ に値をとる）逆系の概念を定義した。
前順序集合が通常の順序をもつ
$\mathbf{N} = \{1, 2, 3, \ldots\}$ である場合、この
$\mathbf{N}$ 上の逆系はしばしば単に{\it 逆系}と呼ばれる。
これは各 $i \in \mathbf{N}$ に対する $\mathcal{C}$ の対象 $M_i$ と、
各 $i > i'$, $i, i' \in \mathbf{N}$ に対する射
$f_{ii'} : M_i \to M_{i'}$ からなる対 $(M_i, f_{ii'})$ であって、
意味をもつ場合にはさらに
$f_{i'i''} \circ f_{ii'} = f_{ii''}$ を満たすものである。
実際には射 $M_2 \to M_1$, $M_3 \to M_2$ などを与えれば十分であることは
明らかである。したがって逆系はしばしば次のように描かれる。
$$
M_1 \xleftarrow{\varphi_2} M_2 \xleftarrow{\varphi_3} M_3 \leftarrow \ldots
$$
さらに、記法から遷移写像 $\varphi_i$ を省略して、
単に「$(M_i)$ を逆系とする」と述べることが多い。

\medskip\noindent
$\mathcal{C}$ に値をとるすべての逆系は、明らかな射の概念により圏をなす。

\begin{lemma}
\label{homology-lemma-inverse-systems-abelian}
$\mathcal{C}$ を圏とする。
\begin{enumerate}
\item $\mathcal{C}$ が加法圏ならば、$\mathcal{C}$ に値をとる逆系の圏は
加法圏である。
\item $\mathcal{C}$ がアーベル圏ならば、$\mathcal{C}$ に値をとる逆系の圏は
アーベル圏である。逆系の列 $(K_i) \to (L_i) \to (M_i)$ が完全であることと、
各 $K_i \to L_i \to N_i$ が完全であることとは同値である。
% P02-E0017: source line 7624 writes N_i although the sequence introduces M_i; the frozen canonical expression is preserved.
\end{enumerate}
\end{lemma}

\begin{proof}
省略する。
\end{proof}

\noindent
このような逆系の極限（Categories, Section
\ref{categories-section-posets-limits} を参照）は $\lim M_i$ または
$\lim_i M_i$ と書く。$\mathcal{C}$ がアーベル群（または集合）の圏ならば、
極限は常に存在し、実際に次のように記述できる。
$$
\lim_i M_i
=
\{(x_i) \in \prod M_i \mid \varphi_i(x_i) = x_{i - 1}, \ i = 2, 3, \ldots\}
$$
Categories, Section \ref{categories-section-limit-sets} を参照されたい。
しかし、アーベル群の逆系の短完全列
$$
0 \to (A_i) \to (B_i) \to (C_i) \to 0
$$
が与えられても、対応する極限の列が常に完全であるとは限らない。
これをさらに論じるため、次の概念を導入する。

\begin{definition}
\label{homology-definition-Mittag-Leffler}
$\mathcal{C}$ をアーベル圏とする。すべての $i$ に対して
$c = c(i) \geq i$ が存在し、すべての $k \geq c$ について
$$
\Im(A_k \to A_i) = \Im(A_c \to A_i)
$$
となるとき、逆系 $(A_i)$ は{\it Mittag-Leffler 条件}を満たす、
または略して {\it ML} であるという。
\end{definition}

\noindent
アーベル群、環上の加群などのアーベル圏で考える限り、
Mittag-Leffler 条件は $\lim$-関手の完全性を保証するのに十分であることが
分かる。A.\ Neeman の論文（\cite{Neeman-Counterexample} を参照）では、
AB4* をもつ（積が完全である）アーベル圏において、この条件が十分に
強くないことが示されている。

\begin{lemma}
\label{homology-lemma-Mittag-Leffler}
$$
0 \to (A_i) \to (B_i) \to (C_i) \to 0
$$
をアーベル群の逆系の短完全列とする。
\begin{enumerate}
\item 常に列
$$
0 \to \lim_i A_i \to \lim_i B_i \to \lim_i C_i
$$
は完全である。
\item $(B_i)$ が ML ならば、$(C_i)$ も ML である。
\item $(A_i)$ が ML ならば、
$$
0 \to \lim_i A_i \to \lim_i B_i \to \lim_i C_i \to 0
$$
は完全である。
\end{enumerate}
\end{lemma}

\begin{proof}
良い演習である。(3) については
Algebra, Lemma \ref{algebra-lemma-Mittag-Leffler} を参照されたい。
\end{proof}

\begin{lemma}
\label{homology-lemma-apply-Mittag-Leffler}
$$
(A_i) \to (B_i) \to (C_i) \to (D_i)
$$
をアーベル群の逆系の完全列とする。系 $(A_i)$ が ML ならば、列
$$
\lim_i B_i \to \lim_i C_i \to \lim_i D_i
$$
は完全である。
\end{lemma}

\begin{proof}
$Z_i = \Ker(C_i \to D_i)$ および $I_i = \Im(A_i \to B_i)$ と置く。
このとき $\lim Z_i = \Ker(\lim C_i \to \lim D_i)$ であり、
系の短完全列
$$
0 \to (I_i) \to (B_i) \to (Z_i) \to 0
$$
を得る。さらに Lemma \ref{homology-lemma-Mittag-Leffler} により
$(I_i)$ は (ML) をもつ。したがって Lemma
\ref{homology-lemma-Mittag-Leffler} をもう一度適用すると
$\lim B_i \to \lim Z_i$ が全射であることが分かり、補題が証明される。
\end{proof}

\noindent
次の本質的に定値な逆系の特徴づけは、特にそのような逆系が ML をもつことを
示す。

\begin{lemma}
\label{homology-lemma-essentially-constant}
$\mathcal{A}$ をアーベル圏とし、$(A_i)$ を $\mathcal{A}$ の逆系で、
極限 $A = \lim A_i$ をもつものとする。
$(A_i)$ が本質的に定値であること
（Categories, Definition
\ref{categories-definition-essentially-constant-diagram} を参照）と、
ある $i$ が存在し、すべての $j \geq i$ に対して直和分解
$A_j = A \oplus Z_j$ が存在して、(a) 写像 $A_{j'} \to A_j$ が直和分解と
両立し、(b) すべての $j$ に対して $Z_{j'} \to Z_j$ が零となる
ある $j' \geq j$ が存在することとは同値である。
\end{lemma}

\begin{proof}
$(A_i)$ が本質的に定値であると仮定する。このとき、$A \to A_i \to A$ が
恒等写像となるようなある $i$ と射 $A_i \to A$ が存在し、
すべての $j \geq i$ に対して、$A_{j'} \to A_j$ が
$A_{j'} \to A_i \to A \to A_j$ として分解するような
$j' \geq j$ が存在する（最後の写像は $A = \lim A_i$ から来る）。
したがって、すべての $j \geq i$ に対して
$Z_j = \Ker(A_j \to A)$ と置けばよい。逆の証明は省略する。
\end{proof}

\noindent
次の補題は More on Algebra, Lemma
\ref{more-algebra-lemma-Mittag-Leffler} で改良される。

\begin{lemma}
\label{homology-lemma-exact-sequence-ML}
$$
0 \to (A_i) \to (B_i) \to (C_i) \to 0
$$
をアーベル群の逆系の完全列とする。
$(C_i)$ が本質的に定値ならば、$(A_i)$ が ML をもつことと
$(B_i)$ が ML をもつこととは同値である。
\end{lemma}

\begin{proof}
番号を付け直すことにより、遷移写像と両立する
$C_i = C \oplus Z_i$ を仮定してよく、さらにすべての $i$ に対して
$Z_{i'} \to Z_i$ が零となるある $i' \geq i$ が存在すると仮定してよい。
Lemma \ref{homology-lemma-essentially-constant} を参照されたい。

\medskip\noindent
まず $C = 0$、すなわち $C_i = Z_i$ であると仮定する。
$Z_{n_{i + 1}} \to Z_{n_i}$ が零となるように
$1 = n_1 < n_2 < n_3 < \ldots$ を選ぶ。
このとき $B_{n_{i + 1}} \to B_{n_i}$ は
$A_{n_i} \subset B_{n_i}$ を経由して分解する。したがって
$j \geq i + 1$ に対して、$A_{n_i}$ の部分集合として
$$
\Im(A_{n_j} \to A_{n_i}) \subset \Im(B_{n_j} \to B_{n_i}) \subset
\Im(A_{n_{j - 1}} \to A_{n_i})
$$
を得る。ゆえに、$j \geq i + 1$ に対して像
$\Im(A_{n_j} \to A_{n_i})$ が安定することと、
像 $\Im(B_{n_j} \to B_{n_i})$ が安定することとは同値である。
これから同値性が従う（細部は省略する）。

\medskip\noindent
$C \not = 0$ ならば、$B_i \to C \oplus Z_i$ による $C$ の逆像を
$B'_i \subset B_i$ と書く。前段落により、$(B'_i)$ が ML をもつことと
$(B_i)$ が ML をもつこととは同値である。したがって $(B_i)$ を
$(B'_i)$ で置き換えてよい。この場合、すべての $i$ に対して完全列
$0 \to A_i \to B_i \to C \to 0$ がある。
したがって、すべての $j \geq i$ に対して
$0 \to \Im(A_j \to A_i) \to \Im(B_j \to B_i) \to C \to 0$
は短完全である。ゆえに、$j \geq i$ に対して像
$\Im(A_j \to A_i)$ が安定することと、
$\Im(B_j \to B_i)$ が安定することとは同値であり、所望の結論を得る。
\end{proof}

\noindent
次の補題の「正しい」版は More on Algebra, Lemma
\ref{more-algebra-lemma-apply-Mittag-Leffler-again} である。

\begin{lemma}
\label{homology-lemma-apply-Mittag-Leffler-again}
$$
(A^{-2}_i \to A^{-1}_i \to A^0_i \to A^1_i)
$$
をアーベル群の複体の逆系とし、その極限を
$A^{-2} \to A^{-1} \to A^0 \to A^1$ と書く。
コホモロジーの逆系を $(H_i^{-1})$, $(H_i^0)$ と書き、
$A^{-2} \to A^{-1} \to A^0 \to A^1$ のコホモロジーを
$H^{-1}$, $H^0$ と書く。$(A^{-2}_i)$ と $(A^{-1}_i)$ が ML であり、
$(H^{-1}_i)$ が本質的に定値ならば、$H^0 = \lim H_i^0$ である。
\end{lemma}

\begin{proof}
$Z^j_i = \Ker(A^j_i \to A^{j + 1}_i)$ および
$I^j_i = \Im(A^{j - 1}_i \to A^j_i)$ と置く。
核をとることは極限と可換するので、
$\lim Z^0_i = \Ker(\lim A^0_i \to \lim A^1_i)$ であることに注意する。
系 $(I^{-1}_i)$ と $(I^0_i)$ は、系 $(A^{-2}_i)$ と
$(A^{-1}_i)$ の商として ML をもつ。
Lemma \ref{homology-lemma-Mittag-Leffler} を参照されたい。
したがって逆系の完全列
$$
0 \to (I^{-1}_i) \to (Z^{-1}_i) \to (H^{-1}_i) \to 0
$$
があり、$(I^{-1}_i)$ は ML をもち、$(H^{-1}_i)$ は仮定により
本質的に定値である。ゆえに Lemma \ref{homology-lemma-exact-sequence-ML} により、
$(Z^{-1}_i)$ は ML をもつ。完全列
$$
0 \to (Z^{-1}_i) \to (A^{-1}_i) \to (I^0_i) \to 0
$$
と Lemma \ref{homology-lemma-Mittag-Leffler} の適用により、
$\lim A^{-1}_i \to \lim I^0_i$ は全射である。
最後に、完全列
$$
0 \to (I^0_i) \to (Z^0_i) \to (H^0_i) \to 0
$$
と Lemma \ref{homology-lemma-Mittag-Leffler} により、
$\lim I^0_i \to \lim Z^0_i \to \lim H^0_i \to 0$ は完全である。
以上を合わせれば結論を得る。
\end{proof}

\noindent
上の補題で、大きな順序数上の極限をとる版が必要になることがある。

\begin{lemma}
\label{homology-lemma-ML-over-ordinals}
$\alpha$ を順序数とする。$K_\beta^\bullet$, $\beta < \alpha$ を、
$\alpha$ 上のアーベル群の複体の逆系とする。
すべての $\beta < \alpha$ に対して複体 $K_\beta^\bullet$ が非輪状であり、
写像
$$
K^n_\beta \longrightarrow \lim_{\gamma < \beta} K^n_\gamma
$$
が全射であるならば、複体
$\lim_{\beta < \alpha} K_\beta^\bullet$ は非輪状である。
\end{lemma}

\begin{proof}
超限帰納法により、すべての順序数 $\alpha$ と補題のようなすべての系に
対して主張が成り立つことを証明する。特に、$\alpha$ に対する結果を
証明する際には、複体
$\lim_{\gamma < \beta} K^n_\gamma$ が非輪状であると仮定してよい。
% P02-E0018: source lines 7877-7878 call the degree-n object limit a complex; the frozen K^n_gamma expression is preserved.

\medskip\noindent
$x \in \lim_{\beta < \alpha} K^0_\beta$ が $\text{d}(x) = 0$ を満たすとする。
$\text{d}(y) = x$ を満たす
$y \in \lim_{\beta < \alpha} K^{-1}_\beta$ を求める。
$\beta < \alpha$ に対し、$x$ の像を $x_\beta \in K_\beta^0$ として
$x = (x_\beta)$ と書く。$y = (y_\beta)$ を超限再帰により構成する。

\medskip\noindent
$\beta = 0$ に対しては、$\text{d}(y_0) = x_0$ を満たす任意の元
$y_0 \in K_0^{-1}$ をとる。

\medskip\noindent
$\beta = \gamma + 1$ が後続順序数であるとする。
遷移写像 $f : K^\bullet_\beta \to K^\bullet_\gamma$ によって
$y_\gamma$ に写り、かつ微分によって $x_\beta$ に写る元 $y_\beta$ を
求めなければならない。最初の近似として
$\text{d}(y'_\beta) = x_\beta$ を満たす $y'_\beta$ を選ぶ。
このとき差 $y_\gamma - f(y'_\beta)$ は微分の核に属するので、
ある $z_\gamma \in K^{-2}_\gamma$ に対して
$\text{d}(z_\gamma)$ に等しい。仮定により写像
$f^{-2} : K^{-2}_\beta \to K^{-2}_\gamma$ は全射である。
したがって $z_\gamma = f(z_\beta)$ と書き、
$y'_\beta$ を $y_\beta = y'_\beta + \text{d}(z_\beta)$ に変えればよい。

\medskip\noindent
$\beta$ が極限順序数ならば、
$\lim_{\gamma < \beta} K^{-1}_\gamma$ の元
$(y_\gamma)_{\gamma < \beta}$ があり、その微分は $x_\beta$ の像である。
したがって、各項で全射な複体の写像
$f : K_\beta^\bullet \to \lim_{\gamma < \beta} K_\gamma^\bullet$
と、帰納法により
$\lim_{\gamma < \beta} K_\gamma^\bullet$ が非輪状であると仮定してよいこと
（証明の第1段落を参照）を用いて、上とまったく同じ仕方で論じられる。
\end{proof}

\section{積の完全性}
\label{homology-section-product-exact}


\begin{lemma}
\label{homology-lemma-product-abelian-groups-exact}
$I$ を集合とする。各 $i \in I$ に対し、
$L_i \to M_i \to N_i$ をアーベル群の複体とする。
$H_i = \Ker(M_i \to N_i)/\Im(L_i \to M_i)$ をそのコホモロジーとする。
このとき
$$
\prod L_i \to \prod M_i \to \prod N_i
$$
は、コホモロジー $\prod H_i$ をもつアーベル群の複体である。
\end{lemma}

\begin{proof}
省略する。
\end{proof}
